\documentclass[11pt]{article}
\usepackage[margin=1in]{geometry}
\usepackage{amsmath,amssymb,mathtools}
\usepackage{booktabs,longtable,array}
\usepackage{enumitem}
\usepackage{microtype}
\usepackage{xurl}
\usepackage[hidelinks]{hyperref}
\setlist{nosep,leftmargin=*}
\setlength{\parindent}{0pt}
\setlength{\parskip}{0.55em}

\title{Reproducing Consumption-Growth Identification from the Stochastic Discount Factor}
\author{}
\date{}

\begin{document}
\maketitle

\begin{abstract}
This manual gives the equations and numerical rules for U.S. consumption-growth analysis. A
stochastic discount factor (SDF) weights future payoffs across states. The analysis builds
expected-SDF and SDF-news series from Treasury yields and term premia; consumption and supplied
return scores enter separately. Term premia compensate for
interest-rate risk; SDF news contains unexpected SDF movements. Principal components summarize
series that move together. An instrument based on realized
yield volatility restricts the consumption-growth mean equation to a mean identified set, the
coefficient values that satisfy the stated restrictions. Four conditional-scale estimator mappings carry that set into
conditional-scale estimates, which describe the size of consumption growth left unexplained by the
mean equation. A 10,000-draw circular moving-block bootstrap resamples consecutive quarter blocks
for inference, and two calculations bound omitted approximation terms. Part I follows only the
current configuration and its comparison and sensitivity checks. Part II catalogues
permitted modifications and unavailable extensions.
\end{abstract}

\tableofcontents
\newpage

\part{Current configured reproduction}

\section{Reproduction target and fixed current choices}
\label{sec:target}

The \emph{current configuration} is the fixed collection of input identities, parameter values,
numerical controls, and recorded decisions used in the present reproduction. Part I follows this
configuration alone.

The analysis studies quarterly consumption growth using expected and unexpected movements in a
stochastic discount factor (SDF). An \emph{SDF} is a random variable that prices a
future payoff by weighting that payoff across states of the world. The analysis constructs its
expected-SDF series and its SDF-news series from Treasury yields and term premia. A \emph{term
premium} is the part of a long yield that compensates an investor for bearing interest-rate risk
rather than for expected future short rates. A \emph{short rate} is an interest rate for a short
maturity; the relevant short rate in this manual has a three-month maturity.

A \emph{principal component}, or \emph{PC}, is a linear combination of the columns of a
rectangular numerical matrix. The chosen combination maximizes sample variance subject to zero
sample covariance with earlier PCs. A \emph{PC score} is its value for one observation.

The analysis summarizes each of the three SDF panel matrices of Section~\ref{sec:pca} with three PC scores. It then studies the conditional
mean of consumption growth and the conditional scale of its structural mean residual. The
\emph{structural mean residual} is observed
consumption growth minus the conditional mean implied by the mean equation. A
\emph{predictor} is an observed variable used to describe an outcome's conditional distribution.
The \emph{conditional mean} is the mean of consumption growth given the observed predictors. A
\emph{conditional scale} describes the size of the structural mean residual given those predictors.
Two conditional-scale estimators describe the \emph{conditional variance}, which is the mean
squared structural mean residual given the predictors. One of the remaining two describes the
conditional mean of the log squared structural mean residual, which differs from the log conditional
variance by a fixed offset. The other describes the \emph{conditional
median absolute-residual scale}, which is the median absolute structural mean residual given the
predictors.

\emph{Ordinary least squares}, abbreviated \emph{OLS}, minimizes the sum of squared errors. It is the
benchmark mean estimator and the building block of the first conditional-scale estimator below, and
is not itself one of the four.

The four conditional-scale estimators are defined here before their
abbreviations first appear.
\emph{Logarithmic least squares}, abbreviated \emph{log-OLS}, applies OLS to the log squared
structural mean residual. \emph{Poisson pseudo-maximum likelihood}, abbreviated
\emph{PPML}, fits an exponential conditional variance through its mean equation without treating
the squared residual as count data. \emph{Harvey normal quasi-maximum likelihood}, abbreviated
\emph{Harvey QML}, uses a normal-likelihood-shaped objective without requiring normal residuals.
\emph{Least-absolute-deviation median regression}, abbreviated \emph{LAD}, minimizes absolute
log-residual errors and estimates a conditional median.

The calculation follows one dependency chain. First, confirm that the preserved inputs are present
and unchanged, before any output location is created or cleared, so that a drifted input stops the
reproduction with the existing outputs intact. Then create the output locations and clear the
conditional-output locations. Load
the quarterly Treasury input and construct the three SDF panel matrices. Then construct consumption
growth, realized yield volatility, and the lagged supplied return scores. Perform
the three PC constructions and fit the consumption-growth mean equation. Calculate the initial mean
identified set and variance shares, followed by log-OLS, the slack profiles of the mean identified set,
PPML, and Harvey QML. Next run the conditional-scale checks, close the residual-dynamics
check, calculate LAD, and conduct bootstrap inference. Render the conditional-scale and
mean-identified-set outputs before calculating approximation-error bounds, the descriptive report,
and the final record of requested and completed conditional calculations.

A \emph{reported output} is a required table, figure,
rendered document, machine-readable table, or reusable result. A \emph{typesetting source} is a text
document that specifies content and layout. A \emph{rendered document} is the final typeset document
produced from a typesetting source. A \emph{machine-readable table} is arranged for direct numerical
reading rather than typeset display. A \emph{reusable result} is a file
that preserves an ordered collection of named values with their types and dimensions for later reuse.

An exact reproduction specifies the observation window, input vintages, numerical conventions, and
random-number stream before any calculation. An \emph{input vintage} is the exact version of an
input available on a stated date. A \emph{random-number stream} is the deterministic sequence
of pseudo-random values produced by a stated pseudo-random-number generator and seed. The
\emph{seed} is the initial numerical value that determines the resulting random-number stream.

Set the observation windows as follows. The PC window runs from 1962Q1 through 2025Q4. The
preserved consumption series runs from 1947Q1 through 2026Q2, which retains the early observations
needed for lags. Date matching and row deletion determine each final sample: remove any row
missing a quantity required for that calculation and compute the observation count from the retained rows.

Preserve the following inputs.

\begin{enumerate}
\item The monthly Treasury input contains U.S. zero-coupon yields and term premia by maturity; the
daily Treasury input contains yields only. Obtain them from the Adrian--Crump--Moench replication archive at
\url{https://github.com/fernando-duarte/ACM_term_premium}. Use the bundled June 2026 monthly release,
whose monthly observations run from June 1961 through May 2026, one per month with no gaps. Use the
preserved July 2026 daily release.
\item The consumption input is the quarterly U.S. real personal consumption expenditures series
identified as PCECC96 and published at \url{https://fred.stlouisfed.org/series/PCECC96}. Read it
from the preserved copy distributed with the reproduction materials. That copy carries one
observation per quarter from 1947Q1 through 2026Q2, measured in billions of chained dollars,
seasonally adjusted and expressed at an annual rate. The reproduction reads this copy from storage
and contacts no network service for it, so its vintage is fixed. The publisher revises this series
and offers no historical vintage at that address, which is why the preserved copy, rather than the
published address, fixes the numbers.
\item The four supplied return scores are columns \(\mathrm{pc1},\ldots,\mathrm{pc4}\) of the
supplied return-score archive distributed with the replication input materials.
Preserve the archive, its date column, and its four score columns. Convert each date to its calendar-quarter end
and label the score row one quarter later. A row dated \(t\) therefore becomes the lagged supplied return
score for quarter \(t+1\). The
financial assets and original PC construction used to create these scores lie outside this analysis.
Treat the four supplied return scores as fixed inputs rather than reconstructed quantities.
\end{enumerate}

The analysis uses three current expected-SDF PC scores, three lagged expected-SDF
PC scores, three SDF-news PC scores, and four supplied return scores.
Use binary64 arithmetic throughout, meaning the 64-bit floating-point representation of real
numbers with 52 explicitly stored fraction bits.

The main maturity range uses every integer maturity from 3 through 117 months. The
approximation-error calculations instead use every third month from 3 through 117 months. Read the monthly Treasury
input at every integer maturity from 1 through 120 months. The upper end is needed because the
expected log bond price at maturity \(i\) requires the yield and the term premium three months beyond
\(i\); the lower end is needed because the next quarter's leg of the news difference evaluates that
same price three months below \(i\), reaching maturities 1 and 2. Read the daily Treasury input at the
same maturities. Only its five-year column enters identification; the descriptive report additionally
displays realized yield volatility at 3, 12, and 117 months.
The news horizon is
one quarter, or three months. Display slack values are 0.05, 0.10, and 0.20; the baseline is 0.05.

Here \emph{slack} scales the permitted departure from a residual-product orthogonality condition
relative to the identifying variation in the corresponding SDF-news reduced-form residual. An
\emph{SDF-news reduced-form residual} is the part of an SDF-news PC score left after removing its
linear relation with an intercept and the three lagged expected-SDF PC scores. An \emph{instrument}
is an observed variable whose covariance with a residual product supplies identifying information.
Residual-product orthogonality requires the instrument to be uncorrelated with the
product of the structural mean residual and an SDF-news reduced-form residual.

\section{Calendar, sequences, and summary operators}
\label{sec:calendar}

Let \(\mathbb R\) denote the set of real numbers. A prime on a vector or matrix denotes its
transpose. Let \(t=1,\ldots,T\) index retained quarters in increasing order. Here \(T\) is the
observation count in the particular sample under discussion.

Assign each observation to the calendar quarter that contains it, and label that quarter by its last
calendar day: 31 March, 30 June, 30 September, or 31 December. Every series enters the analysis under
this single label, so series of different frequencies and different native labelling conventions
merge on identical dates. Within a quarter, sort the monthly observations by date and retain the last
one, relabelled to the quarter-end day. A quarter whose last available monthly observation does not
fall in the quarter's terminal month is \emph{incomplete}; retain it, relabelled the same way, and
flag it. Assign each daily observation to its containing quarter and leave its own date unchanged.

For a scalar sequence \(a_1,\ldots,a_T\), define its sample mean and centered version by
\begin{equation*}
 \bar a=\frac{1}{T}\sum_{t=1}^{T}a_t,
 \qquad \widetilde a_t=a_t-\bar a.
\end{equation*}
Here \(a_t\) is observation \(t\) of the scalar sequence, and \(T\) is its observation count.
The scalar \(\bar a\) is the sample mean, and \(\widetilde a_t\) is the centered value.

For scalar sequences \(a_t\) and \(b_t\), define the divisor-\(T\) covariance and variance by
\begin{equation}
 \widehat{\operatorname{Cov}}_T(a,b)
 =\frac{1}{T}\sum_{t=1}^{T}\widetilde a_t\widetilde b_t,
 \qquad
 \widehat{\operatorname{Var}}_T(a)
 =\widehat{\operatorname{Cov}}_T(a,a).
 \label{eq:covT}
\end{equation}
Here \(b_t\) is observation \(t\) of the second scalar sequence,
\(\widehat{\operatorname{Cov}}_T(a,b)\) is their scalar divisor-\(T\) covariance, and
\(\widehat{\operatorname{Var}}_T(a)\) is the nonnegative scalar divisor-\(T\) variance. When the two
arguments are vector sequences \(a_t\in\mathbb R^{p}\) and \(b_t\in\mathbb R^{q}\) with \(p\) and
\(q\) positive integers, the same divisor-\(T\) centering gives the \(p\)-by-\(q\) matrix
\(\widehat{\operatorname{Cov}}_T(a,b)\) whose entry \((j,\ell)\) is the scalar divisor-\(T\)
covariance of component \(j\) of the first sequence with component \(\ell\) of the second; the
divisor-\(T\) variance of a vector sequence is the case \(b=a\).

A \emph{missing value} is an absent observation. A \emph{nonfinite value} is a missing value, positive
infinity, or negative infinity. A \emph{complete row} has no missing required value; a
\emph{finite row} has only finite required values. A \emph{retained row} is a row that remains
after the rule stated for a calculation. Unless a calculation states otherwise, a mean,
covariance, standard deviation, or regression uses only finite rows.

A \emph{numerical tolerance} is the largest numerical discrepancy treated as zero or as an
acceptable match. An absolute tolerance is measured in the units of the discrepancy; a relative
tolerance divides the discrepancy by a stated numerical scale.

Write
\[
 \operatorname{sd}_T(a)=\{\widehat{\operatorname{Var}}_T(a)\}^{1/2},
 \qquad
 \operatorname{corr}_T(a,b)=
 \frac{\widehat{\operatorname{Cov}}_T(a,b)}
 {\operatorname{sd}_T(a)\operatorname{sd}_T(b)}.
\]
Here \(\operatorname{sd}_T(a)\) is the nonnegative divisor-\(T\) standard deviation, and
\(\operatorname{corr}_T(a,b)\) is the scalar divisor-\(T\) correlation. When the text says ordinary sample
standard deviation, it uses
\[
 \operatorname{sd}_{T-1}(a)=
 \left\{\frac{1}{T-1}\sum_{t=1}^{T}(a_t-\bar a)^2\right\}^{1/2}.
\]
Here \(\operatorname{sd}_{T-1}(a)\) is the nonnegative ordinary sample standard deviation of
\(a_1,\ldots,a_T\).

Let \(\operatorname{med}(a)\) denote the sample median of the finite values of \(a_t\). Define the
median absolute deviation by
\begin{equation}
 \operatorname{MAD}(a)=1.4826\,
 \operatorname{med}\bigl(|a_t-\operatorname{med}(a)|\bigr).
 \label{eq:mad}
\end{equation}
Here \(\operatorname{MAD}(a)\) is the nonnegative median absolute deviation of the finite values of
\(a_t\), and the letters MAD stand for median absolute deviation. The scalar
\(\operatorname{med}(a)\) is their sample median. The constant 1.4826 makes MAD
comparable to the standard deviation under a normal distribution.

For a stated finite index set \(\mathcal A\), \(\operatorname{mean}_{t\in\mathcal A}\) denotes the
arithmetic mean over its indices. The floor \(\lfloor x\rfloor\) is the greatest integer no larger
than a real number \(x\), and the ceiling \(\lceil x\rceil\) is the least integer no smaller than
\(x\).

For a vector \(x=(x_1,\ldots,x_q)'\in\mathbb R^q\), let \(q\) be a positive integer and
define
\[
 \|x\|_2=\left(\sum_{j=1}^q x_j^2\right)^{1/2},
 \qquad
 \|x\|_\infty=\max_{1\le j\le q}|x_j|.
\]
Here \(x\in\mathbb R^q\) is a \(q\)-dimensional vector, \(\|x\|_2\) is its Euclidean norm, and
\(\|x\|_\infty\) is its maximum norm. For scalars \(a_1,\ldots,a_q\),
\(\operatorname{diag}(a_1,\ldots,a_q)\) is the \(q\)-by-\(q\) matrix with those scalars on its
main diagonal and zeros elsewhere; applying \(\operatorname{diag}\) to a vector has the same meaning.

Let \(d_s\) be a positive integer. For a vector sequence
\(s_t\in\mathbb R^{d_s}\), define the four-lag Bartlett long-run
covariance sum
\begin{equation*}
 \widehat M_4(s)=\sum_{t=1}^{T}s_ts_t'
 +\sum_{k=1}^{4}\left(1-\frac{k}{5}\right)
 \sum_{t=k+1}^{T}(s_ts_{t-k}'+s_{t-k}s_t').
\end{equation*}
Here \(s_t\in\mathbb R^{d_s}\) is contribution vector \(t\), and
\(\widehat M_4(s)\in\mathbb R^{d_s\times d_s}\) is its symmetric four-lag Bartlett long-run
covariance sum.

\section{Input observations and quarterly transformations}

\subsection{Treasury yields and term premia}

Use the Adrian--Crump--Moench estimates of the Treasury term structure, meaning Treasury yields and
term premia across maturities. For maturity
\(i\in\{1,\ldots,120\}\) months and calendar date \(\mathsf d\), let \(y_{i,\mathsf d}\) be the
annualized model-implied zero-coupon Treasury yield in percentage points and let
\(p_{i,\mathsf d}\) be the annualized term premium
in percentage points. A zero-coupon yield is the discount rate implied by a bond that makes one
payment at maturity and no interim coupon payments. An annualized rate is expressed at the rate that
would apply over one year.

The June 2026 monthly release supplies 780 observations. The preserved July 2026 daily panel begins on 14 June 1961 and
follows business days. Both panels span maturities from one month through
ten years. Neither panel is trimmed by date when it is read; the observation windows of
Section~\ref{sec:target} are imposed later, on the quarterly date. The daily panel supplies yields
only; term premia are read from the monthly panel alone.

Quarterly yield and term-premium observations are the last monthly values in each quarter. Denote
them by \(y_{i,t}\) and \(p_{i,t}\). Use daily yields separately to construct realized yield volatility.
The term-structure formulas divide percentage points by 100 whenever they require decimal rates.

\subsection{Consumption growth}

Let \(C_t>0\) be quarterly real personal consumption expenditures in quarter \(t\), using the
series identified as PCECC96 by the Federal Reserve Economic Data service. Real means that the
series has been adjusted for changes in its price index. The level is an aggregate, not a per-capita,
series, measured in billions of chained dollars, seasonally adjusted and expressed at an annual rate.
The growth rate below is a ratio of two consecutive levels, so the annual-rate scaling and the
chained-dollar base cancel; the seasonal adjustment does not cancel and is part of the input
identity. The preserved copy carries one observation in every quarter of its span, so \(C_{t-1}\) is
the level in the immediately preceding calendar quarter. Consumption growth is the simple
quarterly percentage change
\begin{equation*}
 g_t=100\left(\frac{C_t}{C_{t-1}}-1\right).
\end{equation*}
Here \(g_t\) is scalar consumption growth in percentage points. The positive scalars \(C_t\) and
\(C_{t-1}\) are the current and preceding quarterly U.S. real personal consumption expenditures
levels. The formula uses levels rather than logarithms and retains the quarterly rate. Mark growth
as missing in the first available quarter.

\subsection{Realized yield volatility and the instrument}

Let the \(D\) daily dates in ascending calendar order be
\[
 \mathsf d_1<\cdots<\mathsf d_D.
\]
Here \(D\) is the positive integer number of daily dates, and \(\mathsf d_d\) is daily date \(d\).
For quarter \(t\), define the corresponding integer-index set
\[
 \mathcal D_t=\{d:\mathsf d_d\text{ falls in quarter }t\}.
\]
Here \(\mathcal D_t\) is the finite set of integer daily-date indices assigned to quarter \(t\).

For maturity \(i\), form daily changes before dividing the observations into quarters. Quarterly
realized yield volatility is
\begin{equation*}
 v_{i,t}=\left(\sum_{d\in\mathcal D_t}
       (y_{i,\mathsf d_d}-y_{i,\mathsf d_{d-1}})^2\right)^{1/2}.
\end{equation*}
Here \(v_{i,t}\) is the nonnegative scalar realized yield volatility for maturity \(i\) in quarter
\(t\), and \(y_{i,\mathsf d_d}\) is the scalar yield observed on daily date \(\mathsf d_d\). The
first difference within a quarter therefore uses the last daily observation of the preceding
quarter. Mark the change for the first global daily observation as missing and omit missing daily
changes from the sum. If a quarter contains no nonmissing daily change at maturity \(i\), mark
\(v_{i,t}\) missing rather than zero, so that the row-deletion rule removes the quarter instead of
admitting it as an extreme-low volatility. The units are percentage points of the annualized yield.
The volatility itself is a per-quarter realized quantity and carries no annualization factor.

The scalar instrument is the centered five-year realized yield volatility. Specifically,
let
\begin{equation*}
 z_t=v_{60,t}-\bar v_{60}.
\end{equation*}
Here \(z_t\) is the scalar instrument in quarter \(t\), \(v_{60,t}\) is five-year realized yield
volatility, and \(\bar v_{60}\) is its sample mean. The subscript 60 is the maturity in months, and
the mean is taken on the \emph{identification sample}, meaning
the quarters with complete consumption growth, the three lagged expected-SDF PC scores, the three
SDF-news PC scores, and the instrument. Here the instrument is one-dimensional. It is measured in
percentage points, the same units as the realized yield volatility from which it is formed. Realized
yield volatility is defined in every quarter of the daily panel; only the quarters of the
identification sample enter the mean \(\bar v_{60}\) and the centered values \(z_t\) used below.

\subsection{Supplied return scores}

Let \(r_{a,t}\) be supplied return score \(a\in\{1,2,3,4\}\) for quarter \(t\). Each supplied
return score is a PC score of nominal financial-asset returns.
Convert each supplied date to its quarter end and then relabel it one quarter forward. The value
joined to row \(t\) is the supplied return score from quarter \(t-1\). The scores are dimensionless
and carry the supplied archive's own scale; apply no transformation when reading them.

\section{Term-structure construction of the SDF}
\label{sec:sdf}

Set the news horizon to \(h=3\) months. For maturity \(i\in\{1,\ldots,120\}\), define
\(\mathcal M_i=i/12\), the scalar maturity in years. For maturity \(i\in\{1,\ldots,117\}\), the
expected log bond price for the next quarter is
\begin{equation}
 n_{i,t}=\frac{\mathcal M_i y_{i,t}-\mathcal M_{i+h}y_{i+h,t}
       +\mathcal M_{i+h}p_{i+h,t}-\mathcal M_i p_{i,t}}{100}.
 \label{eq:nhat}
\end{equation}
Here \(n_{i,t}\) is the scalar expected log bond price for maturity \(i\) in quarter \(t\).
The quantities \(\mathcal M_i\) and \(\mathcal M_{i+h}\) are maturities in years, \(y\) denotes
yields, and \(p\) denotes term premia. The horizon is \(h=3\) months.

At the news-horizon maturity \(i=h\), impose \(p_{h,t}=0\). This fixed reference treats the
one-quarter term premium as zero.
Use the supplied term premia directly as inputs to the expected log bond price for this
three-month news horizon. Their \emph{rollover convention} is the rule that connects a bond before
maturity to a shorter-maturity bond one quarter later. The supplied term premia are not reconciled to
the one-quarter horizon; whatever convention wedge they carry is inherited by \(n_{i,t}\) and is left
there.

Define log bond-price news for quarter \(t+1\) by
\begin{equation*}
 \Delta q_{i,t+1}=n_{i-h,t+1}-n_{i,t}.
\end{equation*}
Here \(\Delta q_{i,t+1}\) is scalar log bond-price news between quarters \(t\) and \(t+1\), and
\(\Delta\) denotes the first-difference operator.

For the boundary term \(i-h=0\), set
\begin{equation}
 n_{0,t+1}=-\frac{h}{12}\frac{y_{h,t+1}}{100}.
 \label{eq:price-news-boundary}
\end{equation}
Here \(n_{0,t+1}\) is the scalar boundary expected log bond price, and \(y_{h,t+1}\) is the
three-month yield in quarter \(t+1\). Equation~\eqref{eq:nhat} is evaluated at every integer maturity
from 1 through 117, so the difference above is defined at maturities 4 and 5 through
\(n_{1,t+1}\) and \(n_{2,t+1}\); the main maturity range is the narrower 3 through 117. Impose the zero
term premium only at the news-horizon maturity \(i=h\)---maturities 1 and 2 are shorter and carry
their supplied term premia---and use the boundary value above only for \(i-h=0\). Every news sequence in this section, the log bond-price news
and the SDF-news series built from it alike, has a missing first observation, because each requires a
previous-quarter quantity. Each therefore has one fewer usable row than the panel has quarters.

The SDF-news series is the centered second-order approximation, meaning the approximation that
retains the linear and squared log bond-price news terms,
\begin{equation*}
 N_{i,t+1}=e^{n_{i,t}}\left(\Delta q_{i,t+1}
                   +\frac{1}{2}\Delta q_{i,t+1}^2\right)-m_i^{N}.
\end{equation*}
Here \(N_{i,t+1}\) is the scalar SDF-news value, and \(m_i^N\) is its scalar centering correction for
maturity \(i\). The scalar centering correction is
\begin{equation*}
 m_i^{N}=\frac{1}{2}\operatorname{mean}_{t\in\mathcal Q_i}
       \left[e^{n_{i,t}}\Delta q_{i,t+1}^2\right].
\end{equation*}
Here \(m_i^N\) is one half of the sample mean of the squared news term over the finite index set
\(\mathcal Q_i\). The set \(\mathcal Q_i\) contains quarters in which both factors are present. Retain
infinite values; any such value makes the expression infinite. The two SDF constructions use
deliberately different admission rules: \(\mathcal Q_i\) admits every non-missing value, including an
infinite one, whereas \(\mathcal G_i\) below admits only finite values. The correction \(m_i^N\)
removes only the sample mean of the squared term.

The expected-SDF series uses an unpaired mean correction that equates two sample means without
shifting either sequence to match the maturity horizon:
\begin{equation}
 E_{i,t}=e^{n_{i,t}}+
 \operatorname{mean}_{t\in\mathcal G_i}
       e^{-\frac{h}{12}y_{h,t}/100}
       -\operatorname{mean}_{t\in\mathcal G_i}e^{n_{i,t}}.
 \label{eq:expected-sdf}
\end{equation}
Here \(E_{i,t}\) is the scalar expected-SDF value for maturity \(i\) in quarter \(t\), and
\(\mathcal G_i\) is the common finite index set for the two sample means.

The unpaired mean correction has four implications. It applies to every integer maturity in the main
maturity range, including integer maturities between consecutive multiples of three. It leaves a subsequently
centered PC unchanged, but remains part of the intermediate series required for exact reproduction. Its
additive form can make \(E_{i,t}\) nonpositive. Because the correction uses every quarter in
\(\mathcal G_i\), its full-sample construction prevents real-time calculation from information
available in each quarter.

\section{Three separate PC constructions}
\label{sec:pca}

Construct three \emph{SDF panel matrices} on quarters from 1962Q1 through 2025Q4. The current
expected-SDF panel matrix contains \(E_{i,t}\) and has dimension \(T_E\times115\). Here \(T_E\)
is its row count, and every row is complete. Form the lagged expected-SDF panel matrix by relabeling each row
of the expected-SDF series one quarter forward and rebuilding a complete matrix. It retains
only complete rows, has row count \(T_{EL}\), and has dimension \(T_{EL}\times115\). The SDF-news panel matrix contains
\(N_{i,t}\) and has dimension \(T_N\times115\). Here \(T_N\) is its row count, and every row is complete. In
each matrix, 115 is the number of integer maturities from 3 through 117. The three SDF panel matrices
have different complete-row sets, because the one-quarter-forward relabelling precedes the fixed
observation window, and they receive three separate PC constructions. An infinite
retained value invalidates the corresponding PC construction.

For any one SDF panel matrix \(\mathsf X^{\mathrm{PC}}\in\mathbb R^{T_s\times115}\), let \(T_s\)
denote that matrix's row count and let
\(\bar X_i^{\mathrm{PC}}\) be maturity column \(i\)'s sample mean.
Let \(s_i^{\mathrm{PC}}>0\) be its ordinary sample standard deviation,
\[
 s_i^{\mathrm{PC}}=\left\{\frac{1}{T_s-1}\sum_{t=1}^{T_s}
 (X_{ti}^{\mathrm{PC}}-\bar X_i^{\mathrm{PC}})^2\right\}^{1/2}.
\]
Here \(s_i^{\mathrm{PC}}>0\) is the ordinary sample standard deviation of maturity column \(i\),
and \(T_s\) is the selected SDF panel matrix's row count. The scalar
\(X_{ti}^{\mathrm{PC}}\) is entry \((t,i)\), and \(\bar X_i^{\mathrm{PC}}\) is column \(i\)'s
sample mean.

Define the standardized matrix
\begin{equation*}
 \widetilde X_{ti}^{\mathrm{PC}}=
 \frac{X_{ti}^{\mathrm{PC}}-\bar X_i^{\mathrm{PC}}}{s_i^{\mathrm{PC}}}.
\end{equation*}
Here \(\widetilde X_{ti}^{\mathrm{PC}}\) is standardized entry \((t,i)\) of the selected SDF panel
matrix, \(\bar X_i^{\mathrm{PC}}\) is column \(i\)'s sample mean, and
\(s_i^{\mathrm{PC}}>0\) is its ordinary sample standard deviation.

Compute the singular-value decomposition with the binary64 divide-and-conquer procedure and
retain the signs returned by that procedure:
\begin{equation*}
 \widetilde{\mathsf X}^{\mathrm{PC}}=\mathsf U\mathsf D\mathsf V'.
\end{equation*}
Here \(r_{\mathrm{PC}}=\min(T_s,115)\) is the positive integer number of returned singular-vector
pairs. The matrices \(\mathsf U\in\mathbb R^{T_s\times r_{\mathrm{PC}}}\) and
\(\mathsf V\in\mathbb R^{115\times r_{\mathrm{PC}}}\) have
orthonormal columns, meaning that every column has length one, and any two distinct columns
have dot product zero. The matrix
\(\mathsf D\in\mathbb R^{r_{\mathrm{PC}}\times r_{\mathrm{PC}}}\) is diagonal with nonincreasing
nonnegative entries called \emph{singular values}. A singular value measures how strongly the
matrix stretches its corresponding pair of singular-vector directions. The calculation applies no
rank-based truncation. The PC score matrix is
\(\widetilde{\mathsf X}^{\mathrm{PC}}\mathsf V=\mathsf U\mathsf D\). Retain its first
three columns.

Define the vectors of three PC scores
\begin{align*}
 F^E_t&=(F^E_{1t},F^E_{2t},F^E_{3t})'\in\mathbb R^3,\\
 L_t&=(F^{EL}_{1t},F^{EL}_{2t},F^{EL}_{3t})'\in\mathbb R^3,\\
 Y_{2t}&=(F^N_{1t},F^N_{2t},F^N_{3t})'\in\mathbb R^3.
\end{align*}
Here \(F_t^E\) contains the three current expected-SDF PC scores, \(L_t\) contains the three lagged
expected-SDF PC scores from their separate PC construction, and \(Y_{2t}\) contains the three
SDF-news PC scores. Each vector belongs to \(\mathbb R^3\).

Apply the \emph{lagged expected-SDF PC score sign alignment}
only to the lagged expected-SDF PC scores: for each \(k\in\{1,2,3\}\), multiply \(F^{EL}_{kt}\) by the sign of its sample
correlation with the one-quarter-lagged \(F^E_{kt}\). If that correlation is zero, retain a zero
multiplier. A \emph{sign convention} chooses consistently between a PC and its negative. A
\emph{PC loading} is a coefficient in the linear combination that defines a PC\@. Apply no further
sign change to the current expected-SDF or SDF-news PC scores. Use the same binary64
divide-and-conquer procedure for all three PC constructions and retain its raw signs. The current
configuration supplies no further mathematical sign normalization. Exact signed-coefficient
reproduction therefore requires the binary64 procedure to return these same signs. The current
procedure does not correct opposite raw signs.

The mean-equation design uses the lagged expected-SDF PC scores:
\begin{equation*}
 X_t=(1,L_t')'\in\mathbb R^4.
\end{equation*}
Here \(X_t\in\mathbb R^4\) is the four-dimensional mean-equation design vector containing an
intercept followed by the three lagged expected-SDF PC scores.

\section{Mean equation and its two samples}
\label{sec:meaneq}

Let \(Y_{1t}=g_t\) denote the scalar consumption growth. The \emph{preliminary mean sample}
contains the complete rows formed from consumption growth, lagged expected-SDF PC scores, and
SDF-news PC scores. Estimate the preliminary mean regression on this sample:
\begin{equation}
 Y_{1t}=X_t'\gamma_X+Y_{2t}'\gamma_N+u_t^{O}.
 \label{eq:prelim-ols}
\end{equation}
Here \(\gamma_X\in\mathbb R^4\) contains the intercept and coefficients on the lagged expected-SDF
PC scores. The vector \(\gamma_N\in\mathbb R^3\) contains the coefficients on the SDF-news PC
scores. The scalar \(u_t^{O}\) is the \emph{preliminary mean residual}: observed consumption growth
minus its fitted value in the preliminary mean regression. The superscript \(O\) denotes ordinary
least squares (OLS). OLS
chooses the seven coefficients to minimize the sum of squared preliminary mean residuals.

The descriptive report uses an OLS covariance that assumes constant preliminary mean residual
variance. Let \(\mathcal I_{\mathrm{pm}}\) denote the preliminary mean sample, with observation count
\(N_{\mathrm{pm}}=|\mathcal I_{\mathrm{pm}}|\). Let \(X_{\mathrm{pm}}\in\mathbb R^{N_{\mathrm{pm}}\times7}\) stack the seven-predictor row vectors of the preliminary mean
sample, one per retained quarter; its rows are not the four-dimensional design vectors \(X_t\) but
the seven-predictor vectors written \(x_t^{O\prime}\) below.
Let \(\widehat u_{\mathrm{pm},t}\) be the fitted preliminary mean residual. Require
\(N_{\mathrm{pm}}>7\) and an invertible matrix \(X_{\mathrm{pm}}'X_{\mathrm{pm}}\). Then
\begin{equation*}
 \widehat s_{\mathrm{pm}}^2=\frac{\sum_{t\in\mathcal I_{\mathrm{pm}}}\widehat u_{\mathrm{pm},t}^2}{N_{\mathrm{pm}}-7},
 \qquad
 \widehat V_{\mathrm{pm}}=\widehat s_{\mathrm{pm}}^2(X_{\mathrm{pm}}'X_{\mathrm{pm}})^{-1}.
\end{equation*}
Here \(\mathcal I_{\mathrm{pm}}\) is the preliminary mean sample, and \(N_{\mathrm{pm}}\) is its observation count. The
scalar \(\widehat u_{\mathrm{pm},t}\) is fitted preliminary mean residual \(t\), and \(\widehat s_{\mathrm{pm}}^2\) is
the nonnegative preliminary mean residual variance estimate. The matrix
\(\widehat V_{\mathrm{pm}}\in\mathbb R^{7\times7}\) is the coefficient covariance matrix formed from the
\(N_{\mathrm{pm}}\)-by-7 design matrix \(X_{\mathrm{pm}}\).

For inference on each coefficient, use the square root of its diagonal covariance entry as its
standard error. Divide the coefficient by that standard error and calculate the two-sided p-value
from a Student \(t\) distribution with \(N_{\mathrm{pm}}-7\) degrees of freedom. The degrees of freedom control
the shape of this reference distribution. The null hypothesis states that the coefficient equals
zero. A \emph{p-value} is the probability, under that hypothesis, that the absolute test statistic
is at least as large as its observed absolute value.

The identification sample additionally requires the scalar instrument \(z_t\). This requirement
removes every quarter in which the instrument is missing and the other three quantities are present. Sort this sample by quarter and refit
equation~\eqref{eq:prelim-ols} as the \emph{mean least-squares benchmark}. All identification
calculations use this sample. If
\(x_t^{O}=(X_t',Y_{2t}')'\in\mathbb R^7\) is its complete vector of seven predictors and
\(\widehat u_t^{O}\) is its fitted mean least-squares benchmark residual, report covariance
\[
\widehat V_O=
 \left(\sum_tx_t^Ox_t^{O\prime}\right)^{-1}
 \widehat M_4(x^O\widehat u^O)
 \left(\sum_tx_t^Ox_t^{O\prime}\right)^{-1}.
\]
Here \(x_t^O\in\mathbb R^7\) is the benchmark predictor vector, and \(\widehat u_t^O\) is its
scalar fitted mean least-squares benchmark residual. The matrix \(\widehat M_4(x^O\widehat u^O)\) is the four-lag Bartlett
long-run covariance sum. The matrix \(\widehat V_O\in\mathbb R^{7\times7}\) is the mean
least-squares benchmark coefficient covariance matrix.
Use the four-lag Bartlett long-run covariance sum directly for the benchmark covariance.

If \(N_O\) is the identification-sample observation count, divide each coefficient estimate by its
reported standard error and calculate its two-sided p-value from a Student \(t\) distribution with
\(N_O-7\) degrees of freedom. The two samples differ whenever the instrument is missing in a quarter
the preliminary mean sample retains, and the preliminary mean regression and the mean least-squares
benchmark then have different coefficient vectors.

Compute two reduced-form specifications. A \emph{specification} is a stated collection of
equations and restrictions. A \emph{reduced form} expresses an observed outcome
as a linear function of the predictors plus a residual. In both specifications, estimate
\begin{equation*}
 Y_{1t}=X_t'\beta_{1R}+W_{1t},
 \qquad \beta_{1R}\in\mathbb R^4.
\end{equation*}
Here \(\beta_{1R}\) is the four-dimensional outcome reduced-form coefficient vector, and \(W_{1t}\)
is the scalar \emph{outcome reduced-form residual}.

Specification B is the reported specification because it removes the linear relation between the
SDF-news PC scores and \(X_t\). Estimate its three equations:
\begin{equation*}
 Y_{2t}=\beta_{2R}X_t+W_{2t},
 \qquad
 \beta_{2R}\in\mathbb R^{3\times4},\quad W_{2t}\in\mathbb R^3.
\end{equation*}
Here \(\beta_{2R}\) is the three-by-four SDF-news reduced-form coefficient matrix, and \(W_{2t}\) is
the three-dimensional SDF-news reduced-form residual vector.

Specification A is the \emph{comparison specification}. It serves only as a \emph{comparison
check}, meaning a calculation that leaves the reported specification unchanged. It leaves the
SDF-news PC scores unresidualized and shows how
that choice changes the mean identified set. It uses the same \(W_{1t}\), sets \(W_{2t}=Y_{2t}\), and imposes
\(\beta_{2R}=0_{3\times4}\), the three-by-four zero matrix. Under either specification, for a candidate structural SDF-news
coefficient vector \(\theta\in\mathbb R^3\), recover the coefficients on the lagged expected-SDF PC scores as
\begin{equation}
 \beta_1(\theta)=\beta_{1R}-\beta_{2R}'\theta.
 \label{eq:structural-recovery}
\end{equation}
Here \(\beta_1(\theta)\in\mathbb R^4\). The first element is the intercept and the remaining three
are the coefficients on the lagged expected-SDF PC scores.

\section{Set identification of the mean equation}
\label{sec:idset}

The identifying restriction uses changes in the variance of the SDF-news reduced-form residuals.
When the instrument changes the
squared SDF-news reduced-form residuals, near orthogonality between the instrument and each product
of the structural mean residual and an SDF-news reduced-form residual restricts the structural
SDF-news coefficient vector. The common slack \(\tau\) permits a proportional departure from exact
orthogonality. The analysis maintains the residual-product orthogonality condition. The identification
tests and checks describe only whether the instrument moves residual variance. The coefficients are
structural because they define the structural mean residual; the restriction by itself leaves
causality unidentified.

\subsection{Sample moments}

A \emph{moment} is a mean of a product or power of quantities. A \emph{sample moment} calculates
that mean from observed quantities. The identification calculation uses the following seven sample
variances and covariances, all computed on the identification sample.
The structural mean residual at a candidate \(\theta\) is
\begin{equation}
 e_t(\theta)=W_{1t}-W_{2t}'\theta.
 \label{eq:mean-residual}
\end{equation}
Here \(e_t(\theta)\) is the scalar structural mean residual at structural SDF-news coefficient
vector \(\theta\in\mathbb R^3\). The scalar \(W_{1t}\) is the outcome reduced-form residual, and
\(W_{2t}\in\mathbb R^3\) is the SDF-news reduced-form residual vector.

For coordinate \(k\in\{1,2,3\}\) of the SDF-news reduced-form residual, define the scalar
\(h_{kt}=W_{1t}W_{2k,t}\). Also define the vector
\(G_{kt}=W_{2t}W_{2k,t}\in\mathbb R^3\) and the scalar \(q_{kt}=W_{2k,t}^2\). Using the centered
divisor-\(T\) operator in equation~\eqref{eq:covT}, calculate seven sample moments:
\begin{align*}
 S_k^{(0)}&=\widehat{\operatorname{Var}}_T(h_k)\in\mathbb R,
 &s_k^2&=\widehat{\operatorname{Var}}_T(q_k)\in\mathbb R,\\
 R_k^{(0)}&=\widehat{\operatorname{Cov}}_T(z,h_k)\in\mathbb R,
 &R_k^{(1)}&=\widehat{\operatorname{Cov}}_T(z,G_k)\in\mathbb R^{1\times3},\\
 P_k^{(0)}&=\widehat{\operatorname{Cov}}_T(z,q_k)\in\mathbb R,
 &S_k^{(1)}&=\widehat{\operatorname{Cov}}_T(G_k,h_k)\in\mathbb R^3,\\
 &&S_k^{(2)}&=\widehat{\operatorname{Var}}_T(G_k)\in\mathbb R^{3\times3}.
\end{align*}
Here \(h_{kt}\) is the scalar residual product, \(G_{kt}\in\mathbb R^3\) is its vector counterpart,
and \(q_{kt}\) is the squared coordinate-\(k\) SDF-news reduced-form residual. This \(q\) carries a
coordinate subscript and a quarter subscript and no argument. It is a different quantity from the
squared structural mean residual \(q_t(\theta)\) of Section~\ref{sec:logvar-scale}, which carries one
subscript and an argument, and from the log bond-price news \(\Delta q_{i,t+1}\) of
Section~\ref{sec:sdf}, which carries the first-difference operator and a maturity
subscript. The scalar sample
moments are \(S_k^{(0)},s_k^2,R_k^{(0)}\), and \(P_k^{(0)}\). The remaining sample moments are the
row vector \(R_k^{(1)}\in\mathbb R^{1\times3}\), the vector
\(S_k^{(1)}\in\mathbb R^3\), and the symmetric matrix
\(S_k^{(2)}\in\mathbb R^{3\times3}\). An \emph{instrument weight} is the coefficient with which the
instrument enters a moment condition. Because the instrument is one-dimensional its weight is one,
which is nonzero and so leaves all three inequalities in force, and the seven expressions require no
further weighting vector. A zero weight would make \(A_k\), \(b_k\) and \(c_k\) below all vanish, so
inequality \(k\) would hold at every candidate and drop out of the mean identified set.

\subsection{Quadratic inequalities}

The correlation form of the restriction requires \(\widehat{\operatorname{Var}}_T(z)>0\) and
\(\widehat{\operatorname{Var}}_T\{e(\theta)W_{2k}\}>0\) to be equivalent to the quadratic
form below. Require \(s_k^2=\widehat{\operatorname{Var}}_T(W_{2k}^2)\) to be finite and strictly
positive for every coordinate \(k\) before assembling the inequalities, and stop the reproduction
otherwise; thereafter evaluate the quadratic form at every candidate without further tests. The slack
\(\tau_k\) scales the permitted residual-product correlation. The relaxed restriction is
\begin{equation}
 \left|\operatorname{corr}_T\{z,e(\theta)W_{2k}\}\right|
 \le \tau_k\left|\operatorname{corr}_T(z,W_{2k}^2)\right|.
 \label{eq:relaxed-correlation}
\end{equation}
Here \(z\) is the scalar instrument sequence, \(e(\theta)\) is the structural mean residual sequence,
\(W_{2k}\) is coordinate \(k\) of the SDF-news reduced-form residual sequence, and
\(\tau_k\in[0,1)\) is its slack parameter. Slack values are admissible only on that half-open unit
interval. The right-hand correlation scales the allowed departure by
the instrument's relation with the squared SDF-news reduced-form residual.

For the slack \(\tau_k\in[0,1)\) of coordinate \(k\) of the SDF-news reduced-form residual, let
\begin{align*}
 Q_k&=R_k^{(1)}\in\mathbb R^{1\times3},
 &V_k&=(P_k^{(0)})^2,\\
 w_k&=\frac{\tau_k^2V_k}{s_k^2}\in\mathbb R,\quad w_k\ge0,
 &A_k&=Q_k'Q_k-w_kS_k^{(2)}\in\mathbb R^{3\times3},\\
 b_k&=-2R_k^{(0)}Q_k'+2w_kS_k^{(1)}\in\mathbb R^3,
 &c_k&=(R_k^{(0)})^2-w_kS_k^{(0)}\in\mathbb R.
\end{align*}
Here \(Q_k\in\mathbb R^{1\times3}\) is the inequality row vector, and \(V_k\ge0\) is the squared
instrument covariance. The scalar \(w_k\ge0\) is the inequality weight, and
\(A_k\in\mathbb R^{3\times3}\) is the symmetric quadratic matrix. The vector
\(b_k\in\mathbb R^3\) is the linear term, and \(c_k\in\mathbb R\) is the scalar constant. Both terms
of \(A_k\) are symmetric in exact arithmetic; replace the computed \(A_k\) by
\((A_k+A_k')/2\) before any later use, because the searches below use an analytic constraint
derivative that assumes exact symmetry. As a reproduction check on the seven sample moments,
coordinate \(k\) of \(Q_k\) equals \(P_k^{(0)}\), since both are the divisor-\(T\) covariance of the
instrument with \(W_{2k}^2\).

The analysis uses a common slack, \(\tau_k=\tau\), for all three coordinates of the SDF-news
reduced-form residual. The mean identified set is
\begin{equation}
 \Theta(\tau)=\left\{\theta\in\mathbb R^3:
 \theta'A_k\theta+b_k'\theta+c_k\le0
 \text{ for }k=1,2,3\right\}.
 \label{eq:identified-set}
\end{equation}
Here \(\Theta(\tau)\subseteq\mathbb R^3\) is the mean identified set at common slack \(\tau\);
\(\theta\) is a structural SDF-news coefficient vector; and \(A_k,b_k,c_k\) are the quadratic,
linear, and constant terms for inequality \(k\). These inequalities are the quadratic form of equation~\eqref{eq:relaxed-correlation}. At
slack \(\tau\), a structural SDF-news coefficient vector is \emph{feasible} when it belongs to
\(\Theta(\tau)\), equivalently when it satisfies all three inequalities. At
\(\tau=0\), they reduce to the linear system
\begin{equation}
 Q\theta=\mathsf m.
 \label{eq:tau-zero-system}
\end{equation}
Here \(Q\in\mathbb R^{3\times3}\) stacks the rows \(Q_k\), and the sample-moment vector is
\(\mathsf m=(R_1^{(0)},R_2^{(0)},R_3^{(0)})'\in\mathbb R^3\).

Use a column-pivoted QR decomposition of \(Q\), meaning a factorization \(Q\Pi=\mathsf{QR}\) in which
\(\Pi\) is a permutation matrix, \(\mathsf Q\) has orthonormal columns, and
\(\mathsf R\) is upper triangular. The permutation leaves the column order unchanged except that a
column judged numerically dependent moves behind the columns that remain independent; it does not
reorder the columns by descending norm. At each stage, call a remaining column numerically
dependent when its remaining Euclidean norm is at most \(10^{-8}\) times its original Euclidean norm. The zero-slack point is available only if all
three columns remain independent and the computed QR solution is finite. It must also satisfy
\begin{equation*}
 \|Q\theta-\mathsf m\|_\infty\le10^{-8}\max(1,\|\mathsf m\|_\infty).
\end{equation*}
Here \(Q\theta-\mathsf m\in\mathbb R^3\) is the zero-slack equation residual, and the right-hand side is
its absolute acceptance tolerance. Report the fast triangular one-norm condition estimate for the
upper-triangular factor of a column-pivoted QR decomposition of \(Q\) taken at relative tolerance
\(10^{-7}\); this is a different tolerance from the \(10^{-8}\) used for the rank test and the
solution above, so the pivoting, and hence the factor, can differ. This estimate approximates
\(\|R\|_1\|R^{-1}\|_1\), where \(R\) is that factor and \(\|\cdot\|_1\) is the largest absolute
column sum. Exclude the estimate from the \emph{acceptance rules}. An acceptance rule is a condition
that must hold before retaining a numerical result. A condition number measures how strongly small
numerical perturbations can change the solution.

Every endpoint has one of four \emph{endpoint states}. \emph{Bounded} means a finite accepted
endpoint. \emph{Unbounded} means the expanding-box rule supports an infinite endpoint.
\emph{Unreliable} means the search evidence is inconclusive. \emph{Failed} means that the enclosing
calculation ended without running the endpoint search to completion, in which case every endpoint and
every point that calculation would have produced carries this state. No full-sample endpoint search
of this section produces it. Endpoint-state severity follows this order: bounded,
unbounded, unreliable, failed.

A \emph{slack profile} records lower and upper values and their endpoint states across stated slack
values. Record zero slack with equal lower and upper entries taken from the zero-slack point, in the
bounded endpoint state when that point passes the stated rank, finiteness, and equation-residual
acceptance rules, and in the unreliable endpoint state otherwise. In the separate coarse slack sequence that
locates the critical slack, record zero slack as bounded with zero total width defined below and run no search.

\subsection{Endpoint searches and endpoint states}

For a \emph{scalar functional} \(\phi(\theta)\), meaning a rule that maps \(\theta\) to one real
number, the lower identified endpoint is the infimum, or greatest lower bound, of \(\phi\) over the
mean identified set \(\Theta(\tau)\) of equation~\eqref{eq:identified-set}. The upper identified
endpoint is the supremum, or least upper bound.
The pair consisting of the lower and upper identified endpoints is the \emph{identified range}.

Two families of functionals are reported. The three \emph{coordinate functionals} are
\(\theta\mapsto\theta_\ell\) for \(\ell\in\{1,2,3\}\); their identified range is taken directly over
\(\Theta(\tau)\). The four \emph{design functionals} come from
equation~\eqref{eq:structural-recovery}: writing \(e_p\in\mathbb R^4\) for the \(p\)-th standard
basis vector and \(\lambda_p=\beta_{2R}e_p\in\mathbb R^3\) for the loading vector of design coefficient
\(p\in\{1,2,3,4\}\), the functional is \(\phi_p(\theta)=\beta_{1R,p}-\lambda_p'\theta\). Because the
loading enters with a minus sign, the two sides reverse: the lower identified endpoint of design
coefficient \(p\) is \(\beta_{1R,p}-\sup_{\theta\in\Theta(\tau)}\lambda_p'\theta\) and its upper identified
endpoint is \(\beta_{1R,p}-\inf_{\theta\in\Theta(\tau)}\lambda_p'\theta\). Each side inherits the endpoint
state of the search that produced it, so the two endpoint states reverse with the two endpoints.

Nonconvex quadratic inequalities define the mean identified set. Search for initial coordinate and
design endpoints over expanding boxes centered at the origin. A \emph{box} is a
rectangular region formed by lower and upper limits on every coordinate. Nonconvex means that a
line segment between two points in the set can leave the set. A separate \emph{widening search}
later widens the displayed identified ranges of the structural SDF-news coefficients. Map accepted
structural SDF-news coefficient vectors through equation~\eqref{eq:structural-recovery} to widen the
corresponding identified ranges of the coefficients on the lagged expected-SDF PC scores. The
widening search repeats the same minimization from several
starting values. A \emph{warm starting value} is a previously attained solution reused as a starting value,
and \emph{attained} means found by the stated numerical search.

Let \(a_\star=\operatorname{med}_k\varrho(A_k)\) and
\(c_\star=\operatorname{med}_k|c_k|\). Here \(\varrho(A_k)\) is the largest absolute eigenvalue
of \(A_k\). An eigenvalue is a scalar by which a square matrix stretches some nonzero vector
without changing that vector's direction; such a vector is an eigenvector. Define the positive search-coordinate scale
\begin{equation*}
 d_\theta=
 \begin{cases}
 (c_\star/a_\star)^{1/2},&a_\star>0, c_\star>0,
 \text{and both are finite},\\
 1,&\text{otherwise}.
 \end{cases}
\end{equation*}
Here \(d_\theta>0\) is the search-coordinate scale, \(a_\star\) is the median largest absolute
eigenvalue, and \(c_\star\) is the median absolute inequality constant. For inequality \(k\), let its raw inequality scale be the largest of
\(\varrho(A_k)d_\theta^2\), \(\|b_k\|_2d_\theta\), and \(|c_k|\). Set any nonfinite raw scale to zero.
Among the strictly positive raw scales, calculate the median, and let the threshold be \(10^{-12}\)
times that median. Replace every raw scale below the threshold, zero included, by the threshold. If no
raw scale is strictly positive, set every scale to one. Denote
the resulting positive scale for inequality \(k\) by \(\omega_k\). For a candidate \(\theta\), define
\begin{equation*}
 g_k(\theta)=\theta'A_k\theta+b_k'\theta+c_k,
 \qquad
 \rho_Q(\theta)=\max_{k\in\{1,2,3\}}\frac{g_k(\theta)}{\omega_k}.
\end{equation*}
Here \(g_k(\theta)\) is the scalar left-hand side of inequality \(k\), and
\(\rho_Q(\theta)\) is the largest normalized inequality value.

The initial identified-endpoint searches for coordinate and design functionals use search boxes with multipliers
\(M_{\mathrm{box}}\in\{10^6,10^9,10^{10}\}\), corresponding to
\(|\theta_\ell|\le d_\theta M_{\mathrm{box}}\) for structural SDF-news coefficient coordinate
\(\ell\in\{1,2,3\}\). Use sequential least-squares quadratic programming, which replaces the
nonlinear problem locally with a quadratic objective and linearized constraints. Each search begins
at the origin. The \emph{objective} is the scalar quantity being minimized or maximized. Each search
uses relative iterate tolerance \(10^{-8}\), permits at most 1,000 objective evaluations, and accepts
normalized feasibility at \(10^{-4}\). The relative iterate tolerance measures the change between successive candidate
coefficient vectors relative to their numerical scale.

Two distinct tests apply. The \emph{feasibility screen} is one-sided: a search result passes it when
its coordinates are finite and \(\rho_Q(\theta)\le10^{-4}\); at the second and third boxes, a result
that fails the screen is discarded and the endpoint takes the unreliable state. The \emph{boundary check} is two-sided:
\(|\rho_Q(\theta)|\le10^{-4}\), which additionally verifies that the search reached the boundary of
the mean identified set rather than stopping strictly inside it, and it is what certifies a finite
endpoint. Call an objective value \emph{accepted} when it passes the feasibility screen. Do not
additionally require the search to report success: on legitimately
ill-conditioned bounds the search reports a roundoff-limited outcome at an otherwise valid solution,
and requiring success would discard it.

Write \(v_{\mathrm{old}}\) and
\(v_{\mathrm{new}}\) for accepted objective values from consecutive
boxes. A candidate lies away from the current box edge when
\begin{equation*}
 |\theta_\ell|<0.99d_\theta M_{\mathrm{box}}
 \quad\text{for every }\ell\in\{1,2,3\}.
\end{equation*}
This away-from-edge test governs the rule that finalizes an endpoint outright and the bounded
endpoint state, and in both it runs on all three coordinates, for coordinate and design functionals
alike. Two later rules instead use a form relaxed to the target coordinate alone, and only for a
coordinate endpoint: such an endpoint counts as away from the edge when its own target coordinate
satisfies the inequality, even if another coordinate lies on an edge. The first is the rule that
carries an accepted value forward across boxes: a box at which the target-coordinate test fails
clears the stored previous value, so the next box that passes the test begins a new pair rather than
being compared with the last box that passed. The second is the final-box rule stated below. A design
functional's endpoint is gated by neither relaxed rule: every accepted value pairs with the
immediately preceding accepted value, and the final box leaves an unresolved design endpoint
unreliable rather than unbounded.

Apply the two tests below in the stated
order: test stability first and stop there when it holds, and test growth only when stability fails.
Call two accepted values stable when
\begin{equation*}
 |v_{\mathrm{new}}-v_{\mathrm{old}}|
 \le10^{-3}\max\{1,|v_{\mathrm{new}}|\},
\end{equation*}
and call the objective magnitude growing when
\begin{equation*}
 |v_{\mathrm{new}}|\ge5\max\{|v_{\mathrm{old}}|,d_\theta\}.
\end{equation*}
The smallest box does not enter the cross-box comparison, and the feasibility screen is not applied
there. Its result finalizes the endpoint outright when its coordinates are finite and it lies away
from every edge, in which case the boundary check alone separates the bounded from the unreliable
state; otherwise the search moves to the second box. The first stability and growth
comparison is therefore between the second and third boxes. Assign the bounded endpoint state to a
candidate that passes both the feasibility screen and the boundary check and that lies away from
every edge or is stable across boxes. Assign the unbounded endpoint state to a growing endpoint. Also assign it to a
coordinate endpoint whose target coordinate remains on the final box edge. Otherwise, assign the
unreliable endpoint state.

Walk the display slacks in increasing order and, at each one, run a separate widening search in the
smallest box. At the smallest display slack, start from the origin, the six signed axis points at
distance \(d_\theta\), and the zero-slack point when it is available. At each later display slack,
start from the origin, the same six signed axis points, and the structural SDF-news coefficient
vectors the immediately preceding display slack accepted; the zero-slack point is not re-added. In
later rounds, start from the vectors accepted in the preceding round, for at most four rounds in
total, and stop earlier as soon as a round produces no starting value that has not already been
solved from. Remove duplicate starting values after rounding each coordinate to six significant
digits.

The widening search uses the same relative iterate tolerance \(10^{-8}\), the same limit of
1,000 objective evaluations, and the same normalized tolerance \(10^{-4}\) as the initial
searches, applied here as the two-sided boundary check \(|\rho_Q(\theta)|\le10^{-4}\): a solve landing
strictly inside the mean identified set is not accepted, neither to move an endpoint nor as a
starting value for the next round.

Move an endpoint only when both endpoints of its coefficient carry
the bounded endpoint state, that is, when the more severe of that coefficient's two endpoint states is
bounded; a coefficient in any other state keeps both endpoint values and both endpoint states, while the feasible coefficient vectors its searches
attain still enter the next round's starting values. Widen the identified ranges of the coefficients
on the lagged expected-SDF PC scores only by applying
equation~\eqref{eq:structural-recovery} to these accepted structural SDF-news coefficient vectors.
Leave unwidened any design coefficient whose loading vector \(\lambda_p\) has every entry no larger in
absolute value than the largest absolute entry of \(\beta_{2R}\) times
\(\sqrt{\epsilon_{\mathrm{mach}}}\), where \(\epsilon_{\mathrm{mach}}=2^{-52}\) is the binary64
machine epsilon, the distance from one to the next larger binary64 number. Such a coefficient is
\emph{point identified}, meaning its identified range is a single value: its loading is zero in exact
arithmetic, and a widening search would separate its endpoints by rounding error alone.

These rules certify only attained feasibility and the stated expanding-box behavior. A global
certificate would also require proof of the smallest or largest value or a mathematical direction
along which feasible coefficient vectors can grow without limit.

Certify a whole mean identified set, rather than a single endpoint, as follows. At the given slack,
solve the lower and upper identified endpoints of all three coordinate functionals, six searches in
all, and let the \emph{total width} be the sum of the three per-coordinate differences between the
upper and lower endpoints. The set is \emph{unreliable} when the total width is missing, and also
when every one of the six sides is finite but at least one of them fails its boundary check. Failing
that, the set is \emph{unbounded} when at least one side is infinite. Failing both, the set is
\emph{bounded}. Note the asymmetry: a set with an infinite side is unbounded even when some other
side failed its boundary check.

Evaluate the coarse slack sequence \(0,0.005,\ldots,0.99\), a finite sequence of evenly spaced slack
values. The critical slack \(\tau^\star\) is the
transition from a numerically certified bounded mean identified set to a numerically certified unbounded mean identified set. A
\emph{numerical certificate} is evidence that satisfies the stated finite-search rules rather than a
proof of global behavior. The coarse slack sequence \emph{brackets} such a transition when one end is
certified bounded and the other is certified unbounded. Ignore unreliable coarse-sequence values when
forming the bracket. Use the largest certified bounded slack below the smallest certified unbounded slack as the initial
lower endpoint, and that smallest certified unbounded slack as the initial upper endpoint. At a midpoint certified bounded or certified unbounded, replace the lower endpoint after a
bounded result and the upper endpoint after an unbounded result. Perform at most 40 bisections. The
procedure imposes no separate bracket-width tolerance. When all 40 bisections run, report the
midpoint of the final bracket, that is, one half of the sum of its lower and upper endpoints. If a midpoint is unreliable, leave the
bracket unchanged, stop immediately, and report that midpoint.
If no mean identified set is certified unbounded by slack 0.99, report the capped value 0.99.

\section{Variance shares implied by the mean equation}
\label{sec:shares}

On the identification sample, use \(L_t\in\mathbb R^3\) for the lagged expected-SDF PC scores and
\(Y_{2t}\in\mathbb R^3\) for the SDF-news PC scores. Define
\begin{equation*}
 C_L=\widehat{\operatorname{Var}}_T(L),\quad
 C_N=\widehat{\operatorname{Var}}_T(Y_2),\quad
 C_{LN}=\widehat{\operatorname{Cov}}_T(L,Y_2),\quad
 V_g=\widehat{\operatorname{Var}}_T(Y_1).
\end{equation*}
Here \(C_L\) and \(C_N\) are three-by-three sample covariance matrices,
\(C_{LN}\) is the three-by-three sample cross-covariance matrix, and \(V_g\) is the nonnegative
scalar sample variance of consumption growth.

Let \(b_L(\theta)\in\mathbb R^3\) be the last three entries of
\(\beta_1(\theta)\). The separate variance shares attribute variation to the lagged expected-SDF PC
scores or to the SDF-news PC scores. The joint variance share also includes their cross-covariance
term; it therefore differs from the sum of the two separate shares by twice that cross-covariance
term, which is nonzero in sample. Express the three shares as percentages:
\begin{align}
 \mathcal S_L(\theta)&=100\frac{b_L(\theta)'C_Lb_L(\theta)}{V_g},\notag\\
 \mathcal S_N(\theta)&=100\frac{\theta'C_N\theta}{V_g},\notag\\
 \mathcal S_{LN}(\theta)&=100\frac{b_L(\theta)'C_Lb_L(\theta)
 +2b_L(\theta)'C_{LN}\theta+\theta'C_N\theta}{V_g}.
 \label{eq:share-combined}
\end{align}
Here \(\mathcal S_L(\theta)\), \(\mathcal S_N(\theta)\), and
\(\mathcal S_{LN}(\theta)\) are scalar percentage variance shares for the lagged expected-SDF PC
scores, SDF-news PC scores, and both sets jointly. The vector \(b_L(\theta)\in\mathbb R^3\) contains
the coefficients on the lagged expected-SDF PC scores. The vector \(\theta\in\mathbb R^3\) contains
the structural SDF-news coefficients.

For PC score coordinate \(k\), the displayed contributions for individual PC scores are
\begin{equation*}
 \mathcal S_{L,k}(\theta)=100\frac{b_{L,k}(\theta)^2C_{L,kk}}{V_g},
 \qquad
 \mathcal S_{N,k}(\theta)=100\frac{\theta_k^2C_{N,kk}}{V_g}.
\end{equation*}
Here \(\mathcal S_{L,k}(\theta)\) and \(\mathcal S_{N,k}(\theta)\) are the scalar percentage
contributions of lagged expected-SDF PC score \(k\) and SDF-news PC score \(k\).

Off-diagonal covariances generally make the sum of the individual contributions differ from their
corresponding variance share. The \emph{largest-remainder display adjustment} aligns the rounded
individual contributions with the rounded variance-share total and leaves the algebraic shares
unchanged. Apply it as follows, in display units of \(0.01\) percentage point, to a block consisting
of one variance share and its three individual contributions.

Multiply each of the four values by 100. Round the block value to the nearest integer, resolving an
exact half to the nearest even integer. Round each of the three component values down to the next
integer no larger than it. Let the \emph{deficit} be the rounded block integer minus the sum of the
three rounded-down component integers. Add one to each of the deficit many components with the largest
fractional parts, taking them in decreasing order of fractional part and breaking ties by row order.
Divide all four integers by 100.

Apply this adjustment to the block of the lagged expected-SDF PC scores and to the block of the
SDF-news PC scores, and only to the two fixed-coefficient columns, namely the mean least-squares
benchmark column and the zero-slack point column. Never apply it to the joint variance share, whose
row is not the sum of the two block rows, and never to a positive-slack range column. Skip the
adjustment for a block in which any of the four values is nonfinite, and print those cells unadjusted.
The adjustment is defined only when the deficit is at least zero and at most three; stop the
reproduction if it is not.

For each display slack, approximate the ranges of all three variance shares over
\(\Theta(\tau)\) with a \(101\times101\times101\) \emph{rectangular lattice}, a finite evenly
spaced grid of structural SDF-news coefficient vectors spanning the finite
\emph{coefficient box}: the rectangle whose side in coordinate \(\ell\) runs from the accepted lower
to the accepted upper identified endpoint of coordinate functional \(\ell\) over \(\Theta(\tau)\) at
that slack. It is a different object from the expanding search boxes of
Section~\ref{sec:idset}. Retain lattice points with \(\rho_Q(\theta)\le10^{-10}\). Use as
starting values the retained points attaining the share minimum and maximum and each of the three
coordinate minima and maxima. After removing duplicates, this rule supplies at most eight starting
values. Stop the reproduction if the lattice retains no feasible point. If either side of the finite
coefficient box is nonfinite, do not compute the range at all: report both endpoints missing and print
the coefficient row's endpoint state in place of an interval.

Run both share minimization and share maximization from every starting value. Use sequential
least-squares quadratic programming with the analytic share gradient and the finite coefficient box as
the coordinate bounds. Divide the structural SDF-news coefficient vector and both coordinate bounds by
\(d_\theta\) and search in those scaled coordinates, leaving the share objective unscaled, so that the
relative iterate tolerance \(10^{-8}\) applies to the scaled coordinates; permit at most 1,000
objective evaluations per starting value. Evaluate \(\rho_Q\) at the unscaled search result, and
retain that result only when the result is finite and \(\rho_Q(\theta)\le10^{-4}\). Evaluate the
retained share at the search result after clamping each of its coordinates into the finite coefficient
box.

Form the lower endpoint of the attained range as the smallest of the retained lattice values together
with the retained minimization results, and the upper endpoint as the largest of the retained lattice
values together with the retained maximization results. A result from one search direction never
enters the endpoint of the other direction. The certificate is feasibility only, so a retained result
can sit on an active inequality or strictly inside the mean identified set.

Each variance share carries the most severe endpoint state among the three structural SDF-news
coefficient rows at that slack, in the severity order given above. Whenever both endpoints of all
four rows of a block are finite, two coherence conditions must hold: the block share's maximum is at least \(0.98\) times
the largest of its three components' maxima, and the block share's minimum is at least \(0.98\) times
the largest of its three components' minima less \(10^{-9}\). Stop the reproduction if either fails.
Nonnegativity of the shares is a consequence of their being quadratic forms in sample covariance
matrices and not a constraint imposed on the search; verify that the joint variance share is
nonnegative in the two fixed-coefficient columns and stop the reproduction if it is not.

For the contribution of individual PC score \(k\), map the corresponding coefficient's identified
range \([\underline b,\overline b]\) through \(b\mapsto100C_{kk}b^2/V_g\), where \(C_{kk}\) is
diagonal entry \(k\) of that block's covariance matrix \(C_L\) or \(C_N\). The upper endpoint of the
contribution is \(100C_{kk}\max(\underline b^2,\overline b^2)/V_g\). Its lower endpoint is
\(100C_{kk}\min(\underline b^2,\overline b^2)/V_g\) when \(\underline b>0\) or \(\overline b<0\), and
is zero when \(\underline b\le0\le\overline b\). The mapped contribution carries the endpoint state of
the coefficient row it came from. Use the identified ranges of the coefficients on the lagged
expected-SDF PC scores for the contributions of the lagged expected-SDF block, and the identified
ranges of the structural SDF-news coefficients for the contributions of the SDF-news block.

Label the variance-share ranges as attained numerical ranges. Reserve
\emph{certified global range} for a range supported by a global proof. A \emph{lattice
projection} is the collection of feasible values found on the stated finite lattice. The same
qualification applies to the three-dimensional mean identified set rendering and other lattice projections.
A local search seeks a smaller or larger objective value near a stated starting value and may miss
the smallest or largest value over the entire mean identified set.

\section{Mean-equation identification tests and checks}
\label{sec:mean-diagnostics}

These tests and checks describe identifying variation and leave the mean identified set
under the reported specification unchanged. \emph{Heteroskedasticity} means that a regression error's
conditional variance changes with observed variables.

Except for the first-order autoregressive conditional heteroskedasticity, or ARCH(1),
Lagrange-multiplier test, all the tests of this section share one null hypothesis:
that the conditional variance of the SDF-news PC score is unrelated to the instrument. Their common
alternative is that the instrument drives that conditional variance. The ARCH(1) Lagrange-multiplier
test has a different null: that the conditional variance is constant over time, against the
alternative that it clusters in time. The White, studentized Breusch--Pagan, Harvey, Glejser,
Breusch--Pagan Lagrange-multiplier, and ARCH(1) Lagrange-multiplier tests are one-sided and use the
upper tail of their reference distribution. The Goldfeld--Quandt test and the studentized Anscombe
test are two-sided.

For each SDF-news PC score \(Y_{2k,t}\), regress the score on an intercept and \(z_t\); call this the
\emph{base regression}. Here \(k\in\{1,2,3\}\) indexes the PC score coordinate. Let \(n_d\) denote
its observation count and \(u_{k,t}^{D}\) its residual. Every diagnostic in this section runs on the
identification sample with no further row deletion, so \(n_d=T=N_O\). Let \(\widehat Y_{2k,t}\) denote
the base regression's fitted value and let \(X_d\in\mathbb R^{n_d\times2}\) contain the intercept and instrument.
An
\emph{auxiliary regression} is a secondary regression used to calculate a test statistic rather
than to estimate the economic equation of interest.

For the White test, regress \((u_{k,t}^{D})^2\) on \(1,z_t,z_t^2\). If
\(u_{k,t}^{W}\) is the auxiliary
residual, calculate
\[
 \mathcal T_W=n_d\left\{1-\frac{\sum_t(u_{k,t}^{W})^2}
 {\sum_t[(u_{k,t}^{D})^2-\overline{(u_k^{D})^2}]^2}\right\}.
\]
Here \(\mathcal T_W\) is the scalar White statistic, \(u_{k,t}^W\) is the auxiliary-regression
residual, \(u_{k,t}^D\) is the base-regression residual, and \(n_d\) is the observation count. Use
the upper tail of a chi-squared distribution with two degrees of freedom. Because the regression has
one nonconstant predictor, these three terms exhaust its products and squares.

For the studentized Breusch--Pagan test, set
\(w_{k,t}^{BP}=(u_{k,t}^{D})^2-\overline{(u_k^{D})^2}\), project
\(w_k^{BP}\) by OLS on \(X_d\), and call the fitted vector \(\widehat w_k^{BP}\). Calculate
\[
 \mathcal T_{BP}=n_d\frac{\sum_t(\widehat w_{k,t}^{BP})^2}
 {\sum_t(w_{k,t}^{BP})^2}.
\]
Here \(\mathcal T_{BP}\) is the scalar studentized Breusch--Pagan statistic, \(w_{k,t}^{BP}\) is
the centered squared base-regression residual, and \(\widehat w_{k,t}^{BP}\) is its fitted value.
Use the upper tail of a chi-squared distribution with one degree of freedom. Studentized means
scaled by an estimated measure of variation.

For the Goldfeld--Quandt test, sort observations by \(z_t\) in ascending order and retain the
original row order among tied values. Let \(r\) be \(n_d/3\) rounded to the nearest integer, with an
exact half-integer rounded to the nearest even integer. Set \(n'=n_d-r\), let \(n_1\) be \(n'/2\)
rounded by the same rule, and set \(n_2=n'-n_1\). Use the first \(n_1\) and last \(n_2\) sorted
observations. If \(\min(n_1,n_2)\le2\), retain the central observations rather than removing them, so
that no observation is dropped and the two groups partition all \(n_d\) sorted rows. For even
\(n_d\),
use two groups of size \(n_d/2\); for odd \(n_d\), use group sizes \((n_d-1)/2\) and
\((n_d+1)/2\). If either resulting group has no more than two observations, stop the reproduction:
unlike the Glejser, Breusch--Pagan Lagrange-multiplier, and ARCH(1)
calculations, whose failures are caught and reported as a missing p-value, a
failure here is treated as a defect rather than as a verdict.

Fit the original two-parameter regression separately in the low and high groups. If group
\(g\in\{1,2\}\) has \(n_g\) rows and residual sum of squares \(RSS_g\), set
\(S_g^2=RSS_g/(n_g-2)\) and
\(\mathcal T_{GQ}=S_2^2/S_1^2\). Let \(\mathcal F_{q_2,q_1}(\cdot)\) be the cumulative distribution function of an \(F\)
distribution whose numerator degrees of freedom are \(q_2=n_2-2\), matching the numerator
\(S_2^2\), and whose denominator degrees of freedom are \(q_1=n_1-2\). The two-sided p-value is
\[
 \min\{2\mathcal F_{q_2,q_1}(\mathcal T_{GQ}),\;2[1-\mathcal F_{q_2,q_1}(\mathcal T_{GQ})]\}.
\]
Here \(\mathcal T_{GQ}\) is the scalar Goldfeld--Quandt statistic,
\(\mathcal F_{q_2,q_1}\) is the cumulative distribution function of an \(F\) distribution, and \(q_2\) and \(q_1\) are its positive integer numerator and denominator degrees of freedom.

For the Harvey heteroskedasticity test, regress
\(\ell_{k,t}^{H}=\log[(u_{k,t}^{D})^2]\) on \(X_d\). Let \(RSS_H\) be the residual sum of
squares from the auxiliary regression and let
\(TSS_H=\sum_t(\ell_{k,t}^{H}-\bar\ell_k^{H})^2\) be its total sum of squares. Calculate
\[
 \mathcal T_H=(TSS_H-RSS_H)/\psi_1(1/2).
\]
Here \(\mathcal T_H\) is the scalar Harvey statistic, and \(TSS_H\) and \(RSS_H\) are the total
and residual sums of squares. The quantity \(\psi_1(1/2)\) is the trigamma function evaluated at one
half. The trigamma function is the second derivative of the logarithm of the gamma function, which
is the continuous extension of the factorial.
Use the upper tail of a chi-squared distribution with one degree of freedom.

For the Glejser test, regress \(|u_{k,t}^{D}|\) on \(X_d\). If \(RSS_G\) and \(TSS_G\) are the
residual and total sums of squares of \(|u_{k,t}^{D}|\), calculate
\[
 \mathcal T_G=\frac{TSS_G-RSS_G}
 {[\sum_t(u_{k,t}^{D})^2/n_d](1-2/\pi)}.
\]
Here \(\mathcal T_G\) is the scalar Glejser statistic; \(TSS_G\) and \(RSS_G\) are the total and
residual sums of squares; and \(\pi\) is the circle constant. Use the upper tail of a chi-squared
distribution with one degree of freedom.

Decide once, for all three SDF-news PC scores together, whether the Anscombe test runs. It runs for
all three when the ratio below exceeds \(10^{-3}\) and for none of them otherwise; when it does not
run, the reported table carries no Anscombe row and its row count falls by one. The ratio is
\[
\max_k\left\{\frac{\operatorname{sd}_{n_d-1}(\widehat Y_{2k})}
{\operatorname{sd}_{n_d-1}(Y_{2k})}\right\}.
\]
Here the displayed quantity is the largest ratio, across the three SDF-news PC scores, of the fitted values'
ordinary sample standard deviation to the outcome's ordinary sample standard deviation.

Define the residual-maker matrix and weighted fitted-value mean by
\[
 P_D^\perp=I-X_d(X_d'X_d)^{-1}X_d'\in\mathbb R^{n_d\times n_d},
 \qquad
 \bar f_{k,w}=\frac{\sum_t P_{D,tt}^\perp\widehat Y_{2k,t}}{n_d-2}.
\]
Here \(P_D^\perp\in\mathbb R^{n_d\times n_d}\) is the residual-maker matrix,
\(\bar f_{k,w}\) is the scalar weighted fitted-value mean, and
\(I\in\mathbb R^{n_d\times n_d}\) is the identity matrix. The residual-maker matrix maps each
SDF-news PC-score vector to its base-regression residual vector.
The studentized Anscombe statistic is
\[
 \mathcal T_A=\frac{\sum_t(u_{k,t}^{D})^2(\widehat Y_{2k,t}-\bar f_{k,w})}
 {\left\{[\sum_t(\widehat Y_{2k,t}-\bar f_{k,w})^2]
 [\sum_t\{(u_{k,t}^{D})^2-\overline{(u_k^{D})^2}\}^2]/(n_d-2)\right\}^{1/2}}.
\]
Here \(\mathcal T_A\) is the scalar studentized Anscombe statistic,
\(\widehat Y_{2k,t}\) is the base-regression fitted value, and \(\bar f_{k,w}\) is its weighted mean.
Let \(\Phi\) denote the standard-normal cumulative distribution function. The p-value is
\(2\{1-\Phi(|\mathcal T_A|)\}\).

Also conduct two auxiliary Lagrange-multiplier (LM) tests. An LM test obtains its statistic from an
auxiliary regression under the null restriction. For the Breusch--Pagan LM test, regress
\((u_{k,t}^{D})^2\) on \(1,z_t\). The coefficient of determination is the fraction of the
outcome's sample variation explained by its fitted values. Denote this fraction by
\(R^2_{BP,k}\), calculate \(LM_{BP,k}=n_dR^2_{BP,k}\), and use a chi-squared distribution with one
degree of freedom. For the ARCH(1) LM test, regress \((u_{k,t}^{D})^2\) on
\(1,(u_{k,t-1}^{D})^2\) over the \(n_d-1\) available pairs. If its coefficient of
determination is \(R^2_{ARCH,k}\), calculate \(LM_{ARCH,k}=(n_d-1)R^2_{ARCH,k}\) and again use a
chi-squared distribution with one degree of freedom.

For each SDF-news PC score, square the score and report
\(\widehat{\operatorname{Cov}}_T(z,Y_{2k}^2)\), the centered divisor-\(T\) covariance of
equation~\eqref{eq:covT}, together with the sample correlation of \(Y_{2k}^2\) with the instrument,
which does not depend on the choice of divisor. Also report the least-squares \(t\)-statistic for the instrument coefficient in the base
regression and \(\operatorname{corr}_T(W_1,Y_{2k})\). Last, report the two normalized quantities
\[
 \frac{\widehat{\operatorname{Cov}}_T(W_1,Y_{2k})}
 {\widehat{\operatorname{Var}}_T(Y_{2k})}
 \frac{\operatorname{sd}_{T-1}(Y_{2k})}{\operatorname{sd}_{T-1}(Y_1)},
 \qquad
 \frac{\widehat{\operatorname{Cov}}_T(W_1,Y_{2k})}
 {\widehat{\operatorname{Var}}_T(Y_{2k})}\operatorname{sd}_{T-1}(Y_{2k}).
\]
Here both displayed quantities are scalar normalized covariances for SDF-news PC score \(k\), reported
in the order shown; the first additionally divides by the ordinary sample standard deviation of
consumption growth.

The reproduction also forms the \emph{unreported relevance summary}, a count of how many of the three
SDF-news PC scores show
instrument-driven conditional heteroskedasticity. It computes this summary, but writes it into no reported output: the summary reaches a displayed exhibit only if a reader carries it there. Only
three p-values enter it: those of the studentized Breusch--Pagan test, the
Goldfeld--Quandt test, and the ARCH(1) LM test. A score \emph{rejects} when at least one of those
three p-values is finite and strictly below five percent. A score counts as \emph{tested} when at
least one of those three p-values is finite; a nonfinite p-value is neither a rejection nor a
non-rejection. Let \(n_{\mathrm{tested}}\) be the number of tested scores and \(n_{\mathrm{reject}}\)
the number of rejecting scores. When \(n_{\mathrm{tested}}=0\), the summary records that relevance is
undetermined. When \(n_{\mathrm{reject}}>0\), it records that the instrument drives significant
conditional heteroskedasticity in \(n_{\mathrm{reject}}\) of \(n_{\mathrm{tested}}\) scores. When
\(n_{\mathrm{tested}}>0\) and \(n_{\mathrm{reject}}=0\), it records no significant conditional
heteroskedasticity, that is, weak relevance. When fewer than three scores are tested, the summary also
reports how many of the three did not run. These calculations use SDF-news PC scores, whereas
identification uses the reported specification's SDF-news reduced-form residuals.

The \emph{one-direction relevance test} asks whether variation related to the instrument is present
in the least-moved direction of the matrix \(\widehat M_z\):
\begin{equation*}
 \widehat M_z=\frac1T\sum_{t=1}^{T}(z_t-\bar z)Y_{2t}Y_{2t}'
 \in\mathbb R^{3\times3}.
\end{equation*}
Here \(\widehat M_z\in\mathbb R^{3\times3}\) is the instrument-weighted second-moment matrix,
\(z_t-\bar z\) is the centered instrument, and \(Y_{2t}\in\mathbb R^3\) is the SDF-news
PC score vector.

Let \(u_z\in\mathbb R^3\) be an eigenvector with Euclidean norm one associated with the eigenvalue of
\(\widehat M_z\) having the smallest absolute value. Call \(u_z\) the \emph{least-moved direction}
because the matrix changes that direction by the smallest absolute eigenvalue. Set
\(d_t^J=(u_z'Y_{2t})^2\). With
\(s_t^J=(z_t-\bar z)(d_t^J-\bar d^J)\), compute
\begin{equation*}
 J=T(\bar s^J)^2/\widehat{\operatorname{LRV}}(s^J).
\end{equation*}
Here \(J\) is the nonnegative scalar one-direction relevance statistic, and \(\bar s^J\) is the
sample mean of \(s_t^J\). The quantity \(\widehat{\operatorname{LRV}}(s^J)\) is its estimated
long-run variance; the letters LRV stand for long-run variance.
It is a Bartlett long-run variance with \(K_z=\lfloor4(T/100)^{2/9}\rfloor\) lags and combines
the sequence's contemporaneous variance with its lagged covariances. Specifically, with
\(\widetilde s_t^J=s_t^J-\bar s^J\), let
\[
 \widehat C_j=\frac1T\sum_{t=1}^{T-j}\widetilde s_t^J\widetilde s_{t+j}^J,
 \qquad
 \widehat{\operatorname{LRV}}(s^J)=\widehat C_0+
 2\sum_{j=1}^{K_z}\left(1-\frac{j}{K_z+1}\right)\widehat C_j.
\]
Here \(\widehat C_j\) is the scalar lag-\(j\) sample covariance of the centered sequence
\(\widetilde s_t^J\), and \(K_z\) is the positive integer lag count. The second expression is the
scalar Bartlett long-run variance estimate.

Compare \(J\) with a chi-squared distribution having one
degree of freedom. The statistic \(J\) tests relevance in one direction rather than the matrix's full
rank. It also
uses SDF-news PC scores. Report the determinant, condition number, and smallest singular value of
\(\widehat M_z\) beside its p-value. The procedure supplies no separate boundary rule for a nonpositive
long-run variance or a nonunique least-moved direction. Either condition leaves the chi-squared
comparison numerically uncertified.
The matrix \(\widehat M_z\) is symmetric, so its singular values are the absolute values of its
eigenvalues. The determinant is a scalar that is zero when a square matrix is singular; here it is
the product of the signed eigenvalues of \(\widehat M_z\). The condition number is the ratio of the
largest absolute eigenvalue to the smallest. The smallest singular value is the smallest absolute
eigenvalue and measures the matrix's weakest linear direction.

\section{Conditional-scale sample}
\label{sec:logvar-scale}

Join the identification sample to the four lagged supplied return scores by quarter and retain the
matched rows in increasing quarter order. These rows form the \emph{conditional-scale sample}. Let the positive integer \(T_v\) be its
observation count. The fixed inputs retain \(T_v=255\) quarters, from 1962Q2 through 2025Q4.
The join deletes no row on account of a nonfinite value, so the sample definition above is exact. A
nonfinite value carried into the conditional-scale design makes the linear system that defines the
conditional-scale estimator mapping unsolvable, and the reproduction stops there. Estimate the two reduced forms on the identification
sample before the join, then restrict the
resulting \(W_{1t}\) and \(W_{2t}\) to the matched rows.

On these \(T_v\) quarters,
center each lagged supplied return score and retain its original scale. Before the PPML, Harvey QML, and LAD conditional-scale searches, but not before the log-OLS search, which runs first and does not
perform this check, verify that every centered score's absolute sample mean is strictly below
\(10^{-8}\); stop the reproduction if any is at least that value. Define
the \emph{conditional-scale design vector}
\begin{equation*}
 X_t^V=(1,\widetilde r_{1t},\widetilde r_{2t},\widetilde r_{3t},
          \widetilde r_{4t})'\in\mathbb R^5.
\end{equation*}
Here \(X_t^V\in\mathbb R^5\) is the conditional-scale design vector containing an intercept and the
four centered lagged supplied return scores. Let \(X^V\in\mathbb R^{T_v\times5}\) stack these row vectors.

For a candidate structural SDF-news coefficient vector \(\theta\in\Theta(\tau)\), restrict
equation~\eqref{eq:mean-residual} to these rows and define the nonnegative squared structural mean residual
\begin{equation*}
 q_t(\theta)=e_t(\theta)^2.
\end{equation*}
Here \(q_t(\theta)\) is the nonnegative scalar squared structural mean residual for quarter \(t\).

The \emph{conditional-scale coefficient} is
\(\eta=(\eta_1,\eta_2,\eta_3,\eta_4,\eta_5)'\in\mathbb R^5\), indexed \(j\in\{1,\ldots,5\}\)
throughout. Its first element is an intercept,
and its last four elements relate the lagged supplied return scores to the estimator-specific log
scale.

Each conditional-scale result has a \emph{mean least-squares benchmark reference column}, a
\emph{zero-slack point column}, and \emph{positive-slack mean-identified-set columns}. The mean
least-squares benchmark reference column takes the mean least-squares benchmark residuals estimated on
the identification sample, restricts them to the \(T_v\) matched rows, and regresses their log squares
on the conditional-scale design over those rows.
The zero-slack point column uses
\(e_t(\widehat\theta_0)\), where \(\widehat\theta_0\) is the accepted solution of
equation~\eqref{eq:tau-zero-system}. Positive-slack mean-identified-set columns apply the chosen conditional-scale
estimator mapping to every accepted \(\theta\in\Theta(\tau)\). A
\emph{conditional-scale estimator mapping} assigns a conditional-scale coefficient vector to each
structural SDF-news coefficient vector.

At each slack, first read the endpoint states of the mean identified set's three structural
SDF-news coefficients. Give each coefficient the more severe of its two endpoint states in the
severity order bounded, unbounded, unreliable. If any of the three is not bounded, perform no
lattice construction and no search at that slack:
report every conditional-scale endpoint value as missing, and assign every conditional-scale
endpoint the unbounded endpoint state when at least one of the three is unbounded, and the
unreliable endpoint state otherwise. This rule belongs to the stage rather than to any one
estimator, so it governs all four conditional-scale estimator mappings, LAD included.

Unless a subsection gives a different rule, construct the conditional-scale lattice at each slack
with 41 equally spaced values along each structural SDF-news coefficient, with the first coefficient
changing fastest. Retain points satisfying \(\rho_Q(\theta)\le10^{-10}\). If fewer than 100 points
remain, replace the lattice once with 81 values along each coefficient. When a calculation has a
lattice-point cap \(C\) and \(m>C\) points remain, set \(h_C=\lceil m/C\rceil\) and retain points
numbered \(1,1+h_C,1+2h_C,\ldots\). Here \(C\) and \(m\) are positive integers, and \(h_C\) is the
positive integer thinning interval. Append the feasible search start point even when it duplicates a
retained point. The search start point is the zero-slack point when that point has no missing coordinate and
satisfies every inequality of the baseline-slack system with every inequality scale
\(\omega_k\) set to one, that is when \(\max_k g_k(\theta)\le0\), a stricter test than the lattice
admission tolerance above. Otherwise the search start point is the first point of the
baseline-slack lattice retained after thinning, which is not in general the current slack's lattice,
except where a subsection states that its estimator has no fallback search start point and leaves the search start point missing.

For PPML and Harvey QML, visit the retained lattice points in greedy nearest-neighbor order. Begin
at the retained point with the smallest squared Euclidean distance to the search start point; from the
current point move to the nearest unvisited retained point; break either tie toward the point that
comes first in lattice order. After each accepted fit, pass its conditional-scale coefficient vector
to the next point as that point's supplied warm starting value. Each of these two estimator mappings
keeps one \emph{fit record}---its store of previously evaluated arguments and their fit results---across
all display slacks, so a fit already recorded at an earlier slack is reused, and rebuilds its
\emph{fresh-fit allowance}---the maximum number of fits at arguments absent from that record---at each
slack. At every slack after the first,
additionally evaluate each attaining structural SDF-news coefficient vector of a preceding slack's
endpoints that received the bounded endpoint state there; skip any such vector that is infeasible or
whose fit is not accepted, and let an accepted fit whose mapped value is more extreme than the
current endpoint replace that endpoint before the refinement begins. Log-OLS evaluates its retained
lattice in lattice order, uses no warm starting values, keeps no fit record across slacks, and
receives no preceding-slack starting values.

For conditional-scale coefficient \(j\in\{1,\ldots,5\}\), evaluate its conditional-scale estimator mapping at every
lattice point retained after thinning that has an accepted fit. Let
\(v_{j,\min}^{\mathrm{lat}}\) and \(v_{j,\max}^{\mathrm{lat}}\) be the smallest and largest finite
mapped values. The lattice points that attain these values are the lower and upper
\emph{attained lattice extremes}. Number the retained lattice points in lattice order, the order in
which the first structural SDF-news coefficient changes fastest. If several points attain the same
value, select the one with the smallest such number. PPML, Harvey QML, and LAD instead use the
\emph{coordinate-by-coordinate tie rule} for ties among attained extreme mapped values; that rule
compares the two tied points'
coordinates from first to last and selects the one with the smaller coordinate at the first
coordinate at which they differ. Define the positive scalar lattice-extreme scale
\begin{equation*}
 g_j=\max\{1,|v_{j,\min}^{\mathrm{lat}}|,|v_{j,\max}^{\mathrm{lat}}|\},
\end{equation*}
taking the maximum over whichever of the two lattice extremes is finite, so that \(g_j=1\) when
neither is.

The \emph{conditional-scale endpoint-refinement procedure} refines one coefficient \(j\) and one
endpoint side at a time. Its starting values are the corresponding attained lattice extreme, then the
search start point when that point has no missing coordinate, then any preceding-slack endpoint arguments the
subsection supplies. A subsection may state in addition that its estimator supplies separated
lattice-derived starting values; where it does, take the retained lattice points in decreasing order
of endpoint extremeness and accept a point only when it lies farther than five percent of the
coefficient box's Euclidean diagonal from every already-accepted point, stopping at the number of
accepted points the invoking subsection states. The first acceptance is the attained lattice extreme
itself, which is already a starting value, so a count of three supplies two further starting values and three
lattice-derived starting values in all. The primary refinements of PPML, Harvey QML, and LAD use three; their
audits use five, as their subsections state. Log-OLS uses no such pool, so for it the attained lattice extreme is the only
lattice-derived starting value.

For log-OLS, PPML, and Harvey QML, refine each lower and upper endpoint by sequential least-squares
quadratic programming with the analytic objective gradient. For LAD, use constrained optimization by
linear approximations, a derivative-free procedure that approximates the objective and constraints
locally by linear functions. Divide the structural SDF-news coefficient vector and both coordinate
bounds by \(d_\theta\) and search in those scaled coordinates, restricting each scaled coordinate to
\([-10^6,10^6]\); clamp the starting value into that scaled box before the search begins. Use relative
iterate tolerance \(10^{-8}\) in the scaled coordinates and permit at most 1,000 objective evaluations
per refinement. A search that ends in an error yields no result for that starting value rather than
stopping the reproduction.

Retain a finite refined value only when \(\rho_Q(\theta)\le10^{-4}\). A retained value whose objective
for coefficient \(j\) is not finite, or whose absolute value exceeds \(5g_j\), is discarded and marks
that side of coefficient \(j\) unreliable. A retained value replaces the endpoint only when it improves it: the lower endpoint
becomes the smallest of the attained lattice value and every retained minimization result, and the
upper endpoint the largest of the attained lattice value and every retained maximization result. If no
refined value is retained on a side, that side is unreliable and the lattice value alone is not
reported. Finally, if the point attaining a side sits outside the finite coefficient box by more than
\(10^{-4}\) relative to the largest of the box's span, the absolute values of its two sides, and one,
that side is unreliable. These rules certify finite feasibility, not a global optimum.

Two further acceptance rules apply across slacks once every slack has been solved. The
\emph{nesting rule} requires each identified range not to shrink as the slack rises, and the
\emph{containment rule} requires the accepted zero-slack point to lie inside every positive-slack
range. Both compare at relative tolerance \(10^{-6}\), and a coefficient side that fails either takes
the unreliable endpoint state.

Four further conditions instead assign the unreliable endpoint state to every conditional-scale
endpoint at a slack and report every endpoint value at that slack as missing. They are: no lattice
point satisfies \(\rho_Q(\theta)\le10^{-10}\); for PPML or Harvey QML, at least one fit failed at a
scanned lattice point; no scanned lattice point produced a successful fit; and the fresh-fit
allowance for that slack was exhausted before the slack finished.

\section{Logarithmic least squares (log-OLS)}
\label{sec:logols}

Log-OLS regresses the log squared structural mean residual on the five conditional-scale design
terms. It provides a direct linear conditional-scale estimator mapping, but becomes undefined at a zero structural mean
residual. When every \(e_t(\theta)\ne0\), define
\begin{equation*}
 \ell_t(\theta)=\log q_t(\theta)=2\log|e_t(\theta)|.
\end{equation*}
Here \(\ell_t(\theta)\) is the scalar log squared structural mean residual for quarter \(t\). The
log-OLS conditional-scale estimator mapping is
\begin{equation*}
 \widehat\eta_L(\theta)=(X^{V\prime}X^V)^{-1}X^{V\prime}\ell(\theta).
\end{equation*}
Here \(\widehat\eta_L(\theta)\in\mathbb R^5\) is the log-OLS conditional-scale coefficient vector,
\(\ell(\theta)\in\mathbb R^{T_v}\) stacks the log squared structural mean residuals, and
\(X^{V\prime}X^V\) must be nonsingular.

The Jacobian measures how the five log-OLS coefficients change with the three structural SDF-news
coefficients:
\begin{equation*}
 \frac{\partial\widehat\eta_L(\theta)}{\partial\theta'}
 =-(X^{V\prime}X^V)^{-1}X^{V\prime}\operatorname{diag}
   \left(\frac{2}{e_t(\theta)}\right)W_2.
\end{equation*}
Here \(W_2\in\mathbb R^{T_v\times3}\) stacks the SDF-news reduced-form residual vectors. The
derivative matrix has dimension \(5\times3\).

A \emph{structural-mean-residual hyperplane} is the set of \(\theta\) for which \(e_t(\theta)=0\) in
one quarter \(t\). If the mean identified set crosses one of these hyperplanes, the corresponding log
squared structural mean residual diverges. The endpoint
search explicitly checks for these crossings and assigns unbounded or unreliable endpoint states to
the affected coefficients.

The direction of divergence follows from the sign of the log-OLS weight on the crossing quarter. Let
\begin{equation*}
 P=(X^{V\prime}X^V)^{-1}X^{V\prime}\in\mathbb R^{5\times T_v},
\end{equation*}
so that \(\widehat\eta_L(\theta)=P\ell(\theta)\), and let \(\mathcal X\) be the set of quarters whose
structural-mean-residual hyperplanes the mean identified set crosses. Here \(P_{jt}\) is the scalar
weight that log-OLS coefficient \(j\) places on quarter \(t\), and \(\mathcal X\) is a finite set of
quarter indices. Because \(\ell_t(\theta)\to-\infty\) as \(e_t(\theta)\to0\), assign log-OLS
coefficient \(j\) the unbounded endpoint state on its lower side when \(P_{jt}>0\) for some
\(t\in\mathcal X\), and on its upper side when \(P_{jt}<0\) for some \(t\in\mathcal X\). A
coefficient can receive the unbounded endpoint state on both sides.

Build the lattice by the general rule stated above, and scan the retained points in groups of 5,000.
For each log-OLS coefficient, record the lattice points attaining its smallest and largest values.
Start the lower-endpoint search from the lower attained lattice extreme and the upper-endpoint search from the
upper attained lattice extreme. Also start each side from the accepted zero-slack structural SDF-news
coefficient vector whenever that vector is available, whether or not it is feasible; feasibility
governs only whether that vector is also appended to that slack's retained lattice, judged against
the current slack's inequalities with every inequality scale \(\omega_k\) set to one. Log-OLS has no fallback
search start point: when the zero-slack vector has a missing coordinate, the search start point is simply missing.

Combine the lattice's structural-mean-residual sign changes with a
\emph{hyperplane-intersection check}. For quarter \(t\), let \(\underline h_t\) and \(\overline h_t\)
be the accepted lower and upper endpoints of \(W_{2t}'\theta\) over the mean identified set. Define
the positive scalar tolerance
\begin{equation*}
 \delta_t=10^{-8}\max\{1,|W_{1t}|\}.
\end{equation*}
After first screening the range of \(W_{2t}'\theta\) over the enclosing coefficient box, declare an
intersection with \(W_{2t}'\theta=W_{1t}\) when
\begin{equation*}
 \underline h_t-\delta_t\le W_{1t}\le\overline h_t+\delta_t.
\end{equation*}
If \(W_{2t}=0\), declare a crossing for quarter \(t\) directly and solve nothing. Such a quarter
passes the box screen only when \(W_{1t}=0\), in which case its structural mean residual is zero at
every \(\theta\).
Apply the box screen exactly, without the tolerance \(\delta_t\), which widens only the comparison
against the accepted functional endpoints. A sign change of the structural mean residual across the
retained lattice, or an exact zero at a retained lattice point, also establishes a crossing. The
endpoint comparison is exact when the mean
identified set is connected, meaning that any two points can be joined by a continuous path inside
the set. For a disconnected set, it favors declaring a crossing. If either endpoint solve for
\(W_{2t}'\theta\) fails to return a certified bound for some quarter \(t\), the check is unresolved:
perform no search at that slack, report every conditional-scale endpoint value as missing, and
assign every conditional-scale endpoint at that slack the unreliable endpoint state.

For the mean least-squares benchmark reference column only, let \(\ell_t^{\mathrm{ref}}\) be the
log squared mean least-squares benchmark residual for quarter \(t\) on the conditional-scale rows,
and stop the reproduction if any of those benchmark residuals is exactly zero. In this column
\(\widehat\eta_L\) denotes the log-OLS coefficient vector fitted to \(\ell^{\mathrm{ref}}\) rather
than to \(\ell(\theta)\) at any \(\theta\). Let
\(\widehat u_t^L=\ell_t^{\mathrm{ref}}-X_t^{V\prime}\widehat\eta_L\) and
\(s_t^L=X_t^V\widehat u_t^L\in\mathbb R^5\). Apply neither prewhitening---the removal of fitted
first-order dependence before the long-run sum, defined among the analytic covariance choices of
Part~II---nor a finite-sample adjustment.
Report
\[
 \widehat V_L=(X^{V\prime}X^V)^{-1}\widehat M_4(s^L)
 (X^{V\prime}X^V)^{-1}.
\]
Here \(\widehat V_L\in\mathbb R^{5\times5}\) is the log-OLS coefficient covariance matrix.
For coefficient \(j\), divide \(\widehat\eta_{L,j}\) by the square root of diagonal element \(j\)
of \(\widehat V_L\) and calculate its two-sided p-value from a Student \(t\) distribution with
\(T_v-5\) degrees of freedom. Leave the statistic and p-value entries blank for the zero-slack point
column and positive-slack mean-identified-set columns.

\section{Exponential conditional-variance estimators}
\label{sec:expcondvar}

Both conditional-scale estimator mappings in this section use
\begin{equation*}
 v_t(\eta)=\exp(X_t^{V\prime}\eta)>0.
\end{equation*}
Here \(v_t(\eta)>0\) is the scalar fitted conditional variance for quarter \(t\) at
conditional-scale coefficient vector \(\eta\in\mathbb R^5\). The fitted conditional standard
deviation is \(\sqrt{v_t(\eta)}=\exp(X_t^{V\prime}\eta/2)\). Use the full exponent without truncation. For PPML, a
nonfinite fitted value causes that fit to fail. For Harvey QML the guard acts on the linear
predictor and fires before any fitted variance is formed: reject a trial conditional-scale
coefficient vector when any \(X_t^{V\prime}\eta\) exceeds the logarithm of the largest finite
binary64 number.

Neither exponential conditional-variance estimator's own search ever produces the unbounded endpoint
state: every endpoint it returns is bounded or unreliable. The shared rule that reads the
mean-equation endpoint states before any search can still assign the unbounded state to all of their
endpoints at a slack.

\subsection{Poisson pseudo-maximum likelihood (PPML)}

The PPML estimator represents the conditional variance with an exponential function and permits the
squared structural mean residual's distribution to differ from a Poisson distribution.
For fixed \(\theta\), choose \(\widehat\eta_P(\theta)\in\mathbb R^5\) to solve
\begin{equation*}
 \sum_{t=1}^{T_v}X_t^V\{q_t(\theta)-v_t(\eta)\}=0.
\end{equation*}
Here \(\widehat\eta_P(\theta)\in\mathbb R^5\) is the PPML conditional-scale coefficient vector,
\(q_t(\theta)\) is the squared structural mean residual, and \(v_t(\eta)\) is the fitted conditional
variance.

The corresponding information matrix is
\begin{equation*}
 A_P(\eta)=\sum_{t=1}^{T_v}v_t(\eta)X_t^VX_t^{V\prime}.
\end{equation*}
The information matrix \(A_P(\eta)\in\mathbb R^{5\times5}\) is symmetric and measures the local
sensitivity of the PPML estimating equation to changes in the conditional-scale coefficient vector.
Allow exact zero squared structural mean residuals, but require at least one positive squared
structural mean residual.

Before applying the conditional-scale estimator mapping to the mean identified set, perform the
\emph{PPML normalization check}. Define the \emph{normalization
anchor} as the accepted zero-slack structural SDF-news coefficient vector when it is feasible at the
baseline slack of 0.05. Otherwise, use the first lattice point retained after thinning at that slack. The
accepted PPML fit at this vector is the \emph{normalization-anchor fit}. Use a PPML normalization
divisor of one for the normalization anchor and the first ten lattice points retained after thinning at
the baseline slack, excluding the normalization anchor when it appears again. Trigger PPML normalization if any of these fits does not converge. Also
trigger it if the normalized information condition number is nonfinite or exceeds \(10^{10}\).
Finally, trigger it if the largest linear predictor lies within 5 of the largest argument for which
the exponential calculation remains finite. Here the linear predictor is the scalar
\(X_t^{V\prime}\eta\). If triggered, divide
every squared structural mean residual by the median positive squared structural mean residual at
the normalization anchor and use that positive PPML normalization divisor unchanged for every
subsequent PPML fit. Stop if the normalization anchor has no positive squared structural mean
residual. After fitting the
normalized squared structural mean residuals, add the logarithm of the PPML normalization divisor
to the intercept; only the intercept changes.

For fixed normalized squared structural mean residuals \(\breve q_t\), try PPML starting values in this order.
First, use the supplied warm starting value when available. Second, use the accepted zero-slack PPML fit when
available; otherwise use the accepted normalization-anchor fit. Third, use
\begin{equation*}
 \left(\log\bar q,0,0,0,0\right)',
 \qquad \bar q=\frac1{T_v}\sum_{t=1}^{T_v}\breve q_t.
\end{equation*}
Here \(\bar q>0\) is the mean normalized squared structural mean residual; the third starting value
exists only when that mean is strictly positive. Skip any candidate starting value whose
coordinates are not all finite or whose implied fitted conditional variances are not all finite, and
continue with the next candidate. Last, initialize fitted
variances to \(\breve q_t+0.1\) and linear predictors to \(\log(\breve q_t+0.1)\). Use at most 100
iteratively reweighted updates. If \(D^{(m)}\) is the quasi-Poisson deviance after update \(m\),
declare convergence when
\begin{equation*}
 \frac{|D^{(m)}-D^{(m-1)}|}{0.1+|D^{(m)}|}<10^{-10}.
\end{equation*}
Let \(\breve v_t^{(m)}>0\) be the fitted normalized conditional variance after update \(m\). Here
\(D^{(m)}\) is twice the sum of
\(\breve q_t\log\{\breve q_t/\breve v_t^{(m)}\}-(\breve q_t-\breve v_t^{(m)})\), with the first term
set to zero when \(\breve q_t=0\).
An iteratively reweighted update solves a weighted least-squares approximation to the nonlinear
estimating equation and refreshes the weights at the new coefficient.

Before fitting, retain the design rows with positive squared structural mean residuals and divide
each nonzero design column by its Euclidean norm. If the singular values of this scaled matrix are
\(s_1\ge\cdots\ge s_5\), define its \emph{numerical rank} as
\begin{equation*}
 \#\{j:s_j>10^{-10}s_1\}.
\end{equation*}
Here the rank is the number of singular values satisfying the displayed inequality. Require rank
five. Accept a fit only when all coefficients and fitted conditional variances are finite and every
fitted variance is positive. Also require convergence and reject a boundary solution, meaning a
final update that required step halving to restore finite deviance or positive finite fitted
variances. Finally, require the largest
coordinate-normalized estimating-equation residual to be at most \(10^{-8}\) and the normalized
information matrix's reciprocal condition number to be at least \(10^{-10}\). A reciprocal
condition number near zero indicates numerical singularity, meaning the matrix cannot be inverted
reliably. Allow a positive squared structural mean residual to approach zero.

Normalize the PPML acceptance quantities as follows. Let \(\breve q_t\) and \(\breve v_t\) denote the squared
structural mean residual and fitted conditional variance after PPML normalization. Let
\(X_{tj}^V\) be entry \(j\) of \(X_t^V\) and set
\(q_{\mathrm{med}+}=\operatorname{med}(\breve q_t:\breve q_t>0)\). For coordinate \(j\), divide the estimating-equation residual by
\begin{equation*}
 c^P_j=\max(1,q_{\mathrm{med}+})\sum_{t=1}^{T_v}|X_{tj}^V|.
\end{equation*}
Here \(c_j^P>0\) is the scalar normalization denominator for coordinate \(j\),
\(\breve q_t\) is the normalized squared structural mean residual, and \(X_{tj}^V\) is design entry
\((t,j)\). The estimating-equation acceptance rule is
\(\max_j|\sum_tX_{tj}^V(\breve q_t-\breve v_t)|/c^P_j\le10^{-8}\). Let
\(d^P_j=\{\sum_t\breve v_t(X_{tj}^V)^2\}^{1/2}\) and the diagonal matrix
\(D_P=\operatorname{diag}(d^P_1,\ldots,d^P_5)\in\mathbb R^{5\times5}\). The conditioning acceptance rule applies to
\(D_P^{-1}\{\sum_t\breve v_tX_t^VX_t^{V\prime}\}D_P^{-1}\).
Here \(d_j^P>0\) is the information-column scale and \(D_P\) is the diagonal matrix containing the
five values \(d_1^P,\ldots,d_5^P\).

Differentiate the estimating equation with respect to the structural SDF-news coefficient vector,
regarding the fitted conditional-scale coefficient as a function of that vector. That differentiation gives
\begin{equation*}
 \frac{\partial\widehat\eta_P(\theta)}{\partial\theta'}
 =A_P(\widehat\eta_P)^{-1}
   \sum_{t=1}^{T_v}X_t^V\{-2e_t(\theta)W_{2t}'\}.
\end{equation*}
Here the displayed Jacobian is a \(5\times3\) matrix measuring the response of the PPML
conditional-scale coefficient vector to the structural SDF-news coefficient vector.

The reciprocal condition number must again be at least \(10^{-10}\); when it is not, no analytic
Jacobian is returned at that structural SDF-news coefficient vector. The refinement procedure is
unchanged by that outcome, because the choice between the gradient-based and the derivative-free
procedure is made once per coefficient from whether the conditional-scale estimator mapping supplies
an analytic Jacobian at all, before any fit runs. Reusing a recorded fit consumes none of the fresh-fit allowance. The primary PPML conditional-scale estimator mapping uses at most 4,000 lattice points and a
fresh-fit allowance of 20,000 at each slack. An independent \emph{PPML coverage audit} uses up to 8,000 lattice
points and a fresh-fit allowance of 40,000 at each slack. For each coefficient and endpoint side, rank
lattice points from most to least extreme. Retain the extreme point and up to four further points whose
Euclidean distance from every retained point exceeds five percent of the coefficient box's Euclidean
diagonal. Use these five separated lattice-derived starting values plus the recorded search start point. The audit
uses one new fit record across display slacks and a new fresh-fit allowance at each slack.

For this PPML audit only, replace lattice thinning with a space-filling selection. Normalize each
feasible lattice coordinate to the unit interval using that coordinate's lattice range, assigning
zero when the range has zero width. Multiply by \(2^{17}-1\) and take the integer floor. Interleave
the 17 binary digits of the first, second, and third coordinates, in that order, to form a 51-bit
key. Sort the points by this key. Break a tie by writing each of the three coordinates as a decimal
string carrying 17 significant digits, joining the three strings in coordinate order with a vertical
bar between consecutive strings, and comparing the joined strings one character at a time by
character code. Let
\(n_M\) be the positive integer number of sorted feasible points, set
\(m_M=\min(8000,n_M)\), and retain the ranks
\begin{equation*}
 1+\left\lfloor j\frac{n_M}{m_M}\right\rfloor,
 \qquad j=0,\ldots,m_M-1.
\end{equation*}
Here \(m_M\) is the retained-point count. Traverse the
retained points in this order.

Let \(v_P\) be a primary endpoint and \(v_A\) its audit counterpart. When both have the bounded
endpoint state, define their nonnegative relative disagreement by
\begin{equation*}
 \Delta(v_P,v_A)=\frac{|v_P-v_A|}{\max\{1,|v_P|,|v_A|\}}.
\end{equation*}
Assign the unreliable endpoint state when \(\Delta(v_P,v_A)>10^{-4}\). Assign it after an audit
failure when the primary endpoint was bounded, and after any disagreement between the primary and
audit endpoint states, whatever those two states are. Always retain the
more extreme certified bounded value: the smaller lower endpoint or larger upper endpoint.
Here coverage describes the breadth of the numerical search; confidence-interval coverage is a
separate concept. Each accepted bounded endpoint must also pass \emph{cold-start endpoint replication}.
This replication runs at the end of the primary search and again at the end of the audit search, and
it is skipped for a side that already has the unbounded or the unreliable endpoint state.
Reevaluate its attaining structural SDF-news coefficient vector from a cold starting value that discards every
prior fit and warm starting value. If \(v_C\) is the cold-start coefficient and \(v_P\) is the accepted
coefficient, assign the unreliable endpoint state when the cold-start fit is not accepted, when
\(v_C\) is not finite, or when
\begin{equation*}
 |v_C-v_P|>10^{-6}\max\{1,|v_P|\}.
\end{equation*}

For either accepted fit with a fixed squared structural mean residual vector---the mean least-squares benchmark
reference fit or the zero-slack point fit---define estimating-equation contribution rows
\(s^P_t=X_t^V\{q_t-v_t\}\in\mathbb R^5\). The reported \emph{analytic covariance}, meaning
a covariance calculated from a formula rather than from resampling, is
\begin{equation*}
 \widehat V_P=A_P^{-1}\widehat M_4(s^P)A_P^{-1}.
\end{equation*}
Here \(\widehat V_P\in\mathbb R^{5\times5}\) is the PPML coefficient covariance matrix, and
\(A_P\) is the PPML information matrix. The matrix \(\widehat M_4(s^P)\) is the four-lag Bartlett
long-run covariance sum of the estimating-equation contribution rows, without a finite-sample
adjustment.
\subsection{Harvey normal quasi-maximum likelihood (Harvey QML)}

The Harvey QML estimator treats the conditional
variance as exponential, but uses the normal variance likelihood only as an estimation objective. For fixed \(\theta\),
choose \(\widehat\eta_H(\theta)\in\mathbb R^5\) to minimize
\begin{equation*}
 Q_H(\eta;\theta)=\frac12\sum_{t=1}^{T_v}
 \left[X_t^{V\prime}\eta+q_t(\theta)e^{-X_t^{V\prime}\eta}\right].
\end{equation*}
Here \(Q_H(\eta;\theta)\) is the scalar Harvey QML objective at conditional-scale coefficient vector
\(\eta\) and structural SDF-news coefficient vector \(\theta\).

Define the nonnegative scalar Harvey variance ratio
\(r_t^H(\theta,\eta)=q_t(\theta)e^{-X_t^{V\prime}\eta}\ge0\). The gradient, meaning the vector of first partial derivatives of the objective,
and the observed information, meaning its local curvature matrix, are
\begin{equation*}
 \nabla Q_H(\eta)=\frac12\sum_{t=1}^{T_v}X_t^V(1-r_t^H),
 \qquad
 A_H(\eta)=\frac12\sum_{t=1}^{T_v}r_t^HX_t^VX_t^{V\prime}.
\end{equation*}
Here \(\nabla Q_H(\eta)\in\mathbb R^5\) is the gradient vector, and
\(r_t^H\ge0\) is the Harvey variance ratio. The symmetric matrix
\(A_H(\eta)\in\mathbb R^{5\times5}\) is the Harvey QML observed-information matrix.
An exact zero squared structural mean residual contributes \(X_t^{V\prime}\eta/2\) to the objective
and zero to \(r_t^H\). Leave it unperturbed.

A \emph{recession direction} can make the objective decrease without bound. For design coordinate
\(j\in\{1,\ldots,5\}\), define
\begin{equation*}
 s_j^H=\max\left\{1,\left(\sum_{t=1}^{T_v}(X_{tj}^V)^2\right)^{1/2}\right\},
 \qquad Z_{tj}^H=\frac{X_{tj}^V}{s_j^H}.
\end{equation*}
Here \(s_j^H>0\) normalizes design column \(j\), and \(Z^H\in\mathbb R^{T_v\times5}\) is the
normalized conditional-scale design matrix. Let \(Z_+^H\) retain its rows with
\(q_t(\theta)>0\) and define \(c^H=\sum_t Z_t^H\in\mathbb R^5\), the sum running over all \(T_v\)
rows. A recession direction \(u\in\mathbb R^5\) satisfies \(Z_+^Hu\ge0\). Define the positive
recession-rate tolerance
\begin{equation*}
 \rho_H=10^{-9}\max\left\{1,\sum_{j=1}^5|c_j^H|\right\}.
\end{equation*}
The scalar \(\rho_H>0\) depends only on \(X^V\), so it takes the same value at every display slack.

The rule that decides whether a vector of squared structural mean residuals admits a recession
direction is the \emph{recession classification}. It returns exactly one of four verdicts: passing,
a certified negative recession direction, a zero-rate direction, and a failure to classify.

Before constructing any lattice at a display slack, subject the recession classification to a
three-part self-test on the frozen design \(X^V\). First, apply it to the vector whose \(T_v\)
entries all equal one; the verdict must be passing. Second, apply it to the vector whose \(T_v\)
entries all equal zero; the verdict must be a certified negative recession direction. Third, take
the trial value \(5\times10^{-9}\) together with the trial stationarity residual
\((10^{-3},10^{-3},10^{-3},10^{-3},10^{-3})'\in\mathbb R^5\); after subtracting the sum of that
residual's absolute coordinates, the corrected value \(5\times10^{-9}-5\times10^{-3}\) must not
exceed \(\rho_H\). If any of the three parts fails, perform no lattice construction and no search at
that slack: report every Harvey QML conditional-scale endpoint value as missing and assign every one
of them the unreliable endpoint state. The self-test reads no data beyond \(X^V\), so its verdict is
the same at every display slack, and because \(X^V\) has full column rank five all three parts hold.

Perform a recession-direction check before minimization exactly when at least one squared structural
mean residual is exactly zero; when every squared structural mean residual is strictly positive, the
objective cannot recede and the minimization proceeds directly. If every squared structural mean
residual is zero, the Harvey QML estimate does not exist and the fit fails.

The numerical rank of any matrix is obtained the same way, by dividing each nonzero column by its
Euclidean norm and counting the singular values of the result that exceed \(10^{-10}\) times the largest singular
value. Two branches of the recession classification are immediate. If every squared structural mean
residual is strictly positive and the design \(X^V\) has numerical rank five, the verdict is
passing. If no squared structural mean residual is strictly positive, the verdict is a certified
negative recession direction, witnessed by the direction with coordinates
\(-\operatorname{sign}(c_j^H)\), whose rate is \(-\sum_{j=1}^5|c_j^H|\).

In every remaining case at least one squared structural mean residual is exactly zero and at least
one is strictly positive, and the verdict comes from a facet certification. Examine all ten facets
indexed by \(j\in\{1,\ldots,5\}\) and \(\sigma\in\{-1,1\}\), where a \emph{facet} is the portion of
the direction box on which coordinate \(j\) equals \(\sigma\). Every nonzero recession direction can
be scaled onto one of these facets, so the direction coordinates are bounded by one. On each facet
minimize \((c^H)'u\) over \(Z_+^Hu\ge0\), \(u_j=\sigma\), and \(-1\le u_\ell\le1\), and classify the
facet from the result: a feasible value below \(-\rho_H\) is a negative recession direction; a
feasible value at most \(\rho_H\) is a zero-rate direction; a feasible value above \(\rho_H\) that a
dual lower bound---a number no larger than the minimum on that facet---confirms is a certified positive
facet; an infeasible facet that a phase-I bound---the matching bound for the smallest constraint
violation---confirms is a certified infeasible facet; anything else is unresolved. The Harvey recession-classification certificate subsection of Part~II states those bounds, their
starting values, and their tolerances. The
verdict is a certified negative recession direction as soon as one facet is a negative recession
direction; otherwise it is a zero-rate direction when at least one facet is a zero-rate direction;
otherwise it is passing when every facet is certified positive or certified infeasible; otherwise it
is a failure to classify.

The verdict decides the fit. A passing verdict lets the minimization proceed. A certified negative
recession direction means no finite Harvey QML estimate exists, and the fit at that structural
SDF-news coefficient vector fails. A zero-rate direction leaves finite attainment unresolved, and
the fit fails. A failure to classify also fails the fit. The Harvey QML domain analysis accounts for
no fit failure, so a fit that fails for any of these three reasons makes every Harvey QML
conditional-scale endpoint at that slack unreliable.

For a vector of squared structural mean residuals that passes the recession check, use a fixed
starting-value sequence in three stages, stopping at the first stage that yields an accepted fit.
The first stage uses the warm Harvey QML starting value: the one the search supplies when it
supplies one, and otherwise the accepted zero-slack Harvey QML fit's coefficients. The second stage
fits PPML at the current structural SDF-news coefficient vector and, when that PPML fit is accepted,
starts from its coefficients. Neither of the first two stages uses any further starting value. The
third stage runs the ladder below in order, omitting an unavailable starting value and stopping at
the first accepted fit. First, the accepted zero-slack PPML coefficients. Next, define the positive
normal log-square adjustment
\begin{equation*}
 \kappa_H=-\{\psi(1/2)+\log2\},
\end{equation*}
where \(\psi\) is the first derivative of the logarithm of the gamma function. If
\(\widehat\eta_{L,\mathrm{ref}}\in\mathbb R^5\) is the log-OLS coefficient vector from the
mean least-squares benchmark reference column, use
\begin{equation*}
 \widehat\eta_{L,\mathrm{ref}}+(\kappa_H,0,0,0,0)'.
\end{equation*}
Finally, use the intercept-only vector
\((\log\bar q,0,0,0,0)'\), where \(\bar q\) is the mean squared structural mean residual for the
current fit; this last rung belongs to the third stage alone and exists only when \(\bar q>0\).

For coordinate \(j\), normalize the first-order-condition residual
\(\sum_tX_{tj}^V(r_t^H-1)\) by \(\sum_t|X_{tj}^V|\). Define
\begin{equation*}
 S_H(\eta)=\max_{1\le j\le5}
 \frac{\left|\sum_tX_{tj}^V(r_t^H-1)\right|}{\sum_t|X_{tj}^V|}.
\end{equation*}
Here \(S_H(\eta)\ge0\) is the first-order-condition norm. Accept a starting value without an
update when \(S_H(\eta)\le10^{-8}\).

At each later iteration, use an observed-Newton direction when the diagonally normalized observed
information has reciprocal condition number at least \(10^{-12}\). The observed-Newton direction
uses the objective's actual local curvature. Otherwise, use a Fisher-scoring direction, which
replaces that curvature by its expectation under the normal conditional-variance model,
\(X^{V\prime}X^V/2\).

For either direction, try step lengths \(2^{-\iota}\) for \(\iota=0,\ldots,30\). Accept a trial whose
objective is strictly smaller. To define a numerically tied objective, let
\(\xi_t=X_t^{V\prime}\eta\) be the scalar linear predictor and let
\(\epsilon_{\mathrm{mach}}\) be the binary64 machine epsilon defined above. Treat the trial and
current objectives as tied when
\begin{equation*}
 |Q_{\mathrm{trial}}-Q_{\mathrm{current}}|
 \le4\epsilon_{\mathrm{mach}}
 \left(1+\sum_t|\xi_t|+\sum_{t:q_t>0}r_t^H\right).
\end{equation*}
Accept such a tie only when
\begin{equation*}
 S_{H,\mathrm{current}}-S_{H,\mathrm{trial}}
 >10\epsilon_{\mathrm{mach}}\max\{1,S_{H,\mathrm{current}}\}.
\end{equation*}
If the observed-Newton direction yields no accepted step, repeat the same line search with the
Fisher-scoring direction. If neither direction yields one, that starting value has stalled: continue
from the next starting value in the sequence. The fit fails only when every starting value has
stalled or otherwise failed.

After an update, require \(S_H(\eta)\le10^{-8}\) and at least one of
\begin{align*}
 |Q_H^{(m)}-Q_H^{(m-1)}|&\le10^{-10}\max\{1,|Q_H^{(m)}|\},\\
 \max_j|\eta_j^{(m)}-\eta_j^{(m-1)}|&\le
 10^{-10}\max\{1,\max_j|\eta_j^{(m)}|\}.
\end{align*}
Here superscript \(m\) indexes the iteration. Allow at most 1,000 iterations. Reaching the cap
without meeting the convergence conditions fails that starting value; continue from the next
starting value in the sequence.

The final fitted conditional variances must be finite and positive. Define the information-column
scale \(d^H_j=\{A_H(\eta)_{jj}\}^{1/2}\) and the diagonal matrix
\(D_H=\operatorname{diag}(d^H_1,\ldots,d^H_5)\in\mathbb R^{5\times5}\). Require the reciprocal
condition number of \(D_H^{-1}A_H(\eta)D_H^{-1}\) to be at least \(10^{-10}\).

The same implicit differentiation applied to the first-order condition gives
\begin{equation*}
 \frac{\partial\widehat\eta_H(\theta)}{\partial\theta'}
 =-A_H(\widehat\eta_H)^{-1}
   \sum_{t=1}^{T_v}X_t^V
   \left\{e_t(\theta)e^{-X_t^{V\prime}\widehat\eta_H}W_{2t}'\right\}.
\end{equation*}
Here the displayed Jacobian is a \(5\times3\) matrix measuring the response of the Harvey QML
conditional-scale coefficient vector to the structural SDF-news coefficient vector.

The primary Harvey QML conditional-scale estimator mapping uses a 4,000-lattice-point cap and a
fresh-fit allowance of 20,000 at each slack. Then run a separate \emph{Harvey sensitivity audit}
with up to 4,000 lattice points and a new fresh-fit allowance of 40,000 at each slack. For each
endpoint, use five separated lattice-derived starting values plus the recorded search start point, following the
PPML separation rule. Use one new fit record
across display slacks. Apply the PPML disagreement measure \(\Delta(v_P,v_A)\) and threshold
\(10^{-4}\). An audit failure, endpoint-state mismatch, or excessive disagreement makes a bounded
primary endpoint unreliable. Retain any more extreme certified bounded audit endpoint. Apply the
same cold-start endpoint replication inequality as for PPML, with threshold \(10^{-6}\).

For either accepted fit with a fixed squared structural mean residual vector---the mean least-squares benchmark
reference fit or the zero-slack point fit---define gradient contributions
\(s_t^H=\tfrac12(1-r_t^H)X_t^V\in\mathbb R^5\). The reported analytic covariance is
\begin{equation*}
 \widehat V_H=A_H^{-1}\widehat M_4(s^H)A_H^{-1}.
\end{equation*}
Here \(\widehat V_H\in\mathbb R^{5\times5}\) is the Harvey QML coefficient covariance matrix,
\(A_H\) is the Harvey QML observed-information matrix, and \(\widehat M_4(s^H)\) is the four-lag
Bartlett long-run covariance sum of the gradient contributions.
In each of the two sandwich covariances above, call the outer matrix \(A_P\) or \(A_H\) the
\emph{outer factor}: it is the symmetric five-by-five matrix that appears inverted on both sides of
\(\widehat V=A^{-1}\widehat M_4(s)A^{-1}\). For the PPML and Harvey QML analytic covariances, invert the outer factor through its normalized form. Let
\(D_A\) be the diagonal matrix whose entries are the square roots of the outer factor's diagonal, form
\(\widetilde A=D_A^{-1}AD_A^{-1}\), and recover \(A^{-1}=D_A^{-1}\widetilde A^{-1}D_A^{-1}\). The
acceptance conditions below apply to \(\widetilde A\); the covariance uses \(A^{-1}\). Require the conditional-scale sample's
observation count to exceed five; every design entry, coefficient, squared structural mean residual, and fitted
conditional variance to be finite; every squared structural mean residual to be nonnegative; every
fitted conditional variance to be strictly positive; every entry of the outer factor to be finite; every
diagonal entry of the outer factor to be strictly positive; the normalized outer factor's reciprocal condition
number to be at least \(10^{-10}\); and its Cholesky factorization, the representation of a symmetric
positive definite matrix as a lower triangular matrix times its own transpose, to succeed. If any of
these
fails, report no covariance matrix and no standard error for that column. Otherwise a coefficient's
standard error is the square root
of its diagonal covariance entry. Divide the coefficient by its standard error and report the
two-sided standard-normal p-value
\(2\{1-\Phi(|\text{coefficient}/\text{standard error}|)\}\).

\section{Conditional-scale checks and residual-dynamics closure}

\subsection{Joint-null compatibility check}

The hypothesis that the supplied return scores do not shift the conditional scale of the structural
mean residual requires the last four entries of a conditional-scale coefficient to be zero. For the log-OLS conditional-scale estimator mapping, let
the \emph{conditional-scale slope vector} \(\eta_R(\theta)\in\mathbb R^4\) collect those four non-intercept coefficients and let
\(D_R=\operatorname{diag}(1/\operatorname{sd}_{T_v-1}(r_1),\ldots,
1/\operatorname{sd}_{T_v-1}(r_4))\in\mathbb R^{4\times4}\). Measure their standardized squared
distance from zero by
\begin{equation*}
 D_0(\theta)=\frac12\|D_R^{-1}\eta_R(\theta)\|_2^2.
\end{equation*}
Here \(D_0(\theta)\) is the nonnegative scalar standardized squared distance from the joint null.
The matrix \(D_R\in\mathbb R^{4\times4}\) is the diagonal scale matrix for the supplied return
scores, and \(\eta_R(\theta)\in\mathbb R^4\) contains the four non-intercept log-OLS coefficients. Minimize this
distance over \(\Theta(\tau)\). Measure the same standardized departure in the maximum norm by
\begin{equation*}
 \varsigma(\theta)=\left\|D_R^{-1}\eta_R(\theta)\right\|_\infty
 =\max_{j\in\{1,\ldots,4\}}\operatorname{sd}_{T_v-1}(r_j)\,\left|\eta_{R,j}(\theta)\right|.
\end{equation*}
Here \(\varsigma(\theta)\ge0\) is the largest effect on a log-OLS conditional-scale coefficient of a
one-standard-deviation move in a supplied return score, and \(r_j\in\mathbb R^{T_v}\) is the
conditional-scale sample of supplied return score \(j\). Set the \emph{root tolerance} to
\(10^{-6}\). The objective \(D_0\) orders candidates; \(\varsigma\) decides them.

At zero slack the mean identified set is the single accepted zero-slack structural SDF-news
coefficient vector, so evaluate the mapping once at that vector and carry out no lattice, no
refinement, no reproduction, no perturbation, and no stability assessment, recording the
reproduction, perturbation, and stability entries as not applicable.

At every slack, this one included, record whether the slack is \emph{box-limited}, meaning that the
candidate the slack selects lies on a limit of the coefficient box over which it was minimized. Write
\(\underline\theta_k\in\mathbb R\) and \(\overline\theta_k\in\mathbb R\) for the scalar lower and upper limits of coordinate
\(k\) in the structural SDF-news coefficient box of the mean identified set at that slack, and write
\(\widehat\theta\) for the selected candidate, which at zero slack is the accepted zero-slack vector
and at a positive slack is the candidate selected by the rules below. The slack is box-limited when
\(|\widehat\theta_k-\underline\theta_k|\le10^{-6}\max\{|\overline\theta_k-\underline\theta_k|,1\}\) or
\(|\widehat\theta_k-\overline\theta_k|\le10^{-6}\max\{|\overline\theta_k-\underline\theta_k|,1\}\)
holds for at least one coordinate \(k\).
The flag is a statement about where the candidate landed, so record it even at a slack whose joint-null outcome
is unreliable for another reason. At zero slack the box is the single accepted vector itself, so its
two limits coincide with the candidate; an explicit exception records zero slack as not box-limited.

Let \(\operatorname{med}_{t}|e_{t}^{\mathrm{ref}}|\) be the
median absolute aligned mean least-squares benchmark residual defined below. The zero-slack result is
unreliable, with compatibility not demonstrated, when some structural mean residual is exactly zero,
when \(\min_t|e_t|\le10^{-6}\operatorname{med}_{t'}|e_{t'}^{\mathrm{ref}}|\), when \(D_0\) is
nonfinite, or when \(\rho_Q\) is nonfinite or exceeds \(10^{-4}\). Otherwise it is bounded, and it is
an \emph{attained compatibility witness}---a candidate certified as demonstrating compatibility,
whose positive-slack form the eight conditions below give---exactly when
\(\varsigma\le10^{-6}\). Every remaining rule in this
subsection governs the three positive slacks.

Begin with a 41-point lattice along each structural SDF-news coefficient, with the first coefficient
changing fastest, and retain points satisfying \(\rho_Q(\theta)\le10^{-10}\). If fewer than 100
points survive, replace this lattice once with 81 points along each coefficient. If fewer than 100
points survive the replacement as well, the slack is unreliable with compatibility not demonstrated
and none of the rules below is carried out. Pool the finite
lattice candidates with every minimum carried from a preceding slack. Each preceding slack contributes
its selected candidate together with every refined candidate accepted there, and each carried vector
carries the objective recomputed at that vector. Stop the reproduction if the smallest objective
attained at the current slack exceeds the smallest carried objective. Sort the pool first by \(D_0(\theta)\)
and then by the coordinate-by-coordinate tie rule of Section~\ref{sec:logvar-scale}, and select at
most five starting values. To determine
separation, use the same coordinate limits \(\underline\theta_k\) and \(\overline\theta_k\). Transform coordinate \(k\) to
\((\theta_k-\underline\theta_k)/(\overline\theta_k-\underline\theta_k)\); use width one when
\(\overline\theta_k-\underline\theta_k\) is nonpositive or nonfinite.
Accept a candidate whose transformed Euclidean distance from every previously selected starting value is at
least \(0.05\sqrt3\).

Treat each feasible starting value as a candidate and refine it by sequential least-squares
quadratic programming. Retain only finite candidates with \(\rho_Q(\theta)\le10^{-4}\). Among tied
objective values, select the first vector under the same coordinate-by-coordinate tie rule. Replace
the best pooled candidate, meaning the best of the lattice and carried vectors, with the best
refined candidate only when the refined objective is strictly smaller.

The vector retained by that comparison is the selected candidate \(\widehat\theta\) of this slack.
Measure distances in the perturbation rules below with the search-coordinate scale \(d_\theta>0\) of
Section~\ref{sec:idset}, formed from the same quadratic inequalities at the same slack. Only a refined
value that is itself finite, satisfies \(\rho_Q(\theta)\le10^{-4}\), and has a finite objective is
eligible to reproduce \(\widehat\theta\). Such a refined value reproduces \(\widehat\theta\) when every coordinate gap
is at most \(10^{-6}\max(d_\theta,|\widehat\theta_k|)\) and its objective gap is at most
\begin{equation*}
 \max\{64\epsilon_{\mathrm{mach}}\max(1,|D_0|,|\widehat D_0|),
       4(10^{-6})^2\}.
\end{equation*}
Here \(\epsilon_{\mathrm{mach}}\) is the binary64 machine epsilon defined above, \(D_0\) is the replicated
objective, and \(\widehat D_0\) is the selected objective. Require at least two reproducing starting values.
Distinct feasible local minima reached from separate starting values do not by themselves invalidate the check.

Let \(e_{t}^{\mathrm{ref}}\) be the scalar quarter-\(t\) residual from the mean least-squares
benchmark, aligned to the conditional-scale sample. Call quarter \(t\) active when
\(|e_t(\widehat\theta)|\le10^{-6}\operatorname{med}_{t'}|e_{t'}^{\mathrm{ref}}|\). Construct perturbation
directions from both signs of the three coordinate axes and, for every active quarter, both signs of
\(W_{2t}/\|W_{2t}\|_2\). To canonicalize an unsigned direction, find its first coordinate whose
absolute value exceeds \(10^{-14}\) and multiply the vector by minus one when that coordinate is
negative. Treat two canonical directions less than \(10^{-12}\) apart as duplicates. Call a
direction \emph{matched} when all three perturbations at radii
\(0.01d_\theta\), \(0.03d_\theta\), and \(0.10d_\theta\) are feasible and have finite conditional-scale estimator mappings.
The matched directions must span
\(\mathbb R^3\), where a singular value counts toward rank only when it exceeds \(10^{-10}\) times
the largest singular value. For every active quarter, at least one sign of that quarter's normal direction
must be matched.

Two rejection rules close the perturbation protocol, both stated in unstandardized
slope units. Reject the candidate when some matched direction has both \(\|\eta_R\|_2\) and
\(\min_t|e_t|\) at radius \(0.01d_\theta\) below \(0.15\) times their values at radius \(0.10d_\theta\);
a simultaneous collapse of the slope length and of the distance to the nearest
structural-mean-residual hyperplane between the two extreme radii is a cancellation signature rather
than an attained root. When at least one quarter is active, reject the candidate unless at least one
direction is feasible at radius \(0.01d_\theta\) and every direction feasible at that radius has a
finite \(\|\eta_R\|_\infty\) of at most \(10^{-4}\); a genuine isolated root stays near zero one small
step off every hyperplane. Either of two further conditions rejects the candidate before any
direction is constructed:
\(\min_t|e_t(\widehat\theta)|\le10^{-8}\operatorname{med}_{t'}|e_{t'}^{\mathrm{ref}}|\), which places
it closer to a hyperplane than the arithmetic can resolve, or an active quarter with
\(\|W_{2t}\|_2=0\), which leaves that quarter without a perturbation normal.

Carry out the reproduction and perturbation rules exactly when \(\varsigma(\widehat\theta)\le10^{-6}\)
or \(\min_t|e_t(\widehat\theta)|\le10^{-8}\operatorname{med}_{t'}|e_{t'}^{\mathrm{ref}}|\). When
neither holds, the slack is bounded with compatibility not demonstrated, and the reproduction,
perturbation, and stability entries record that they were not required, the stability entry being the
perturbation protocol's overall outcome. When one of them holds,
\(\widehat\theta\) is the attained compatibility witness if and only if all eight of the
following hold, and the slack is then bounded:
\begin{itemize}
\item every \(e_t(\widehat\theta)\) is finite;
\item every \(e_t(\widehat\theta)\) is nonzero;
\item \(\min_t|e_t(\widehat\theta)|>10^{-8}\operatorname{med}_{t'}|e_{t'}^{\mathrm{ref}}|\);
\item \(\varsigma(\widehat\theta)\le10^{-6}\);
\item \(\rho_Q(\widehat\theta)\) is finite and at most \(10^{-4}\);
\item recomputing the mapping at \(\widehat\theta\) returns the selected objective exactly, and every
 recomputed unstandardized slope agrees with the selected slope within
 \(10^{-10}\max\{1,\|\eta_R(\widehat\theta)\|_\infty\}\);
\item at least two refined starting values reproduce \(\widehat\theta\);
\item every perturbation rule above is satisfied.
\end{itemize}
If any of the eight fails, the slack is unreliable with compatibility not demonstrated. Record the
basin count as one plus the number of refined candidates that were feasible with a finite objective
and did not reproduce \(\widehat\theta\).

Record a maximum-norm sensitivity alongside the selected candidate. Pool the selected candidate, the
best lattice candidate, every separated starting value, the best refined candidate, and every refined candidate, keeping the finite ones; evaluate \(D_0\) and \(\varsigma\) at each. Let
\(\varsigma^\star\) be the smallest pooled value of \(\varsigma\), attained at the pooled candidate
\(\theta^\star\). Record \(\varsigma(\widehat\theta)\), \(\varsigma^\star\), the provenance of
\(\widehat\theta\) and of \(\theta^\star\), the rank of \(\widehat\theta\) among the pooled candidates
ordered by \(\varsigma\), the rank of \(\theta^\star\) among them ordered by \(D_0\), the absolute gap
\(\varsigma(\widehat\theta)-\varsigma^\star\), and the relative gap
\(\{\varsigma(\widehat\theta)-\varsigma^\star\}/\max(10^{-6},\varsigma^\star)\), ranking ties by the
smallest rank. The sensitivity triggers when \(\|\theta^\star-\widehat\theta\|_2>10^{-12}\), the
absolute gap exceeds \(10^{-5}\), and the relative gap exceeds \(0.10\). On a trigger, refine
\(\varsigma^\star\) by minimizing an auxiliary scalar \(\varpi\ge0\) subject to
\(\operatorname{sd}_{T_v-1}(r_j)|\eta_{R,j}(\theta)|\le\varpi\) for \(j=1,\ldots,4\) and
\(\theta\in\Theta(\tau)\), by derivative-free constrained minimization from both \(\widehat\theta\)
and \(\theta^\star\), in coordinates that divide \(\theta\) by \(d_\theta\) and \(\varpi\) by
\(\max(10^{-6},\varsigma^\star)\), with relative iterate tolerance \(10^{-10}\), at most 500
objective evaluations, and every normalized coordinate confined to \([-10^{6},10^{6}]\). Accept a
solution only when its recomputed point satisfies \(\rho_Q\le10^{-4}\), its recomputed \(\varpi\) is
at least \(-10^{-8}\), and its recomputed \(\varsigma\) is at most
\(\varpi+10^{-8}\max(1,\varpi)\); keep it only if it strictly lowers \(\varsigma\). Record the
sensitivity as refined when a solution is kept, as attained from the starting values alone otherwise,
and as not triggered whenever the slack is unreliable. The joint-null compatibility check remains a
diagnostic rather than a statistical hypothesis test or a certificate of the global minimum.

\subsection{Stacked-moment replication check}

The \emph{stacked-moment replication check} verifies that writing the log-OLS normal equations in
the larger ordered moment representation does not change the log-OLS conditional-scale estimator
mapping. Stacking means placing several sample-moment vectors one after another in one larger vector.
The current check evaluates only the log-OLS normal equations. It adds no identifying restriction and
produces no new estimate.

At a feasible structural SDF-news coefficient vector \(\theta\in\mathbb R^3\), define
\begin{equation*}
 g_L(\theta,\eta)=\frac{1}{T_v}X^{V\prime}
 \left[\ell(\theta)-X^V\eta\right].
\end{equation*}
Here \(g_L(\theta,\eta)\in\mathbb R^5\) is the log-OLS sample-moment vector,
\(\ell(\theta)\in\mathbb R^{T_v}\) is the log squared structural mean residual vector defined in
Section~\ref{sec:logols}, and \(\eta\in\mathbb R^5\) is a conditional-scale coefficient vector.
Evaluate this vector at \(\eta=\widehat\eta_L(\theta)\), the log-OLS conditional-scale estimator mapping.

Use the finite zero-slack point. Then, when the baseline mean identified set's coefficient box has a
finite lower and a finite upper limit in every coordinate and its upper limit is at least its lower
limit in every coordinate, add the first at most eight points of a three-by-three-by-three lattice
over that box, retaining a lattice point only when \(\rho_Q(\theta)\le10^{-10}\). When the box fails
either condition, the zero-slack point is the only replication point. If no replication point
survives, the reproduction stops. Preserve lattice order and
retain the zero-slack point again if it also occurs in the lattice. This box is that of the mean
identified set at the baseline slack. The joint-null compatibility check reads the same box at that
slack, and the corresponding box at each of the other two display slacks as well. At
each point, require finite residuals and a finite log-OLS coefficient vector. Record the check as
attained only when at least one point is finite and
\begin{equation*}
 \max_{\theta}\left\|g_L\{\theta,\widehat\eta_L(\theta)\}\right\|_\infty<10^{-8},
\end{equation*}
where the maximum spans the retained replication points. The vector inside the norm contains the
log-OLS residual mean in its intercept coordinate and the four sample covariances between the
centered lagged supplied return scores and the log-OLS residual. Also record the smallest absolute
structural mean residual across these points. Write the current three-row record as the graph-replication row, whose recorded numerical status is
the outcome of the displayed comparison---attained when it holds and unreliable when it does
not---followed by two fixed skipped rows. Write the disabled projection record with no
rows. These records document the fixed current decision and define no additional calculation. The
stacked-moment replication check is a diagnostic rather than a joint estimation procedure.

\subsection{Residual-dynamics check and current closure}

At an accepted zero-slack point, define the log-OLS residual
\begin{equation*}
 u_t^L=2\log|e_t(\widehat\theta_0)|-X_t^{V\prime}\widehat\eta_L(\widehat\theta_0).
\end{equation*}
Here \(u_t^L\) is the scalar zero-slack log-OLS residual in quarter \(t\), and
\(\widehat\theta_0\in\mathbb R^3\) is the accepted zero-slack structural SDF-news coefficient
vector. The \emph{residual-dynamics check} applies a fixed serial-correlation calculation to this
residual sequence. Sort the conditional-scale quarters into ascending calendar order first. Stop the
reproduction if the recomputed zero-slack log-OLS coefficient vector differs from the frozen
zero-slack log-OLS coefficient vector by more than \(10^{-10}\) in any coordinate; this identity
proves the residual sequence is the paper's own and not a re-estimate on a drifted sample. Record the
\emph{unreliable residual-dynamics outcome}---every Ljung--Box statistic, every p-value, every autocorrelation, and the residual table left empty---when any coordinate of the zero-slack point is
unavailable, when some structural mean residual is exactly zero, or when some log-OLS residual is
nonfinite, recording the offending quarters in the last two cases.

Otherwise let \(n_L=T_v\) be the
number of zero-slack log-OLS residuals, one per conditional-scale quarter, require \(n_L\ge5\), and
define their sample
autocorrelation at lag \(j\) by
\begin{equation*}
 \widehat\rho_j^L=
 \frac{\sum_{t=j+1}^{n_L}(u_t^L-\bar u^L)(u_{t-j}^L-\bar u^L)}
      {\sum_{t=1}^{n_L}(u_t^L-\bar u^L)^2}.
\end{equation*}
Here \(\bar u^L\) is the sample mean of the log-OLS residuals, and
\(\widehat\rho_j^L\) is their scalar lag-\(j\) sample autocorrelation. For
\(m\in\{1,4,8\}\), calculate the Ljung--Box statistic
\begin{equation*}
 \mathcal T_{LB,m}=n_L(n_L+2)\sum_{j=1}^{m}\frac{(\widehat\rho_j^L)^2}{n_L-j}.
\end{equation*}
Here \(\mathcal T_{LB,m}\ge0\) is the lag-\(m\) Ljung--Box statistic. Its p-value is the upper-tail
probability of a chi-squared random variable with \(m\) degrees of freedom. Do not reduce those
degrees of freedom for fitted parameters. The check holds the zero-slack point fixed and ignores its
estimation uncertainty.

Compare only the lag-4 p-value with 0.05. The fixed inputs give a lag-4
p-value of 0.330. The \emph{non-reject residual-dynamics outcome} means that the decision p-value is
not below 0.05. Record this outcome because \(0.330>0.05\),
close the residual-dynamics check, and leave the four dynamic-volatility output locations empty. A
\emph{dynamic-volatility output location} is a reserved location that the current route never fills.

\section{Least-absolute-deviation (LAD) median-regression estimator}

LAD estimates the conditional median of the log squared structural mean residual. It supplies a
conditional median absolute-residual scale rather than a conditional standard deviation. Its
conditional-scale estimator mapping is
\begin{equation*}
 \widehat\eta_A(\theta)
 =\arg\min_{\eta\in\mathbb R^5}
   \sum_{t=1}^{T_v}|2\log|e_t(\theta)|-X_t^{V\prime}\eta|.
\end{equation*}
Here \(\arg\min\) denotes a coefficient vector attaining the smallest displayed objective. The
Barrodale--Roberts selection below resolves multiple minimizing vectors. The fitted conditional median
absolute-residual scale is \(\exp(X_t^{V\prime}\eta/2)\).

Using the aligned mean least-squares benchmark residual \(e_t^{\mathrm{ref}}\) defined above, set
\(m_e=\operatorname{med}(|e_t^{\mathrm{ref}}|)\), the same median absolute aligned benchmark residual
the joint-null compatibility check uses, except for the fallback that follows. If this value is
nonpositive, use
the median of its strictly positive finite absolute residuals. If that value is also unavailable,
use their finite maximum. LAD is unavailable if no strictly positive finite value remains.

The LAD domain excludes every structural-mean-residual hyperplane
\(e_t(\theta)=W_{1t}-W_{2t}'\theta=0\). A candidate satisfying
\(0<|e_t(\theta)|\le10^{-12}m_e\) for any \(t\) belongs to the mathematical domain, but is
numerically unresolved. Reject that candidate without replacing or truncating its structural mean
residual. At an accepted candidate, set \(z_t^A(\theta)=2\log|e_t(\theta)|\) and minimize
\(\sum_t|z_t^A(\theta)-X_t^{V\prime}\eta|\) at the median quantile.

Select the minimizing coefficient vector with the Barrodale--Roberts vertex rule. A
\emph{vertex solution} is a candidate coefficient vector determined by five linearly independent
design rows whose fitted residuals equal zero. Those five rows form the
\emph{five-row basis}. This finite-step rule moves among vertex solutions by a
\emph{pivot}, which replaces one basis row with another whenever the replacement reduces
the sum of absolute residuals. It stops when no permitted pivot reduces that sum. Use the
\emph{LAD numerical reference}, release 6.1, to fix the deterministic pivot and tie ordering. A
\emph{pivot table} is the rectangular numerical array that records the current basis
transformation. At each pivot, scan eligible working columns in their current pivot-table order and
retain the first column whose entry is strictly largest under the pivot criterion that release 6.1 of
the LAD numerical reference fixes. Then scan eligible
working rows in their current pivot-table order and retain the first row that attains the strictly
smallest ratio. Every pivot swaps rows and columns, so later ties follow the updated pivot-table
order. Verify release 6.1, preserve the initial observation and design-column orders, give every
observation equal weight, and calculate no confidence interval from this fit. Set the Barrodale--Roberts numerical tolerance to
\(\epsilon_{\mathrm{mach}}^{2/3}=2^{-104/3}\). Stop if the design has fewer than five linearly
independent columns.

Use a Frisch--Newton interior-point fit only as a nonuniqueness check. This fit stores and operates
on every matrix entry. The Frisch--Newton method approaches the
minimum through the interior by solving successive smooth barrier approximations. A
smooth barrier approximation adds a differentiable penalty that grows near the boundary and thereby
keeps each intermediate coefficient vector in the interior. At every iteration, multiply the
largest admissible primal step and the largest admissible dual step by the step fraction 0.99995
and cap each resulting step at one. A primal step updates the two positive quantities that measure
the lower- and upper-bound residuals. A dual step updates their two positive multiplier quantities. An
admissible step retains strict positivity of all four quantities. Let
\(x_t^{FN},s_t^{FN}>0\), for \(t=1,\ldots,T_v\), denote the lower- and upper-bound residual quantities,
and let \(z_t^{FN},w_t^{FN}>0\) denote their corresponding multipliers. Define the nonnegative scalar complementarity gap
\begin{equation*}
 G_{FN}=\sum_{t=1}^{T_v}\left(z_t^{FN}x_t^{FN}+w_t^{FN}s_t^{FN}\right).
\end{equation*}
Stop when \(G_{FN}\le10^{-6}\), or when the iteration limit fixed by release 6.1 of the LAD numerical
reference is reached. If
\(Q_{BR}\) and \(Q_{FN}\) are the two
objective values, call them tied when
\begin{equation*}
 |Q_{BR}-Q_{FN}|\le10^{-8}\max(1,|Q_{BR}|).
\end{equation*}
Here \(Q_{BR}\) is the scalar Barrodale--Roberts objective, and \(Q_{FN}\) is the scalar
Frisch--Newton objective. When the objectives are tied, call the two coefficient selections materially
different when
\begin{equation*}
 \max_j|\widehat\eta_{A,j}^{BR}-\widehat\eta_{A,j}^{FN}|
 >10^{-4}\max\{1,\max_j|\widehat\eta_{A,j}^{BR}|\}.
\end{equation*}
The vectors \(\widehat\eta_A^{BR},\widehat\eta_A^{FN}\in\mathbb R^5\) are the two coefficient
selections. Always retain the Barrodale--Roberts selection.

The current LAD decision---the approved, declined, or unanswered value catalogued in Part~II that
determines whether LAD enters the reproduction---is approved and requires the LAD numerical
reference. The reproduction
therefore includes every LAD calculation in this section. Use the LAD numerical reference for every
LAD fit. It fixes the coefficient selection when several vectors minimize the same objective.
Use a fresh-fit allowance of 45,000 at each display slack. Within that total, set activity allowances
of 20,000 for lattice scanning, 16,000 for structural-mean-residual-hyperplane approaches, 2,000 for
nonuniqueness checks, 26,000 for endpoint refinement, and 4,000 for cold-start endpoint replication.
Both the total allowance and each activity allowance apply. Use one fit record across display slacks,
so a fit recorded at one slack is reused at any later slack, and start a new fresh-fit allowance at
each display slack.

Use a lattice-point cap of 20,000, both when building the baseline-slack lattice from which the
fallback search start point is taken and at every display slack. Use
the feasible zero-slack point as the search start point; if it is unavailable, use the first
lattice point retained after thinning. For each coefficient and endpoint side, rank lattice points
first by endpoint extremeness and then by the coordinate-by-coordinate tie rule of
Section~\ref{sec:logvar-scale}. Retain the extreme point and up to two further
points whose Euclidean distance from every retained point exceeds five percent of the structural
SDF-news coefficient box's Euclidean diagonal, giving three lattice-derived starting values per side.
The LAD independent endpoint audit repeats the same search with a new fit record, a new fresh-fit
allowance carrying the same activity allowances, and, per side, the extreme point plus up to four
further separated lattice-derived starting values. An endpoint differs
materially from its independent counterpart when their gap exceeds
\(10^{-4}\max(1,|v_P|)\), where \(v_P\) is the primary endpoint. If that audit finds a
strictly more extreme argument, repeat the primary search once from that argument and then compare
again. Assign the unreliable endpoint state to a bounded primary endpoint side in each of three
cases: that audit does not hold that side bounded; it holds it bounded at a value that is
neither strictly more extreme nor materially equal; or that audit does not complete, in
which case every bounded primary endpoint side is demoted. Demotion changes the endpoint state only:
keep the endpoint value and its attaining argument unchanged, and recompute the slack's state from
its two side states.

Apply the Frisch--Newton nonuniqueness check in a fixed order: for each coefficient in turn, its
bounded lower endpoint argument then its bounded upper endpoint argument; then the search start point; then
every argument promoted by that audit. Skip an argument whose fit is not accepted, without
consuming the allowance. Call an argument \emph{decisive} when it is an endpoint argument or a
promoted argument; the search start point is not decisive. Assign every endpoint at that slack the
unreliable endpoint state when the two coefficient selections are materially different at a decisive
argument, or when the 2,000 nonuniqueness-check activity allowance stops the order before a decisive
argument has been reached. A materially different selection at the search start point alone changes no
endpoint state. Keep every endpoint value and attaining argument unchanged. Each bounded endpoint side must also pass cold-start endpoint replication. Refit its attaining
argument once from a cold starting value that discards every recorded fit and warm starting value,
and assign that side the unreliable endpoint state when the refitted coefficient is nonfinite or when
\(|v_C-v_P|>10^{-9}\max\{1,|v_P|\}\), with \(v_C\) the cold-start coefficient and \(v_P\) the
accepted coefficient. Skip this replication on a side already carrying the unbounded or unreliable
endpoint state. LAD supplies identified ranges, a table, and figures of the conditional median
absolute-residual scale. Exclude LAD from all bootstrap calculations.

Before any search at a display slack, apply the hyperplane-intersection check of
Section~\ref{sec:logols} to every conditional-scale quarter. It returns the quarters whose
structural-mean-residual hyperplanes meet the mean identified set, called the \emph{certified
crossing quarters}, and the quarters it leaves uncertified. Apply every rule in the remainder of this
section to exactly the certified crossing quarters. If any certified crossing quarter fails to yield
a verified crossing by the rules below, or if the check leaves any quarter uncertified, assign every
endpoint at that slack the unreliable endpoint state and carry out no further work at that slack.
During the lattice scan, tolerate a candidate rejected for lying on or numerically adjacent to a
structural-mean-residual hyperplane only when every quarter implicated in the rejection is a certified
crossing quarter; any other rejected candidate assigns every endpoint at that slack the unreliable
endpoint state.

For a structural-mean-residual hyperplane in quarter \(t\), use the existing expression
\(e_t(\theta)=W_{1t}-W_{2t}'\theta\). Let \([\underline\theta_k,\overline\theta_k]\) be coordinate
\(k\)'s side of the coefficient box and \(\mu_k^B=(\underline\theta_k+\overline\theta_k)/2\). Here
\(\mu^B=(\mu_1^B,\mu_2^B,\mu_3^B)'\) is the box midpoint. Define the all-lower corner
\(\underline\theta^B=(\underline\theta_1,\underline\theta_2,\underline\theta_3)'\in\mathbb R^3\) and
the all-upper corner
\(\overline\theta^B=(\overline\theta_1,\overline\theta_2,\overline\theta_3)'\in\mathbb R^3\). Every distance in the crossing machinery below is measured in the search-coordinate scale \(d_\theta\). No per-coordinate scale enters.
If
\begin{equation*}
 \|W_{2t}\|_2\le10^{-12}\max\{1,\max_{t',k}|W_{2k,t'}|\},
\end{equation*}
where the inner maximum spans all conditional-scale quarters \(t'\) and coordinates
\(k\in\{1,2,3\}\),
then classify the hyperplane as unavailable when \(|W_{1t}|\le10^{-12}m_e\) and as having no crossing
otherwise. For every other quarter, construct nine base starting values: \(\mu^B\),
\(\underline\theta^B\), \(\overline\theta^B\),
and the six points formed from \(\mu^B\) by adding or subtracting \((\overline\theta_k-\underline\theta_k)/4\) in coordinate
\(k\in\{1,2,3\}\) while leaving the other coordinates unchanged. For each base starting value
\(a\in\mathbb R^3\), also construct its Euclidean projection
\begin{equation*}
 p_t(a)=a+W_{2t}\frac{W_{1t}-W_{2t}'a}{W_{2t}'W_{2t}}.
\end{equation*}
Here \(p_t(a)\in\mathbb R^3\) is the Euclidean projection of \(a\) onto the quarter-\(t\)
hyperplane.
Remove duplicate starting values. From each remaining value, minimize
\(\{e_t(\theta)/m_e\}^2\) over the mean identified set by sequential least-squares quadratic
programming. Use relative iterate tolerance \(10^{-8}\), permit at most 2,000 objective evaluations
per starting value, and retain a candidate only when \(\rho_Q(\theta)\le10^{-8}\). Among retained candidates,
select the one with the smallest \(|e_t(\theta)|/m_e\). Denote it by \(b\in\mathbb R^3\) and form its
Euclidean projection \(b_c=p_t(b)\) onto the quarter-\(t\) hyperplane, using the expression for
\(p_t(\cdot)\) displayed above.
Here \(b_c\in\mathbb R^3\) lies on the quarter-\(t\) hyperplane. Accept it only when
\(|e_t(b_c)|/m_e\le10^{-13}\) and \(\rho_Q(b_c)\le10^{-8}\); an accepted \(b_c\) is a
\emph{verified crossing}.

From an accepted crossing, search each residual-sign side over the candidate vectors
\(\theta\in\mathbb R^3\) in the ball
\begin{equation*}
 \|\theta-b_c\|_2\le0.05\,d_\theta.
\end{equation*}
On each side, maximize the signed residual by sequential least-squares quadratic programming subject
to this neighborhood and the mean identified set. Use relative iterate tolerance \(10^{-8}\) and
permit at most 1,000 objective evaluations. For a proposed \emph{LAD side anchor}
\(b^*\in\mathbb R^3\), set
\(b_a=b^*\) when \(\rho_Q(b^*)\le0\). When \(\rho_Q(b^*)>0\), bisect the segment fraction interval
\([0,1]\) 50 times. At each bisection, test \(b_c+\upsilon(b^*-b_c)\) at the midpoint fraction
\(\upsilon\); replace the lower fraction when \(\rho_Q\le0\) and the upper fraction otherwise. Set
\(b_a\in\mathbb R^3\) to the point at the final lower fraction; when no midpoint is ever feasible
this leaves \(b_a=b_c\), which the residual-magnitude condition below then rejects. Retain \(b_a\) as the LAD side anchor only when
\begin{equation*}
 \frac{\|b_a-b_c\|_2^2}{d_\theta^2}\le0.05^2+10^{-8},
\end{equation*}
its residual has the intended sign, and \(|e_t(b_a)|\ge10^{-8}m_e\). At least one side must
survive. From each retained LAD side anchor \(b_a\), evaluate
\begin{equation*}
 b_s=b_c+2^{-(s-1)}(b_a-b_c),\qquad s=1,\ldots,40.
\end{equation*}
Here \(b_s\in\mathbb R^3\) is approach point \(s\). Stop the approach at the first \(s\) that is
infeasible, that satisfies \(|e_t(b_s)|\le10^{-12}m_e\), or whose LAD fit is not accepted or has a
nonfinite coefficient, and record no point at or beyond that \(s\). For every earlier point set
\(M_s=-2\log\{|e_t(b_s)|/m_e\}\).
Here \(M_s\in\mathbb R\) is the scalar transformed log distance from the quarter-\(t\)
structural-mean-residual hyperplane. Record the points in approach order with no further filtering;
finiteness and strict increase of \(M_s\) are required window by window in the eligibility rule
below.

A \emph{coefficient tail} is one LAD coefficient evaluated along the strictly increasing sequence
of finite \(M_s\) values. Associate every fit with a \emph{tail label}, a categorical label used only
to determine whether consecutive fits belong to one tail. All accepted fits in this reproduction
receive the same neutral tail label. An eligible window has strictly increasing finite \(M_s\), a
finite coefficient at every point, an accepted fit at every point, and one unchanged tail label.

For coefficient \(j\in\{1,\ldots,5\}\) at approach point \(s\), let
\(\widehat\eta_{A,j,s}\in\mathbb R\) be the scalar LAD coefficient. For the last six eligible
points, regress this coefficient on \(M_s\) and call the fitted slope \(d_6\). Define
\begin{equation*}
 h_6=10^{-6}\max\{1,\max_{s\ \mathrm{in\ the\ six\ points}}|\widehat\eta_{A,j,s}|\}.
\end{equation*}
The tail is stable and finite when every one of its last four increments has absolute value at most
\(h_6\) and
\begin{equation*}
 |d_6|(\max M_s-\min M_s)<h_6
\end{equation*}
over those six points.

Persistent divergence requires the last eight eligible points. Let \(d_8\) be the corresponding
eight-point slope and define
\begin{equation*}
 h_8=10^{-6}\max\{1,\max_{s\ \mathrm{in\ the\ eight\ points}}|\widehat\eta_{A,j,s}|\}.
\end{equation*}
Both slopes must be nonzero and have the same sign, and they must satisfy
\begin{equation*}
 \frac{|d_6-d_8|}{\max(|d_6|,|d_8|,10^{-12})}\le0.20.
\end{equation*}
Every increment among the last six points must have exactly the sign of \(d_6\). The eight-point
\(M_s\)-span must be at least 12, and the eight-point linear regression must have \(R^2\ge0.98\).
The fitted six-point movement must exceed both \(10h_8\) and
\(0.05\max\{1,\operatorname{med}_s|\widehat\eta_{A,j,s}|\}\), where the median uses the eight
points. Classify each of the five coefficients, on each approach path of each verified crossing, into
exactly one of four states, tested in this order. The tail is \emph{uninformative} when the path
recorded fewer than six approach points, including the case of no recorded point at all. Otherwise it
is \emph{stable and finite} when the six-point rule above holds. Otherwise it shows \emph{persistent
divergence} when the eight-point rule above holds. Otherwise it is \emph{unresolved}.

The four states have different consequences. Persistent divergence with a positive \(d_6\) gives that
coefficient's upper endpoint the unbounded endpoint state, and with a negative \(d_6\) gives its lower
endpoint that state. A stable and finite tail assigns no endpoint state: record it separately as a
one-sided crossing-limit approximation carrying the coefficient, the slack, the crossing quarter, the
residual side, the approach-path index, the span of \(M_s\), the fitted six-point slope, the fitted six-point movement, and the last recorded coefficient value, and never merge it into an identified
range. An uninformative tail contributes no evidence and is passed over. An unresolved tail assigns
the unreliable endpoint state to both sides of that coefficient. A path that cannot be classified at
all assigns the unreliable endpoint state to both sides of every coefficient.

Carry out this classification on all five coefficients on every approach path, including a path that
recorded no approach point, on which all five are uninformative. After every approach path of every
verified crossing at a slack has been classified, and provided at least one crossing was verified
there, assign the unreliable endpoint state to both sides of every coefficient that received no
stable, divergent or unresolved classification on any path. Keep every finite attained
numerical value separate from this one-sided classification. Here \(R^2\) is the fraction of the
tail coefficient's variation explained by its fitted linear relation with \(M_s\).

\section{Circular moving-block bootstrap}
\label{sec:bootstrap}

Use \(B=10{,}000\) bootstrap draws. Here \(B\) is the positive integer number of bootstrap draws.
Confirm this value before the reproduction begins: the run overwrites the paper's typeset numerical
results, and it does not proceed at any other draw count without the explicit draft
acknowledgement described in Part~II.

A \emph{circular moving-block bootstrap} repeatedly samples
consecutive runs of quarters and wraps the sample end back to its beginning, thereby preserving
short-run dependence. A \emph{bootstrap draw} is one sampled index sequence and every estimate
calculated from it. The \emph{original sample} is the corresponding observations before bootstrap sampling.

The present reproduction starts without a reusable bootstrap result. Confirm that no reusable
bootstrap result exists, perform every bootstrap draw, and preserve the new result only after full
validation.

Distribute the bootstrap draws across \emph{parallel workers}, independent evaluators
that each calculate assigned bootstrap draws. Generate every sampled index sequence before parallel
evaluation, so neither the number of workers nor their scheduling can change any bootstrap draw.
The number of workers is therefore free and leaves every reported quantity unchanged.

Resampling acts on the identification sample, so every block-length rule below uses its observation
count \(N_O\), not the smaller conditional-scale count \(T_v\). Set the primary block length to
\begin{equation*}
 m_{\mathrm{pri}}=\left\lceil1.5N_O^{1/3}\right\rceil.
\end{equation*}
Here \(m_{\mathrm{pri}}\) is the positive integer primary block length, and \(N_O\) is the
identification-sample observation count defined in Section~\ref{sec:meaneq}. The rule gives
\(m_{\mathrm{pri}}=10\) for every \(N_O\) between 217 and 296, which covers the present sample.
Set the requested sensitivity block length to
\(m_{\mathrm{sens}}=2m_{\mathrm{pri}}=20\). Here \(m_{\mathrm{sens}}\) is a positive integer. It is
below \(N_O\), so its effective block length is also 20. The bootstrap draws
using \(m_{\mathrm{pri}}\) form the \emph{primary bootstrap-index family}; the bootstrap draws using
\(m_{\mathrm{sens}}\) form the \emph{sensitivity bootstrap-index family}. For either requested
block length \(m\), set \(m_{\mathrm{eff}}=\min(m,N_O)\). Here \(m_{\mathrm{eff}}\) is its positive
integer effective block length. Draw \(K_B=\lceil N_O/m_{\mathrm{eff}}\rceil\) starting positions independently
and uniformly with replacement, allowing the same position to appear more than once. From each starting position,
take \(m_{\mathrm{eff}}\) consecutive positions, wrapping from \(N_O\) to 1. Concatenate the blocks and retain the
first \(N_O\) indices.
Here \(K_B\) is the positive integer number of sampled blocks.
The sensitivity bootstrap-index family changes the block length to measure how strongly the
bootstrap conclusions depend on that choice. Each family contains \(B\) bootstrap draws, so
\(20{,}000\) sampled index sequences are generated in all. The two families do not carry the same
content. Every bootstrap draw of the primary family produces the mean-equation quantities and the
conditional-scale quantities together, from a single system estimate. Every bootstrap draw of the
sensitivity family produces conditional-scale quantities only.

Use the 32-bit MT19937 Mersenne--Twister generator for uniform random numbers. Set the
normal-random-number method to inversion and set discrete sampling to rejection. Initialize each
bootstrap-index family independently with seed 20260708. This construction draws no normal
random numbers; therefore, the inversion setting remains part of the random-number rules without changing
the sampled indices.

The following equations fully determine the random-number stream across languages. All state quantities are
unsigned 32-bit integers, all arithmetic wraps modulo \(2^{32}\), \(\oplus\) denotes bitwise
exclusive-or, \(\mathbin{\&}\) denotes bitwise and, \(\mathbin{|}\) denotes bitwise or, and
\(\gg,\ll\) denote right and left bit shifts. Define
\begin{equation*}
 \mathcal L(s)=(69069s+1)\bmod 2^{32}.
\end{equation*}
Starting from \(s=20260708\), apply \(\mathcal L\) 50 times. Discard the next update and fill the
624-word state \(\mathfrak m_0,\ldots,\mathfrak m_{623}\) with the following 624 successive updates.
Set the state position to 624.

Whenever the position reaches 624, replace every state word, using the old state on every
right-hand side, by
\begin{align*}
 y_i&=(\mathfrak m_i\mathbin{\&}2^{31})
      \mathbin{|}(\mathfrak m_{i+1}\mathbin{\&}(2^{31}-1)),\\
 \mathfrak m_i
 &\leftarrow\mathfrak m_{i+397}\oplus(y_i\gg1)
   \oplus
   \begin{cases}
    2567483615,&y_i\text{ is odd},\\
    0,&y_i\text{ is even}.
   \end{cases}
\end{align*}
Here all state subscripts are reduced modulo 624. Apply these two assignments in place and in
increasing order of \(i\), so that a right-hand side already replaced earlier in the same sweep
contributes its replaced value. For \(i\ge227\) the subscript \(i+397\) reduces to \(i-227\), which
this sweep has already replaced, and at \(i=623\) the subscript \(i+1\) reduces to 0, likewise
already replaced. Reading the unreplaced state instead would reproduce only the first 227 draws.
Reset the position to zero. For the next state word
\(w\), apply the four tempering transformations
\begin{align*}
 y&=w\oplus(w\gg11),\\
 y&=y\oplus\{(y\ll7)\mathbin{\&}2636928640\},\\
 y&=y\oplus\{(y\ll15)\mathbin{\&}4022730752\},\\
 y&=y\oplus(y\gg18).
\end{align*}
The resulting uniform draw is \(U=y/2^{32}\) when \(y>0\) and
\(U=\{2(2^{32}-1)\}^{-1}\) when \(y=0\). Advance the state position by one.

For a uniform choice among \(n\) positions, let
\(q_n^{\mathrm{bit}}=\lceil\log_2 n\rceil\). Form a \(q_n^{\mathrm{bit}}\)-bit
nonnegative integer from the low-order bits of successive 16-bit pieces
\(\lfloor2^{16}U\rfloor\), where \(U\in[0,1)\) is the next Mersenne--Twister uniform draw. Consume one
such piece for each of \(n'=0,16,32,\ldots\) with \(n'\le q_n^{\mathrm{bit}}\), accumulating
\(v\mapsto2^{16}v+\lfloor2^{16}U\rfloor\), and then take the low \(q_n^{\mathrm{bit}}\) bits of
\(v\). If the
integer is at least \(n\), discard it and repeat; otherwise select that integer plus one. This
rejection rule gives all \(n\) positions equal probability. Different block lengths consume
different numbers of selections. Therefore, bootstrap draw \(b\) in the two bootstrap-index families generally uses
different positions in its family's random-number stream. Preserve the complete random-number state
after constructing each bootstrap-index family, and restore the state that existed before construction.

For each bootstrap draw, sample the entire identification sample in the generated order and estimate
the shared mean system once on all resampled rows. Recompute all seven sample moments and the mean
identified set. Only afterward, retain for the conditional-scale calculation the rows on which no
lagged supplied return score is missing or undefined. An infinite score is neither missing nor
undefined and does not remove its row. Apply this one row mask to the re-estimated outcome
reduced-form residual, the three re-estimated SDF-news reduced-form residuals, the quarter label, and
the four lagged supplied return scores. Residual finiteness is not an additional row-selection condition. Fit PPML
first and Harvey QML second because the Harvey QML initialization depends on the PPML fit.

At zero slack, solve the linear system and evaluate each requested coefficient functional directly.
At each positive slack, construct a 41-by-41-by-41 coordinate lattice and retain feasible points in
lattice order, with the first structural SDF-news coefficient changing fastest. If fewer than 100
points are feasible, replace this lattice once with an 81-by-81-by-81 lattice. If \(m_G>1500\)
feasible lattice points remain, set \(h_G=\lceil m_G/1500\rceil\) and retain points numbered
\(1,1+h_G,1+2h_G,\ldots\). Here \(m_G\) is the number of feasible lattice points, and \(h_G\) is
the positive integer thinning interval. Append the feasible search start point afterward, even when it
duplicates a retained point. The scan therefore contains at most 1,500 retained lattice points plus
one appended search start point. Use a fresh-fit allowance of 8,000 across the scan and endpoint searches. Do not perform cold-start endpoint
replication. Recompute every lower and upper endpoint and record its endpoint state.

Keep two failure counts distinct. A \emph{bootstrap-draw calculation failure} means that the
mean-equation calculation or conditional-scale calculation within a bootstrap draw returned no
structured result. A structured
calculation result contains that calculation's required estimates and endpoint states in their declared
arrangement. A bootstrap draw with a bootstrap-draw
calculation failure is neither retried nor discarded. It keeps its position among the \(B\) bootstrap
draws, and every endpoint it would have produced carries the failed endpoint state and no numerical
value, on both sides, at every slack, and at the zero-slack point. Count such bootstrap draws
separately for the mean-equation calculation and for the conditional-scale calculation, and within
the primary bootstrap-index family only. Each of the two counts must be at most
\(\lfloor0.25B\rfloor\), which is 2,500 bootstrap draws. Separately, a structured calculation result may contain
endpoints with the failed endpoint state. For every combination of PPML or Harvey QML, slack value,
and bootstrap-index family, pool both endpoints, all \(B\) bootstrap draws, and all \(d_\eta=5\)
reported coefficients. Here \(d_\eta\) is the positive integer number of reported conditional-scale
coefficients. At most \(0.25(2Bd_\eta)\) pooled endpoints may have the failed endpoint
state. The slack values in this rule are \(0\), \(0.05\), \(0.10\), and \(0.20\). At zero slack both
sides of a bootstrap draw are exact copies of its zero-slack point, so that point's endpoint state
enters the pool twice. Breaching either the bootstrap-draw limit or this pooled limit stops the
reproduction; neither is a warning. The \emph{stability denominator} is the number of endpoints with bounded, unbounded, or
unreliable endpoint states used to calculate the proportion having the bounded endpoint state. Exclude
endpoints with the failed endpoint state from this denominator.

\section{Bootstrap confidence intervals from identified endpoints}

This section converts sampling variation in the lower and upper identified endpoints into a
confidence interval for one coefficient at one positive slack. Standardizing the two endpoint
deviations permits a common critical value and accommodates different lower- and upper-endpoint
scales.

Select one coefficient and one positive slack. Let \(\widehat\phi_{\mathrm{lo}}\) and
\(\widehat\phi_{\mathrm{hi}}\) be its lower and upper identified endpoints from the original sample.
Mean-equation and conditional-scale coefficients obtain \(\widehat\phi_{\mathrm{lo}}\) and
\(\widehat\phi_{\mathrm{hi}}\) differently. For a mean-equation coefficient they are the
identified endpoints of Section~\ref{sec:idset}. For a conditional-scale coefficient they
are obtained instead by carrying out the per-bootstrap-draw calculation of
Section~\ref{sec:bootstrap} once on the original sample, unresampled, under exactly the per-draw
budgets stated there: a lattice cap of 1,500 retained points, a fresh-fit allowance of 8,000, and no
cold-start endpoint replication. Do not substitute the full-budget conditional-scale identified
endpoints, which use a 4,000-point lattice cap, a fresh-fit allowance of 20,000, and cold-start
endpoint replication. The zero-slack conditional-scale coefficient used further below comes from
this same unresampled calculation. For
bootstrap draw \(b\in\{1,\ldots,B\}\), let \(\phi_{\mathrm{lo},b}\) and \(\phi_{\mathrm{hi},b}\) be its bootstrap
endpoints when the corresponding endpoint state is bounded. Let
\begin{equation*}
 s_{\mathrm{lo}}=\operatorname{MAD}(\phi_{\mathrm{lo},b}),\qquad
 s_{\mathrm{hi}}=\operatorname{MAD}(\phi_{\mathrm{hi},b}).
\end{equation*}
Here \(s_{\mathrm{lo}}\) and \(s_{\mathrm{hi}}\) are the nonnegative scalar MAD scales. Each is
taken over that side's \emph{usable} bootstrap draws: those whose endpoint state for that side is
bounded \emph{and} whose endpoint value for that side is finite. The two conditions are separate
requirements, not one; every rule below that speaks of bounded draws for a side means usable draws in
this sense. Require both scales to be finite and strictly
positive. Require each endpoint to have at least \(\lceil0.5B\rceil\) usable bootstrap draws. Among
bootstrap draws whose endpoint state is bounded, unbounded, or unreliable for that endpoint, require
at least 85 percent to be usable.
Exclude bootstrap draws with a failed endpoint state from this stability denominator, but retain them in the
requirement of at least \(\lceil0.5B\rceil\) bootstrap draws.

Define the standardized endpoint deviations
\begin{equation*}
 Z_{\mathrm{lo},b}=\frac{\phi_{\mathrm{lo},b}-\widehat\phi_{\mathrm{lo}}}{s_{\mathrm{lo}}},\qquad
 Z_{\mathrm{hi},b}=\frac{\widehat\phi_{\mathrm{hi}}-\phi_{\mathrm{hi},b}}{s_{\mathrm{hi}}}.
\end{equation*}
Here \(Z_{\mathrm{lo},b}\) and \(Z_{\mathrm{hi},b}\) are the scalar standardized lower- and
upper-endpoint deviations in bootstrap draw \(b\). Each is defined on the bootstrap draws whose
endpoint state is bounded on its own side, that is, on the same bootstrap draws that enter that side's
scale. Both are positive when the bootstrap draw's identified range is narrower than the
original-sample one on the side in question. Call a bootstrap draw \emph{paired} when both of its
endpoints have the bounded endpoint state. Every quantile of the two-sided construction that follows
is taken over the paired bootstrap draws only. The two scales are not: each of
\(s_{\mathrm{lo}}\) and \(s_{\mathrm{hi}}\) also uses bootstrap draws bounded on its own side alone,
which the paired set excludes. For that two-sided construction, require at least
\(\lceil0.5B\rceil\) paired values. Five requirements stop the reproduction when they fail for a
reported coefficient at a displayed positive slack: the two scale requirements, the usable-draw count
on either side, the 85 percent usable share on either side, and the paired-value count. On any of
these, stop and report the coefficient, the slack, the requirement that failed, the usable share on
each side, and the number of paired bootstrap draws, rather than leaving the cell blank and
continuing. Requirements stated later in this section for the two comparison calibrations do not
stop it; failing them leaves those comparison values missing. For any finite sequence of length \(n_q\), the
conservative \(1-\alpha\) empirical quantile is its \(k_q\)-th order statistic with
\begin{equation}
 k_q=\min\{n_q,\lceil(n_q+1)(1-\alpha)\rceil\},
 \qquad \alpha=0.10.
 \label{eq:conservative-order}
\end{equation}
Here \(\alpha=0.10\) is the scalar tail probability, and \(n_q\) is the sequence length. The
positive integer \(k_q\) is the order-statistic index. The \(k_q\)-th order statistic is the
\(k_q\)-th smallest value after sorting the sequence.

The reported
\emph{pointwise-coverage critical value} applies separately to one coefficient and one slack value.
It uses the scalar interpolation weight
\(\lambda\in[0,1]\) and
\begin{equation*}
 \mathcal C_b(\lambda)=\max\left\{0,
 Z_{\mathrm{lo},b}-\lambda\frac{\widehat\phi_{\mathrm{hi}}-\widehat\phi_{\mathrm{lo}}}{s_{\mathrm{lo}}},
 Z_{\mathrm{hi},b}-(1-\lambda)\frac{\widehat\phi_{\mathrm{hi}}-\widehat\phi_{\mathrm{lo}}}{s_{\mathrm{hi}}}\right\}.
\end{equation*}
Here \(\mathcal C_b(\lambda)\ge0\) is the scalar pointwise endpoint-deviation target for
bootstrap draw \(b\) at interpolation parameter \(\lambda\in[0,1]\).

For each \(\lambda\), let \(g(\lambda)\) be the conservative quantile of
\(\mathcal C_b(\lambda)\) defined by equation~\eqref{eq:conservative-order}. Define
\begin{equation*}
 L=\max\left\{
 \frac{\widehat\phi_{\mathrm{hi}}-\widehat\phi_{\mathrm{lo}}}{s_{\mathrm{lo}}},
 \frac{\widehat\phi_{\mathrm{hi}}-\widehat\phi_{\mathrm{lo}}}{s_{\mathrm{hi}}}
 \right\}.
\end{equation*}
Here \(L\ge0\) is a scalar Lipschitz constant: changing \(\lambda\) by \(d\ge0\) changes
\(g(\lambda)\) by at most \(Ld\). For a retained interval \([\lambda^{\flat},\lambda^{\sharp}]\subseteq[0,1]\), evaluate
\(g(\lambda^{\flat})\) and \(g(\lambda^{\sharp})\) and define its certified upper bound as
\begin{equation*}
 U(\lambda^{\flat},\lambda^{\sharp})=\frac{g(\lambda^{\flat})+g(\lambda^{\sharp})+L(\lambda^{\sharp}-\lambda^{\flat})}{2}.
\end{equation*}
Here \(U(\lambda^{\flat},\lambda^{\sharp})\) is the scalar upper bound for the interval
\([\lambda^{\flat},\lambda^{\sharp}]\). The best evaluated value is
\begin{equation*}
 G=\max_{\lambda\ \mathrm{evaluated}}g(\lambda),
\end{equation*}
where \(G\) is a scalar lower bound on \(\sup_{\lambda\in[0,1]}g(\lambda)\).

Let \(c_S\) be the conservative critical value from
\(\max(0,Z_{\mathrm{lo},b},Z_{\mathrm{hi},b})\). If \(L\) is zero or is not finite, perform no
search: set \(c_{\mathrm{CI}}=c_S\) and go straight to the reported interval below. Otherwise
proceed as follows.

Begin the branch-and-bound search with the single retained interval \([0,1]\). Evaluate \(g(0)\) and \(g(1)\), so that at the
outset \(G=\max\{g(0),g(1)\}\). Then repeat these two steps in this order.
\begin{enumerate}
\item Form \(\max_{[\lambda^{\flat},\lambda^{\sharp}]}U(\lambda^{\flat},\lambda^{\sharp})\) over the retained intervals. If
\begin{equation*}
 \min\left\{c_S,\max_{[\lambda^{\flat},\lambda^{\sharp}]}U(\lambda^{\flat},\lambda^{\sharp})\right\}-G\le10^{-4},
\end{equation*}
stop, and carry this same value of \(\max_{[\lambda^{\flat},\lambda^{\sharp}]}U(\lambda^{\flat},\lambda^{\sharp})\) into the critical value below.
\item Otherwise choose the retained interval with the largest \(U(\lambda^{\flat},\lambda^{\sharp})\), breaking a tie by the
smallest left endpoint. Split it at \((\lambda^{\flat}+\lambda^{\sharp})/2\), evaluate \(g\) there, replace the parent interval
with its two children, and update \(G\).
\end{enumerate}
The stopping test therefore precedes the first split, and a search that stops at once has evaluated
\(g\) only at \(\lambda=0\) and \(\lambda=1\). The certified upper endpoint is
\begin{equation*}
 c_{\mathrm{CI}}=\min\left\{c_S,\max_{[\lambda^{\flat},\lambda^{\sharp}]}U(\lambda^{\flat},\lambda^{\sharp})\right\}.
\end{equation*}
Here \(c_{\mathrm{CI}}\ge0\) is the scalar confidence-interval critical value. Record alongside it
the interpolation weight \(\lambda\) attaining \(G\), and record the search as interior when
\(G\) exceeds \(\max\{g(0),g(1)\}\) by more than \(10^{-4}\) and as not interior otherwise.

Record one further critical value, formed once for each conditional-scale estimator at each positive
display slack rather than once per coefficient. Within that estimator and slack, call an endpoint
side \emph{contributing} when its original-sample endpoint state is bounded and it passes the
regularity requirements above. Pool the bootstrap draws that are bounded with a finite value on every
contributing side at once. For each bootstrap draw in that pool, take the larger of zero and the
largest standardized endpoint deviation over all contributing sides of all five reported
conditional-scale coefficients, and apply equation~\eqref{eq:conservative-order} to that quantity. Call the result the
\emph{simultaneous-coverage critical value}. If no side contributes, it is missing. It is recorded
only in the machine-readable conditional-scale inference table, reaches no displayed table cell, and
never replaces \(c_{\mathrm{CI}}\) in any interval.

Report the open confidence
interval
\begin{equation*}
 (\widehat\phi_{\mathrm{lo}}-c_{\mathrm{CI}}s_{\mathrm{lo}},\;
 \widehat\phi_{\mathrm{hi}}+c_{\mathrm{CI}}s_{\mathrm{hi}}).
\end{equation*}
The two displayed quantities are the lower and upper confidence-interval endpoints. An open
confidence interval excludes these endpoints. Record the two halves of the branch-and-bound result
separately: the attained lower bound \(G\) and the certified upper bound
\(\min\{c_S,\max_{[\lambda^{\flat},\lambda^{\sharp}]}U(\lambda^{\flat},\lambda^{\sharp})\}\). Record also the minimum usable draw count
\(\lceil0.5B\rceil\) that the requirements above impose, the gap between the certified upper bound
and \(G\), and the number of evaluations of \(g\) the search performed. At a positive display slack
each recorded per-side statistic is that side's original-sample endpoint divided by that side's MAD
scale, and is missing whenever the scale is not finite or not strictly positive; the
zero-slack statistic defined below occupies those positions only in the zero-slack rows. Record for
the zero-slack point, in addition, the MAD scale of its bootstrap draws, the analytic standard error
of the corresponding fit, and the ratio of the first to the second.

Three further interval families are recorded for comparison and enter no displayed interval. Compute
all three only when both endpoint sides pass their regularity requirements and both usable sets are
nonempty; otherwise leave all six values missing. With \(\alpha=0.10\), let \(q_{\mathrm{lo}}(p)\)
and \(q_{\mathrm{hi}}(p)\) be the \(p\) quantiles of the usable lower and upper bootstrap
endpoints, taken by linear interpolation: sort the \(n\) values in increasing order, set
\(\zeta_p=(n-1)p+1\), and interpolate linearly between the order statistics at
\(\lfloor\zeta_p\rfloor\) and \(\lceil\zeta_p\rceil\). This interpolation rule differs from the
conservative order statistic of equation~\eqref{eq:conservative-order}, which these three families do
not use. The \emph{normal-approximation interval} is
\begin{equation*}
 \left(\widehat\phi_{\mathrm{lo}}-\Phi^{-1}(1-\tfrac{\alpha}{2})s_{\mathrm{lo}},\;
 \widehat\phi_{\mathrm{hi}}+\Phi^{-1}(1-\tfrac{\alpha}{2})s_{\mathrm{hi}}\right).
\end{equation*}
The \emph{percentile interval} is
\(\left(q_{\mathrm{lo}}(\tfrac{\alpha}{2}),\,q_{\mathrm{hi}}(1-\tfrac{\alpha}{2})\right)\).
The \emph{basic interval} reflects each draw quantile through its original-sample endpoint,
\begin{equation*}
 \left(2\widehat\phi_{\mathrm{lo}}-q_{\mathrm{lo}}(1-\tfrac{\alpha}{2}),\;
 2\widehat\phi_{\mathrm{hi}}-q_{\mathrm{hi}}(\tfrac{\alpha}{2})\right),
\end{equation*}
so the basic and percentile intervals move in opposite directions when the bootstrap endpoints are
skewed.

At zero slack the lower and upper entries of a coefficient's row both carry that coefficient's
zero-slack point value, and both of its recorded statistics carry the zero-slack statistic \(t_0\)
defined below.
The construction just given applies when both original-sample endpoints have the bounded endpoint
state. Report a half-infinite confidence interval exactly when one original-sample endpoint has the
bounded endpoint state, the other has the unbounded endpoint state, and the bounded side passes the
regularity requirements above. Should that side fail any of them, the stopping rule already stated
applies: the reproduction stops rather than reporting a half-infinite interval. Call the bounded one the
\emph{live side}. On such a cell the truth may sit anywhere along the infinite ray, so the width
credit vanishes and no interpolation search is performed: take \(c_{\mathrm{CI}}\) to be the
conservative \(1-\alpha\) quantile of \(\max(0,Z_b)\), where \(Z_b\) is the live side's own
standardized endpoint deviation and the quantile is taken over the bootstrap draws whose live-side
endpoint state is bounded and whose live-side endpoint value is finite. The other endpoint's state
never restricts that set of bootstrap draws. Pad the live original-sample endpoint outward by
\(c_{\mathrm{CI}}\) times the live side's scale and leave the other side infinite. Suppress the
confidence interval when both original-sample endpoints have the unbounded endpoint state, and
suppress it whenever either original-sample endpoint has a state other than bounded or unbounded,
whatever the other endpoint's state.

Mean-equation coefficient tables do not display half-infinite identified ranges. Conditional-scale
tables may display them and their associated half-infinite confidence intervals.

Two calculations based on normal reference distributions serve only as comparison checks. Calculate
both for mean-equation coefficients only, and only at those coefficient-and-slack cells whose
original-sample identified range has strictly positive finite width and whose two endpoint scales are
both finite and strictly positive. At every other cell leave the two comparison critical values
blank; the paired endpoint correlation is recorded at every positive-slack displayed row regardless,
so a positive-slack zero-width cell carries a correlation beside two blank critical values. The
zero-slack rows carry no correlation and no comparison critical values: all three entries are
missing. For the
Imbens--Manski comparison, use the standard-normal cumulative distribution function \(\Phi\). Define
\(d_\phi=\max(\widehat\phi_{\mathrm{hi}}-\widehat\phi_{\mathrm{lo}},0)\) and
\(s_{\max}=\max(s_{\mathrm{lo}},s_{\mathrm{hi}})\). The
Imbens--Manski comparison critical value solves
\begin{equation*}
 \Phi(c_{IM}+d_\phi/s_{\max})-\Phi(-c_{IM})=1-\alpha.
\end{equation*}
Here \(c_{IM}\) is the scalar Imbens--Manski comparison critical value, \(d_\phi\ge0\) is the width of
the original identified range, and \(s_{\max}>0\) is the larger endpoint scale. Solve for \(c_{IM}\in[0,10]\)
to tolerance \(10^{-10}\).

For the Stoye comparison, pair the bootstrap draws whose lower and upper endpoint values are both
finite, and require at least eight such pairs. From these paired values alone, calculate a trimming
MAD for each endpoint as defined in equation~\eqref{eq:mad}, together with each endpoint's median
over the same pairs. Require both trimming MADs to be finite and positive. These trimming MADs are
distinct from \(s_{\mathrm{lo}}\) and \(s_{\mathrm{hi}}\), which the critical-value equation below
uses and which are taken over each side's own bounded bootstrap draws. Retain a pair only when each
endpoint's absolute deviation from its median is strictly less than six of its own trimming MADs,
and again require at least eight pairs. Let
\(r_{\mathrm{lo},\mathrm{hi}}\) be their ordinary sample correlation. Let
\((X,Y)\) be a standard bivariate-normal pair. The two variables are jointly normal, and each has mean
zero and variance one. Set their correlation to \(r_{\mathrm{lo},\mathrm{hi}}\). The Stoye comparison
critical value solves
\begin{equation*}
 \min\left\{
 \begin{aligned}
  &\Pr(X\le c_{St},\;Y\ge-c_{St}-d_\phi/s_{\mathrm{hi}}),\\
  &\Pr(X\le c_{St},\;Y\ge-c_{St}-d_\phi/s_{\mathrm{lo}})
 \end{aligned}
 \right\}=1-\alpha.
\end{equation*}
Here \(c_{St}\) is the scalar Stoye comparison critical value, \(X\) and \(Y\) are the standard
bivariate-normal variables, and \(r_{\mathrm{lo},\mathrm{hi}}\) is their correlation. Solve for
\(c_{St}\in[10^{-6},8]\) to tolerance \(10^{-8}\). For \(|r_{\mathrm{lo},\mathrm{hi}}|\le0.999\),
let \(x_0\in\mathbb R\) and \(y_0\in\mathbb R\) be arbitrary scalar arguments and calculate
\begin{equation*}
 \Pr(X\le x_0,Y\ge y_0)=
 \int_{-\infty}^{x_0}\varphi_N(t)
 \left[1-\Phi\left\{\frac{y_0-r_{\mathrm{lo},\mathrm{hi}}t}
 {\sqrt{1-r_{\mathrm{lo},\mathrm{hi}}^2}}\right\}\right]dt.
\end{equation*}
Here \(\varphi_N\) is the standard-normal density. Evaluate this improper integral with relative and
absolute tolerances \(10^{-9}\) and at most 100 subintervals. If
\(r_{\mathrm{lo},\mathrm{hi}}>0.999\), use
\(\max\{0,\Phi(x_0)-\Phi(y_0)\}\). If \(r_{\mathrm{lo},\mathrm{hi}}<-0.999\), use
\(\Phi\{\min(x_0,-y_0)\}\). Neither comparison critical value
determines a reported confidence interval.

At zero slack, let \(\widehat\psi\) be an original-sample zero-slack point coefficient. Require
\(\widehat\psi\) itself to be finite. Retain only
finite bootstrap values \(\psi_b\) with the bounded endpoint state. Require at least
\(\lceil0.5B\rceil\) retained bootstrap draws. They must represent at least 85 percent of all zero-slack bootstrap draws
without the failed endpoint state, and \(\operatorname{MAD}(\psi_b)\) must be finite and positive.
Let \(n_0\) be the positive integer number of retained finite bootstrap values with the bounded
endpoint state.
If any acceptance rule fails, do not report the statistic. Otherwise, define the
\emph{zero-slack bootstrap statistic} \(t_0\) and the
\emph{empirical p-value} \(p_0\) by
\begin{equation}
 t_0=\frac{\widehat\psi}{\operatorname{MAD}(\psi_b)},
 \qquad
 p_0=\frac{1+\sum_b
 \mathbf 1\{|\psi_b-\widehat\psi|\ge|\widehat\psi|\}}
 {1+n_0}.
 \label{eq:tau-zero-inference}
\end{equation}
Here \(t_0\) is the scalar zero-slack bootstrap statistic, \(p_0\in[0,1]\) is its empirical
p-value, \(\widehat\psi\) is the original-sample coefficient, and \(\psi_b\) is its retained
bootstrap value. The sum and \(n_0\) use only the retained bootstrap draws and
\(\mathbf 1\{\cdot\}\) is an indicator,
equal to one when its condition is true and zero otherwise. The p-value \(p_0\) is empirical;
it is not a normal approximation. An unbounded endpoint state is impossible for a zero-slack point.

Run the mean-equation endpoint bootstrap for the reported and comparison specifications using the
same primary bootstrap-index family. Compare their paired endpoints as a comparison check. Only the
reported specification reaches the reported outputs.

\section{Plug-in approximation-error variance bounds}
\label{sec:approximation}

A \emph{plug-in approximation-error variance bound} replaces unknown moments from the population
of possible observations with their sample counterparts to bound variance from omitted approximation terms.
These calculations use maturities \(i=3,6,\ldots,117\) and the \emph{approximation sample}, which is
the full available quarterly Treasury panel. They do not use the PC
window. They quantify omitted higher-order terms in the SDF approximations and do not
alter the predictors in the mean equation or any conditional-scale estimator mapping. Throughout this
section \(n_{i,t}\) is the expected log bond price of equation~\eqref{eq:nhat}. The approximation
sample runs from 1961Q2 through 2026Q2, one observation per quarter with no gaps, the last quarter
being the retained incomplete one. Let \(T_a\) be its observation count and let \(k_i=i/h\) be
the integer number of quarterly steps associated with maturity \(i\).

\subsection{SDF-news plug-in approximation-error variance bound}

Two candidate bounds bound the same omitted exponential remainder by different routes, and neither is
uniformly smaller than the other. The leading-term candidate bound uses fourth moments and an
exponential envelope. The cancellation-aware candidate bound retains the exact exponential remainders
after the first-order cancellation, discarding nothing, but pays the full unprojected carrier variance
in each of its two legs. Report the smaller candidate bound at each maturity.

For the leading-term candidate bound, let
\begin{equation*}
 \mathcal E_i^N=\max_{\substack{1\le t\le T_a-k_i\\ n_{i,t}\ \mathrm{not\ missing}}} e^{2n_{i,t}}.
\end{equation*}
Here \(\mathcal E_i^N\ge0\) is the scalar exponential envelope for maturity \(i\), and the maximum
uses its nonmissing maturity-trimmed dates. Let \(K_{1i}\) be the mean of
\(\{-\tfrac{h}{12}y_{h,t+k_i}/100-n_{i-h,t+1}\}^4\) over the indices \(t=1,\ldots,T_a-k_i\) at which
both \(y_{h,t+k_i}\) and \(n_{i-h,t+1}\) are not missing. Let \(K_{2i}\) be the
mean of \((\Delta q_{i,t+1})^4\) over its nonmissing values with
\(t=1,\ldots,T_a-k_i\). The three components \(\mathcal E_i^N\), \(K_{1i}\) and \(K_{2i}\) share the
index range \(t=1,\ldots,T_a-k_i\) and each applies its own nonmissing mask; the three masks are in
general different sets. At \(i=h\),
set \(K_{1i}=0\). The
leading-term candidate bound is
\begin{equation*}
 U_{i,\mathrm{lead}}^N=\frac14\mathcal E_i^N(K_{1i}+K_{2i}).
\end{equation*}
Here \(U_{i,\mathrm{lead}}^N\ge0\) is the scalar leading-term candidate bound,
\(\mathcal E_i^N\) is its exponential envelope, and \(K_{1i},K_{2i}\ge0\) are its fourth-moment
terms.

For the cancellation-aware candidate bound, retain indices \(t=1,\ldots,T_a-k_i\) for which
\(n_{i,t}\), \(x_{i,t}^h\), and \(n_{i,t}^{+}\) are finite. The two candidate bounds use different
retained sets: the leading-term components drop only missing values, so an infinite expected log bond
price is retained and makes the envelope infinite, whereas the cancellation-aware components drop
every nonfinite value. Evaluate every difference of the form \(e^{z}-1\) below by a rule that is exact
for small \(|z|\), rather than by subtracting one from a computed exponential. Here \(x_{i,t}^h\) is the
horizon-paired short-rate exponent and \(n_{i,t}^{+}\) is the next-quarter expected log bond price
at the shorter maturity:
\begin{equation*}
 x_{i,t}^h=-\frac{h}{12}\frac{y_{h,t+k_i}}{100},\qquad
 n_{i,t}^{+}=n_{i-h,t+1}.
\end{equation*}
Here \(x_{i,t}^h\) is the scalar horizon-paired short-rate exponent and \(n_{i,t}^+\) is the scalar
next-quarter expected log bond price at the shorter maturity. At \(i=h\), use
equation~\eqref{eq:price-news-boundary} for \(n_{i,t}^{+}\). Define
\begin{align*}
 \mathcal R_{0,i,t}^N&=e^{n_{i,t}}\{e^{x_{i,t}^h-n_{i,t}}-1-(x_{i,t}^h-n_{i,t})\},\\
 \mathcal R_{1,i,t}^N&=e^{n_{i,t}^{+}}\{e^{x_{i,t}^h-n_{i,t}^{+}}-1-(x_{i,t}^h-n_{i,t}^{+})\},\\
 D_{i,t}^{N}&=n_{i,t}^{+}-n_{i,t},\\
 \mathcal R_{g,i,t}^N&=e^{n_{i,t}}\{e^{D_{i,t}^{N}}-1-D_{i,t}^{N}-(D_{i,t}^{N})^2/2\}.
\end{align*}
Here \(\mathcal R_{0,i,t}^N\) and \(\mathcal R_{1,i,t}^N\) are scalar exponential remainders at
the two adjacent expected log bond prices. The scalar \(D_{i,t}^N\) is their log-price gap, and
\(\mathcal R_{g,i,t}^N\) is the scalar higher-order gap remainder.

Let \(n_i^C\) be the observation count after retaining quarters for which all input quantities are
finite. Do not delete
a retained quarter merely because a transformed quantity fails to be finite; instead, set that
candidate bound to infinity. Replace a negative computed value of a divisor-\(n_i^C\) centered
variance by zero before taking its square root. With divisor-\(n_i^C\) centered standard deviations on
this retained set, define
\begin{equation*}
 U_{i,\mathrm{cancel}}^N=
 \{\operatorname{sd}_{n_i^C}(\mathcal R_{1,i}^N)+\operatorname{sd}_{n_i^C}(\mathcal R_{0,i}^N)
   +\operatorname{sd}_{n_i^C}(\mathcal R_{g,i}^N)\}^2.
\end{equation*}
Here \(U_{i,\mathrm{cancel}}^N\ge0\) is the scalar cancellation-aware candidate bound for maturity
\(i\).

The reported SDF-news plug-in approximation-error variance bound is the pointwise minimum
\begin{equation*}
 U_i^N=\min\{U_{i,\mathrm{lead}}^N,U_{i,\mathrm{cancel}}^N\}.
\end{equation*}
Here \(U_i^N\ge0\) is the scalar reported SDF-news plug-in approximation-error variance bound at
maturity \(i\). Pointwise means that the minimum is taken separately at each maturity. Both candidate
bounds are plug-in estimates of population approximation-error variance bounds rather than literal
finite-sample upper bounds; the leading-term candidate additionally discards a higher-order remainder
and a cross term. Both are measured in squared units of the stochastic discount factor level, which is dimensionless; \(\mathcal E_i^N\) is in squared units of a bond price.

If no quarter is retained for a component, the corresponding candidate bound is undefined; stop the
reproduction if that state is reached. At every one of the 39 maturities the leading-term candidate
must be finite and strictly positive, the cancellation-aware candidate must be strictly positive but is permitted to be infinite, and the reported minimum must be finite and strictly positive. Stop the
reproduction if any of these fails.

\subsection{Paired expected-SDF plug-in approximation-error variance bound}

The paired expected-SDF plug-in approximation-error variance bound measures the exponential
remainder after pairing the short rate with the maturity horizon. It differs from the expected-SDF
series, which uses an unpaired mean correction. For \(t=1,\ldots,T_a-k_i\), reuse the horizon-paired
short-rate exponent \(x_{i,t}^h\) defined above. The expected-SDF level gap is
\(D_{i,t}^{E,\mathrm{lev}}=e^{x_{i,t}^h}-e^{n_{i,t}}\). Retain exactly the
quarters in which \(D_{i,t}^{E,\mathrm{lev}}\) is finite and let \(n_i^E\) be their observation
count. On this one common set, define
\begin{equation*}
 D_{i,t}^{E,\mathrm{log}}=x_{i,t}^h-n_{i,t},\qquad
 \mathcal R_{i,t}^{E}=e^{n_{i,t}}(e^{D_{i,t}^{E,\mathrm{log}}}-1-D_{i,t}^{E,\mathrm{log}}).
\end{equation*}
Here \(D_{i,t}^{E,\mathrm{log}}\) is the scalar expected log bond-price gap and
\(\mathcal R_{i,t}^E\) is its scalar exponential remainder.

If either candidate bound is nonfinite on this common set, treat that candidate bound as infinite rather
than deleting further rows.
Both the maximum and the fourth-moment mean below run over the \(n_i^E\) retained quarters and over no
others. Then calculate
\begin{equation*}
 U_i^E=\min\left\{
 \frac14\max_t e^{2n_{i,t}}\operatorname{mean}_t (D_{i,t}^{E,\mathrm{log}})^4,
 \widehat{\operatorname{Var}}_{n_i^E}(\mathcal R_i^E)
 \right\}.
\end{equation*}
Here \(U_i^E\ge0\) is the scalar paired expected-SDF plug-in approximation-error variance bound at
maturity \(i\), in squared units of the stochastic discount factor level. Like the two SDF-news
candidates, it is a plug-in estimate of a population bound and not a literal finite-sample upper
bound. It must be finite and strictly positive at every one of the 39 maturities; stop the
reproduction otherwise. The resulting \(U_i^E\)
is a paired expected-SDF plug-in approximation-error variance bound; it does not automatically bound
the error in the three lagged expected-SDF PC scores.

For each of the two reported sequences \(U_i^N\) and \(U_i^E\), summarize the mean, median, minimum,
maximum, and ordinary sample standard deviation across the 39 maturities. Do not summarize the two
SDF-news candidate bounds separately. Let
\(B_{\max}\) denote the maximum approximation-error standard deviation in basis points. Calculate it
as
\begin{equation*}
 B_{\max}=10^4\max_i\{\sqrt{U_i^N},\sqrt{U_i^E}\}.
\end{equation*}
Here \(B_{\max}\ge0\) is the scalar maximum across maturities and across the two plug-in
approximation-error variance bounds. One basis point is one ten-thousandth of the stochastic discount
factor level, which is the decimal unit in which \(\sqrt{U_i^N}\) and \(\sqrt{U_i^E}\) are measured.

For the quoted-yield calculation, let \(y_{h,t}/400\) be the annualized three-month yield converted
from percentage points to a one-quarter decimal rate. Use all \(T_a\) quarters of the approximation
sample, in which the three-month yield must have no missing value, giving \(T_a-1\) first differences.
Its typical quarterly move is the ordinary sample standard deviation of those first differences, and
must be finite and strictly positive. Divide the maximum error standard deviation in decimal units by
that typical move and multiply by 100. The comparison is meaningful because a move of the one-quarter
yield by \(\mathrm dy\) moves the one-quarter bond price, and hence the stochastic discount factor
level, by approximately \(\mathrm dy\) times a price close to one.

The \emph{quoted-number acceptance rules} compare three reproduced quantities with fixed acceptance
values and impose two additional inequalities. The fixed acceptance values are \(4.19\times10^{-9}\) for the maximum SDF-news plug-in approximation-error variance bound,
\(3.19\times10^{-9}\) for the maximum paired expected-SDF plug-in approximation-error variance bound, and 0.00239 for the typical quarterly
yield move. Use absolute tolerance \(5\times10^{-12}\) for each variance bound and
\(5\times10^{-6}\) for the typical quarterly yield move. The maximum approximation-error standard
deviation must also remain below 0.7 basis points and
below 3 percent of the typical move. Check every one of these five rules before producing any
approximation-error output, and stop the reproduction if any of them fails. These five rules are predeclared,
not newly observed results in this manual.

\section{Reported outputs and formatting}
\label{sec:outputs}

A reproduction is complete only after it constructs every required reported output.

The reproduction emits exhibit bodies, not finished exhibits. Of the five published tables, only the
conditional-scale estimator table pages carry their own title, cross-reference label, and note; the
other four table bodies and the three descriptive tables are emitted without any of the three, and no
figure carries a title. The document that displays an exhibit supplies whatever the reproduction
does not.

Wherever the title and note are supplied, they must be self-contained. Use plain-language row and
column names, and define every series and every numerical row in the note. State the sample, units,
observation count, specification, controls, and coefficient of determination whenever they apply. An
inferential table must identify the covariance or bootstrap rule, the reference distribution, the
null hypothesis, and the p-value or confidence level, and its note must define the significance
stars. A figure's title and note must define its symbols, sample, and units, name each axis in plain
language with its units, and state whether a displayed set is an attained lattice projection, a
numerically certified range, or a confidence interval. These presentation requirements apply in
addition to the validation requirements in this section.

\begin{enumerate}
\item Produce a descriptive report over the following 22 series, in this order: scalar consumption
growth; the three current expected-SDF PC scores; the three lagged expected-SDF PC scores; the three
SDF-news PC scores; the expected SDF at maturities 3, 12, 60, and 117 months; the SDF news at the
same four maturities; and realized yield volatility at the same four maturities. This order fixes the
rows of the summary table, the rows and columns of the correlation matrix, and the panel order of
both figure grids. The report contains, for each series, the nonmissing count, mean, ordinary sample
standard deviation, minimum, median, and maximum. Add pairwise-complete correlations, each calculated
from all rows in which both series are observed, the
preliminary mean regression in equation~\eqref{eq:prelim-ols} with OLS standard errors, one
time-series plot, and a histogram for each series. Use 30 bars when the finite observations vary
and one bar when they are all equal. For a series with \(n\) finite observations,
where \(n\) is a positive integer, denote its minimum by \(x_{(1)}\) and its maximum by
\(x_{(n)}\). Set \(w=(x_{(n)}-x_{(1)})/29\). When \(w>0\), use the 31 boundaries
\(b_j=x_{(1)}-w/2+jw\), for \(j=0,\ldots,30\). Here \(w>0\) is the scalar interval width, and
\(b_j\) is the \(j\)th scalar boundary. Draw bar \(j\) from \(b_{j-1}\) to \(b_j\), for
\(j=1,\ldots,30\). To assign observations, set \(\delta_h=10^{-8}w\), define the adjusted first
boundary as \(\widetilde b_0=b_0-\delta_h\), and define every later adjusted boundary as
\(\widetilde b_j=b_j+\delta_h\). Here \(\delta_h>0\) is the scalar boundary adjustment. The first adjusted
interval includes both endpoints; every later adjusted interval excludes its lower endpoint and
includes its upper endpoint. Thus an observation no more than \(\delta_h\) above an interior nominal
boundary enters the interval on its left. When all finite observations equal the scalar value \(c\),
set the lower boundary to \(0.05+0.1\lfloor(c-0.05)/0.1\rfloor\) and the upper boundary to
\(0.15+0.1\lfloor(c-0.05)/0.1\rfloor\).
Use \(\delta_h=10^{-9}\) for assignment and apply the same endpoint rule. Omit missing observations from the
counts. The height of a bar is its assigned finite-observation count. Pairwise correlations may use
observation counts that differ from the per-series summary counts. Print the summary-statistics
table, the preliminary-regression table, and the reported coefficient of determination and adjusted coefficient of determination to three decimal
places, and the correlation matrix to two decimal places. The adjusted coefficient of determination
is one minus the ratio of the residual sum of squares divided by \(N_{\mathrm{pm}}-7\) to the total
sum of squares divided by \(N_{\mathrm{pm}}-1\), where seven is the number of estimated
coefficients. The preliminary mean regression is not a substitute for the mean least-squares
benchmark's four-lag covariance calculation.
\item Produce five published tables. The table identified by \nolinkurl{hetero_tests} reports the
mean-equation identification tests and checks. The table identified by
\nolinkurl{structural_var_inference} reports combined mean-equation and conditional-scale inference.
The table identified by \nolinkurl{structural_var_estimators} combines mean-equation and
conditional-scale results for the conditional-scale estimators. The table identified by
\nolinkurl{var_share} reports variance shares. The table identified by
\nolinkurl{variance_bounds_summary} reports the plug-in approximation-error variance-bound
summaries. Accompany each published table with a standalone typesetting source and a
standalone rendered document.
Standalone means that the typesetting source produces its rendered document without being embedded
in another document. The
comparison between the reported and comparison specifications is a comparison check. Critical slack appears
in the slack-dependent calculations and figures, not in an additional published table.
\item Produce one figure of identified ranges of mean-equation coefficients as a function of slack,
and two projected figures of mean-equation coefficients. Draw the first on the \emph{bounds-by-slack
grid}, built once from the critical slack \(\tau^\star\) and shared by the mean-equation stage and
the figure. Form 25 equally spaced values from zero through \(\tau^\star\) and discard both
endpoints, leaving the strictly interior backbone of 23 values. Split the backbone at
\(0.9\tau^\star\): the values strictly below that point form the coarse part, and the values at or
above it form the tail. Both parts must be nonempty. Take the largest coarse value followed by every
tail value as refinement endpoints, divide each adjacent interval between them into four equal
subintervals, and keep every division point of that interval except its left endpoint, so each
interval contributes three new interior points together with its own right endpoint. The grid is the
sorted union of the coarse part and the kept points. It is strictly increasing, every value lies
strictly between zero and \(\tau^\star\), and it contains the whole backbone, so its largest value
equals the largest backbone value exactly and refinement never extends the grid's range. Under the
current controls the coarse part holds 21 values and the refinement contributes 8, giving 29 slacks.
Run the widening search of Section~\ref{sec:idset} at every grid slack in increasing order, carrying
the accepted coefficient vectors forward from one grid slack to the next as warm starting values and
starting the first from the zero-slack point when that point is available; this widens only the
identified ranges of the structural SDF-news coefficients, leaving those of the design coefficients as
the origin-started solve returns them.
Plot two further slacks alongside the grid, recovered from the stored tables rather than re-solved: zero
slack, entered as the degenerate range whose lower and upper values both equal the zero-slack point,
and the baseline slack, entered as the identified range stored for that slack. Drop every coefficient-and-slack
row whose endpoint state is not bounded, so each band is truncated where certification ends, and mark
the baseline slack with a vertical reference line. Also produce three-dimensional figures of
the mean identified set at slacks 0.05, 0.10, 0.20, and 0.30. At each slack, render both coefficient units and standard-deviation-scaled
units, each with and without the mean least-squares benchmark. Label every finite rendering as
a lattice projection of the mean identified set rather than as a globally certified mean identified set.
\item Produce figures of conditional-scale coefficient identified ranges for log-OLS, PPML, Harvey
QML, and LAD. The associated table pages follow these inference rules. In
mean least-squares benchmark reference columns, log-OLS uses standard errors from its four-lag Bartlett analytic
covariance and Student \(t\) p-values. PPML and Harvey QML use standard errors from their four-lag
Bartlett analytic covariances and standard-normal p-values. In zero-slack
point columns, PPML and Harvey QML use the zero-slack bootstrap statistic and empirical p-value
defined in equation~\eqref{eq:tau-zero-inference}; log-OLS and LAD leave those entries blank.
Positive-slack PPML and Harvey QML mean-identified-set columns contain bootstrap confidence intervals.
The log-OLS and LAD pages have no row for a bootstrap confidence interval.
\item Produce quarter-indexed envelopes of the fitted conditional standard deviation for PPML and
Harvey QML\@. Draw these baseline envelopes at the baseline slack 0.05, reading that slack's
coefficient box from the identified ranges stored for that slack rather than solving it afresh.
For each quarter, minimize and maximize the fitted log scale over the mean
identified set, then apply the increasing exponential transformation. Also construct an envelope for
the fitted conditional median absolute-residual scale, by the separate rule stated below: it is the
range attained on a finite feasible sample of the mean identified set rather than a searched
projection of that set.

In addition to the baseline figures, render sweep figures of fitted conditional standard deviations
for PPML and Harvey QML\@. Solve the mean-equation coefficient boxes for these figures over the
sorted union of the display slacks and the sweep slacks, running the same warm-start endpoint
refinement of Section~\ref{sec:idset} at every slack of that union, so a sweep slack off the display
grid receives a full endpoint solve rather than an interpolated box. Drop any sweep slack that is not
below the critical slack \(\tau^\star\), and render no figure for it. The sweep slacks are 0.05, 0.10, 0.20,
0.30, and 0.40. Also render combined level and logarithmic-value figures and combined normalized
lower- and upper-endpoint figures. These fitted envelopes provide pointwise rather than simultaneous
coverage across quarters. Check that an available fitted path at the zero-slack point lies inside
every corresponding two-sided envelope.
\item Produce machine-readable tables for mean-equation inference and conditional-scale inference.
Also produce them for the comparison check between the reported and comparison specifications, the joint-null compatibility witness, and the stacked-moment replication check.
Produce a machine-readable table of the LAD one-sided closure limits at verified structural-mean-residual
crossings, which are approximate limits rather than attained endpoints. Preserve the bootstrap draws,
the residual-dynamics outcome, the joint-null and stacked-moment results, the dynamic-volatility
routing outcome, and the route-status result as exact reusable results.
\item Produce a plug-in approximation-error variance bound figure and summary table across the 39 maturities spaced three months apart,
plus a machine-readable table and a prose record of the quoted-number acceptance rules.
\end{enumerate}

For standard-deviation-scaled mean identified set figures, let
\(s_k^N=\operatorname{sd}_{T-1}(\{W_{2k,t}\}_{t=1}^{T})>0\) be the ordinary sample standard deviation
of coordinate \(k\) of the SDF-news reduced-form residual vector over the identification sample, and
let \(D_N=\operatorname{diag}(s_1^N,s_2^N,s_3^N)\). Here \(D_N\in\mathbb R^{3\times3}\) is the
diagonal matrix of SDF-news reduced-form residual standard deviations on the identification sample.
Plot the scaled coordinate \(\vartheta=D_N\theta\), where \(\vartheta\in\mathbb R^3\). In these
coordinates, replace each quadratic inequality's terms by
\begin{equation*}
 A_k^*=D_N^{-1}A_kD_N^{-1},\qquad b_k^*=D_N^{-1}b_k,\qquad c_k^*=c_k.
\end{equation*}
Here \(A_k^*,b_k^*,c_k^*\) are the quadratic, linear, and constant terms in scaled coordinates.

For a two-coordinate projection, write inequality \(k\), conditional on displayed coordinates
\((x,y)\), as
\begin{equation*}
 a_k\theta_u^2+d_k(x,y)\theta_u+\chi_k(x,y)\le0,
\end{equation*}
where \(\theta_u\in\mathbb R\) is the undisplayed coefficient, and
\(a_k,d_k(x,y),\chi_k(x,y)\in\mathbb R\) are the quadratic, linear, and constant coefficients in
\(\theta_u\). Require \(a_k>0\); stop the reproduction if any \(a_k\le0\). When the discriminant
\(d_k(x,y)^2-4a_k\chi_k(x,y)\) is nonnegative, let \(L_k(x,y)\) and \(H_k(x,y)\) be the smaller and
larger quadratic roots. When the discriminant is negative, set that inequality's lower root to
\(10^{6}\) and its upper root to \(-10^{6}\), which empties its feasible interval and excludes the
point while keeping the difference \(L(x,y)-H(x,y)\) defined below finite, so that the contour drawn
from it stays well defined.
Define \(L(x,y)=\max_kL_k(x,y)\) and \(H(x,y)=\min_kH_k(x,y)\). Include \((x,y)\) in the exact
projection precisely when \(L(x,y)\le H(x,y)\) and draw its boundary where
\(L(x,y)-H(x,y)=0\). Evaluate each displayed projection on a 300-by-300 lattice. Draw three panels in a single row. Panel \(j\) projects out coordinate \(j\) and displays the
remaining two in increasing order, so the panels show the coordinate pairs \((\theta_2,\theta_3)\),
\((\theta_1,\theta_3)\), and \((\theta_1,\theta_2)\) from left to right, and each coordinate appears
in exactly two panels. The displayed projection slacks are 0.05, 0.10, and 0.20, so the largest
is 0.20. All displayed
panels share one horizontal range and one vertical range. Take the coefficient box of the mean
identified set at the largest displayed projection slack. The horizontal range spans the box extents
of the two coefficients that ever appear on a horizontal axis, the vertical range spans the box
extents of the two that ever appear on a vertical axis, and each range is widened by six percent of
its width at both ends. In every panel and at every displayed projection slack, also draw the rectangle whose
side in each displayed coordinate is that coordinate's \emph{drawn projection extent}: the
union, over the two panels that display the coordinate, of the extent of the drawn projection
boundary in that coordinate.

For the PPML and Harvey QML fitted conditional-scale envelopes and each quarter \(t\), evaluate the fitted log scale
\(f_t(\theta)=X_t^{V\prime}\widehat\eta(\theta)\) at every lattice point retained after thinning that
has an accepted fit. Let \(f_{t,\min}^{\mathrm{lat}}\) and \(f_{t,\max}^{\mathrm{lat}}\) be the
smallest and largest finite values and call their attaining points the lower and upper attained
lattice extremes for quarter \(t\). Define the positive fitted-log-scale lattice-extreme scale
\begin{equation*}
 g_t^f=\max\{1,|f_{t,\min}^{\mathrm{lat}}|,|f_{t,\max}^{\mathrm{lat}}|\}.
\end{equation*}
Minimize and maximize \(f_t(\theta)\) over the mean identified set. Use the shared conditional-scale
endpoint-refinement procedure with the corresponding attained lattice extreme and search start point, and no
further lattice-derived starting values. Replace \(g_j\) by \(g_t^f\) in the
unreliable endpoint-state rule. Use a fresh-fit allowance of 80,000 at each slack. Convert the fitted log scale to a standard deviation with
\(\exp(X_t^{V\prime}\eta/2)\). A finite zero-slack fitted value lies inside its two-sided envelope when its
distance beyond either endpoint is no more than \(10^{-6}\max(1,|s_t^0|)\), where \(s_t^0>0\) is
the zero-slack fitted scale.

Build the fitted conditional median absolute-residual envelope by a different rule, which uses no
part of the endpoint-refinement procedure above. Form a fresh lattice of 31 equally spaced values
along each structural SDF-news coefficient, spanning that coefficient's range in the baseline-slack
coefficient box, and retain the points satisfying \(\rho_Q(\theta)\le10^{-10}\). Thin the retained
points to a cap of 3,000 by the thinning rule of Section~\ref{sec:logvar-scale}, which keeps at most
that many and in general fewer. Then prepend the zero-slack point when it satisfies every inequality
of the baseline-slack system with every inequality scale \(\omega_k\) set to one; do not append it. Fit the
conditional median absolute-residual estimator afresh at every retained point, carrying no warm
starting value from one point to the next, and keep a point only when its fit is accepted with every
coefficient finite. At least one point must be kept; stop the reproduction otherwise. For each
quarter \(t\), the lower and upper envelope values are the smallest and largest of
\(f_t(\theta)=X_t^{V\prime}\widehat\eta(\theta)\) over the kept points, so this envelope is the range
attained on that finite sample and lies inside the exact projection of the mean identified set rather
than reproducing it. Assign both sides the bounded endpoint state, downgrading a side to the
unreliable endpoint state when either \(\exp(X_t^{V\prime}\eta)\) or \(\exp(X_t^{V\prime}\eta/2)\)
fails to be finite, of which the first is the binding condition at any finite log value. Carry out no
endpoint refinement, no lattice-extreme scale, no cold-start endpoint replication, no fresh-fit allowance and
no containment check: prepending a feasible zero-slack point already places its fitted path inside
the envelope, and when that point is infeasible or its fit is not accepted the zero-slack path is
missing and there is nothing to contain. Convert the fitted log scale to the fitted conditional median
absolute-residual scale with the same \(\exp(X_t^{V\prime}\eta/2)\). That quantity is a conditional
median absolute residual, not a standard deviation, because the median regression applies no
normality scale correction.

A combined logarithmic figure changes only its vertical axis to base-ten
logarithms; the plotted quantity remains the fitted scale. For a normalized lower or upper endpoint
figure, standardize each endpoint series and the zero-slack series separately by its own sample mean
and ordinary sample standard deviation, retaining missing gaps.

\begin{sloppypar}
The \emph{output root} is the top-level folder containing all reported outputs. The
\emph{relative output location} is a reported output's folder and basename below the output root.
The basename is the final name, including its extension. The literal basenames below are fixed output
identifiers; their abbreviations do not change the terminology used for the mathematical concepts in
this manual. The extension \nolinkurl{.tex} denotes a typesetting source, \nolinkurl{.pdf} a rendered
document, \nolinkurl{.csv} a comma-separated table, \nolinkurl{.rds} a typed reusable result,
\nolinkurl{.svg} a vector graphic, and \nolinkurl{.md} a plain-text report. An \emph{output
inventory} is the list of required reported outputs. The \emph{current output inventory} applies the
current configuration. Construct this exact inventory:

\begin{itemize}
\item In \nolinkurl{tables/descriptive_statistics}, create
\nolinkurl{summary_stats.tex}, \nolinkurl{correlations.tex}, and \nolinkurl{ols_mean_eq.tex}.

\item In \nolinkurl{tables}, create each named typesetting source, a second source obtained by adding
\nolinkurl{_standalone} before the extension, and a rendered document with the same
\nolinkurl{_standalone} basename and extension \nolinkurl{.pdf}. Apply this rule to
\nolinkurl{hetero_tests.tex}, \nolinkurl{structural_var_inference.tex},
\nolinkurl{structural_var_estimators.tex}, \nolinkurl{var_share.tex}, and
\nolinkurl{variance_bounds_summary.tex}. These five triples and the preceding three files give 18
table files.

\item In \nolinkurl{reports}, create \nolinkurl{descriptive_stats.tex},
\nolinkurl{descriptive_stats.pdf}, and \nolinkurl{approximation_error_quoted_numbers.md}.

\item In \nolinkurl{diagnostics}, create \nolinkurl{set_id_inference_diagnostics.csv},
\nolinkurl{set_id_spec_comparison.csv}, \nolinkurl{log_var_eq_set_inference_diagnostics.csv},
\nolinkurl{log_var_eq_joint_null.csv}, \nolinkurl{log_var_eq_joint_null.rds},
\nolinkurl{log_var_eq_joint_gmm.csv}, \nolinkurl{log_var_eq_joint_gmm.rds},
\nolinkurl{approximation_error_quoted_numbers.csv}, and
\nolinkurl{log_var_eq_lad_closure.csv}.

\item In \nolinkurl{state}, always create \nolinkurl{bootstrap_stage_draws.rds},
\nolinkurl{log_var_eq_dynamics_gate.rds}, \nolinkurl{log_var_eq_egarch_status.rds}, and
\nolinkurl{conditional_route_status.rds}.

\item In \nolinkurl{figures}, create \nolinkurl{set_id_bounds_tau.svg},
\nolinkurl{set_id_projections_sd.svg}, \nolinkurl{set_id_projections_sd_ols.svg},
\nolinkurl{log_var_eq_bounds_tau_logols.svg}, \nolinkurl{log_var_eq_bounds_tau_ppml.svg},
\nolinkurl{log_var_eq_bounds_tau_harvey.svg}, \nolinkurl{log_var_eq_fitted_volatility_ppml.svg},
\nolinkurl{log_var_eq_fitted_volatility_harvey.svg}, \nolinkurl{variance_bounds.svg},
\nolinkurl{log_var_eq_bounds_tau_lad.svg}, and \nolinkurl{log_var_eq_fitted_volatility_lad.svg}.
Also create
\nolinkurl{figures/descriptive_statistics/figures.pdf} in the indicated subfolder.

\item For each estimator label \nolinkurl{ppml} and \nolinkurl{harvey}, create nine additional
files in \nolinkurl{figures}. Their common form is
\nolinkurl{log_var_eq_fitted_volatility_[estimator]_[suffix].svg}. Replace
\nolinkurl{[estimator]} with the estimator label and \nolinkurl{[suffix]} successively with
\nolinkurl{tau0p05}, \nolinkurl{tau0p1}, \nolinkurl{tau0p2}, \nolinkurl{tau0p3},
\nolinkurl{tau0p4}, \nolinkurl{tau_combined}, \nolinkurl{tau_combined_log},
\nolinkurl{tau_combined_lower_normalized}, and \nolinkurl{tau_combined_upper_normalized}. This rule
creates 18 sweep figures of fitted conditional standard deviations.

\item For each suffix \nolinkurl{tau0p05}, \nolinkurl{tau0p1}, \nolinkurl{tau0p2}, and
\nolinkurl{tau0p3}, create four additional files in \nolinkurl{figures}:
\nolinkurl{set_id_region_3d_sd_[suffix].svg}, \nolinkurl{set_id_region_3d_b_[suffix].svg},
\nolinkurl{set_id_region_3d_ols_projected_sd_[suffix].svg}, and
\nolinkurl{set_id_region_3d_ols_projected_b_[suffix].svg}. Replace \nolinkurl{[suffix]} with the
current suffix. This rule creates 16 three-dimensional region figures.
\end{itemize}

The current output inventory contains 18 table files, three report files, nine diagnostic files,
four state files, and 46 figure files, for a total of 80 files.
\end{sloppypar}

\subsection{Machine-readable output schemas}

An \emph{output schema} is the ordered list of fields in a machine-readable output, together with
each field's mathematical type. A \emph{field} is one named entry. An \emph{outer field} is a field
at the highest level. A \emph{record} is an ordered group of named fields. A \emph{collection record} is one
record within a collection, and a \emph{collection} is an ordered sequence of records with the
same fields. In the lists below, C denotes a character string, meaning a sequence of text
characters; D denotes a binary64 real number; I denotes an integer; L denotes a logical
value that is either true or false; N denotes a present field with no value; Q denotes a quarter
label; R denotes a record; K denotes a
collection; and T denotes a typed table. A superscript gives a vector length or matrix dimensions. A star in a superscript
denotes a finite length that the surrounding rule determines.

Store
comma-separated tables in UTF-8, a standard text-character encoding. Print a finite real number with
17 significant digits, an undefined real value as the literal label \nolinkurl{NaN}, an infinite
value as the literal label \nolinkurl{Inf} or \nolinkurl{-Inf}, a logical value as the literal label
\nolinkurl{TRUE} or \nolinkurl{FALSE}, and a missing value as the literal label \nolinkurl{NA}.
Enclose text values in quotation marks and leave every other value unquoted. A text value that
equals the missing-value label, or that already begins with the \emph{escape character}, a single
backslash, carries one extra leading escape character in storage, which retrieval removes. Preserve
the stated column order.

Store each
typed reusable result in version 3 of the binary format carried by the reusable-result extension
named in the output legend, a public format for ordered typed values. An exact
retrieval match requires identical field names, field order, mathematical types, dimensions,
values, and missing-value locations.

The literal field names below are fixed output identifiers. They do not create alternate terms for
the mathematical concepts defined in this manual.

Two of the machine-readable records carry fields belonging to moment blocks that the current
configuration disables. Those fields are written so the record's shape does not depend on which
blocks ran, and every one of them that carries a value is an inert control value carried from the single joint-moment
specification, a count fixed by the model's dimensions, or a status or reason entry recording that
the block did not run. None enters any
calculation on the current route, and none is reproducible from Part~I, because Part~I describes only
what the route performs. Reproduce them by writing the record with the declared fields present:
record the numerical-status entry as skipped, the coverage-status entry as not applicable, and the
reason entry as the declined gate; leave the search-status entry at its missing value; copy in the
three numerical controls of the joint-moment specification, which are the same values the
always-written replication row carries, together with those moment and parameter counts the model's
dimensions fix for that block, which the combined log-and-exponential block carries and the
additional-instrument block leaves missing; and leave every remaining block field at its typed
missing value.
Reconstruct nothing further about a block that does not run.

\begin{sloppypar}
The table \nolinkurl{set_id_inference_diagnostics.csv} has 28 rows: seven coefficients at zero slack,
then seven coefficients at each of the three positive display slacks. Its ordered columns are
\nolinkurl{inference_version}[C], \nolinkurl{nominal_alpha}[D],
\nolinkurl{minimum_valid_draw_share}[D], \nolinkurl{coef}[C], \nolinkurl{tau}[D],
\nolinkurl{set_lower}[D], \nolinkurl{set_upper}[D], \nolinkurl{width}[D],
\nolinkurl{set_status}[C], \nolinkurl{set_lower_status}[C], \nolinkurl{set_upper_status}[C],
\nolinkurl{n_bounded_lower}[I], \nolinkurl{n_unbounded_lower}[I],
\nolinkurl{n_unreliable_lower}[I], \nolinkurl{n_failed_lower}[I],
\nolinkurl{n_bounded_upper}[I], \nolinkurl{n_unbounded_upper}[I],
\nolinkurl{n_unreliable_upper}[I], \nolinkurl{n_failed_upper}[I], \nolinkurl{min_reps}[D],
\nolinkurl{t_lower}[D], \nolinkurl{t_upper}[D], \nolinkurl{se_lower}[D],
\nolinkurl{se_upper}[D], \nolinkurl{n_lower}[I], \nolinkurl{n_upper}[I],
\nolinkurl{n_common}[I], \nolinkurl{n_non_failed_lower}[I],
\nolinkurl{n_non_failed_upper}[I], \nolinkurl{frac_lower}[D], \nolinkurl{frac_upper}[D],
\nolinkurl{gate_lower}[L], \nolinkurl{gate_upper}[L], \nolinkurl{side}[C], \nolinkurl{c_s}[D],
\nolinkurl{c_p_lower}[D], \nolinkurl{c_p_upper}[D], \nolinkurl{c_p_gap}[D],
\nolinkurl{c_p_evals}[I], \nolinkurl{c_p_lambda}[D], \nolinkurl{c_p_interior}[L],
\nolinkurl{ci_lower}[D], \nolinkurl{ci_upper}[D], \nolinkurl{reason}[C],
\nolinkurl{ci_normal_lower}[D], \nolinkurl{ci_normal_upper}[D],
\nolinkurl{ci_pct_lower}[D], \nolinkurl{ci_pct_upper}[D], \nolinkurl{ci_basic_lower}[D],
\nolinkurl{ci_basic_upper}[D], \nolinkurl{rho}[D], \nolinkurl{c_stoye}[D],
\nolinkurl{c_im}[D], \nolinkurl{point_n_bounded}[I], \nolinkurl{point_n_unbounded}[I],
\nolinkurl{point_n_unreliable}[I], \nolinkurl{point_n_failed}[I],
\nolinkurl{point_n_valid}[I], \nolinkurl{point_n_non_failed}[I],
\nolinkurl{point_frac_bounded}[D], and \nolinkurl{p_value}[D].

The table \nolinkurl{set_id_spec_comparison.csv} has 56 rows. Put
the 28 reported-specification rows first and the 28 comparison-specification rows second. Begin
with \nolinkurl{inference_version}[C], \nolinkurl{nominal_alpha}[D],
\nolinkurl{minimum_valid_draw_share}[D], \nolinkurl{spec}[C], \nolinkurl{published}[L], and
\nolinkurl{impose_null}[L]. Follow them with every column from \nolinkurl{coef} through
\nolinkurl{p_value} in the preceding schema, in the same order and with the same types.

The table \nolinkurl{log_var_eq_set_inference_diagnostics.csv} has 30 rows: two estimators, three
positive display slacks, and five coefficients. Its ordered columns are
\nolinkurl{inference_version}[C], \nolinkurl{nominal_alpha}[D],
\nolinkurl{minimum_valid_draw_share}[D], \nolinkurl{estimator}[C], \nolinkurl{tau}[D],
\nolinkurl{coef}[C], \nolinkurl{se_lower}[D], \nolinkurl{se_upper}[D], \nolinkurl{n_lower}[I],
\nolinkurl{n_upper}[I], \nolinkurl{n_common}[I], \nolinkurl{n_non_failed_lower}[I],
\nolinkurl{n_non_failed_upper}[I], \nolinkurl{frac_lower}[D], \nolinkurl{frac_upper}[D],
\nolinkurl{gate_lower}[L], \nolinkurl{gate_upper}[L], \nolinkurl{side}[C], \nolinkurl{c_s}[D],
\nolinkurl{c_p_lower}[D], \nolinkurl{c_p_upper}[D], \nolinkurl{c_p_gap}[D],
\nolinkurl{c_p_evals}[I], \nolinkurl{c_p_lambda}[D], \nolinkurl{c_p_interior}[L],
\nolinkurl{ci_lower}[D], \nolinkurl{ci_upper}[D], \nolinkurl{reason}[C],
\nolinkurl{ci_normal_lower}[D], \nolinkurl{ci_normal_upper}[D],
\nolinkurl{ci_pct_lower}[D], \nolinkurl{ci_pct_upper}[D], \nolinkurl{ci_basic_lower}[D],
\nolinkurl{ci_basic_upper}[D], \nolinkurl{c_sim}[D], \nolinkurl{anchor_lower}[D],
\nolinkurl{anchor_upper}[D], \nolinkurl{published_lower}[D], \nolinkurl{published_upper}[D],
\nolinkurl{published_status}[C], \nolinkurl{sens_ci_lower}[D], \nolinkurl{sens_ci_upper}[D],
\nolinkurl{point_estimate}[D], \nolinkurl{point_se}[D], \nolinkurl{point_statistic}[D],
\nolinkurl{point_p_value}[D], \nolinkurl{point_n_bounded}[I],
\nolinkurl{point_n_unbounded}[I], \nolinkurl{point_n_unreliable}[I],
\nolinkurl{point_n_failed}[I], \nolinkurl{point_n_valid}[I],
\nolinkurl{point_n_non_failed}[I], \nolinkurl{point_reason}[C], \nolinkurl{tau0_sd_boot}[D],
\nolinkurl{tau0_se_analytic}[D], \nolinkurl{tau0_ratio}[D], \nolinkurl{failed_count}[D], and
\nolinkurl{sens_failed_count}[D]. Here \nolinkurl{anchor_lower} and \nolinkurl{anchor_upper} are the
endpoints \(\widehat\phi_{\mathrm{lo}}\) and \(\widehat\phi_{\mathrm{hi}}\) of the unresampled per-draw-budget calculation; \nolinkurl{published_lower},
\nolinkurl{published_upper}, and \nolinkurl{published_status} are the full-budget conditional-scale identified endpoints and set state,
which the confidence interval does not use.

The table \nolinkurl{log_var_eq_joint_null.csv} has four rows, ordered by slack: \(0\), \(0.05\),
\(0.10\), and \(0.20\). Its ordered columns are \nolinkurl{schema_version}[C], \nolinkurl{estimator}[C],
\nolinkurl{diagnostic}[C], \nolinkurl{sample_id}[C], \nolinkurl{inference_status}[C],
\nolinkurl{tau_order}[I], \nolinkurl{tau_display}[D], \nolinkurl{tau_internal}[D],
\nolinkurl{scaled_l2}[D], \nolinkurl{scaled_linf}[D], \nolinkurl{b1}[D], \nolinkurl{b2}[D],
\nolinkurl{b3}[D], \nolinkurl{theta0}[D], \nolinkurl{thetaR1}[D], \nolinkurl{thetaR2}[D],
\nolinkurl{thetaR3}[D], \nolinkurl{thetaR4}[D], \nolinkurl{min_abs_eps}[D],
\nolinkurl{nearest_crossing_qtr}[Q or missing], \nolinkurl{constraint_residual}[D],
\nolinkurl{raw_gradient_norm}[D], \nolinkurl{grid_points}[I],
\nolinkurl{n_unavailable_grid}[I], \nolinkurl{polish_starts}[I],
\nolinkurl{successful_polishes}[I], \nolinkurl{argmin_source}[C],
\nolinkurl{winning_start_type}[C], \nolinkurl{n_agreeing_starts}[I],
\nolinkurl{start_agreement_max_scaled}[D], \nolinkurl{basin_count}[I],
\nolinkurl{replication_status}[C], \nolinkurl{perturbation_status}[C],
\nolinkurl{stability_status}[C], \nolinkurl{n_near_crossing_rows}[I], \nolinkurl{status}[C],
\nolinkurl{membership_result}[C], \nolinkurl{root_tol}[D], \nolinkurl{scale1}[D],
\nolinkurl{scale2}[D], \nolinkurl{scale3}[D], \nolinkurl{scale4}[D],
\nolinkurl{box_active}[L], \nolinkurl{sparse_grid}[L], \nolinkurl{metric_l2_source}[C],
\nolinkurl{metric_linf_source}[C], \nolinkurl{l2_rank_under_linf}[I],
\nolinkurl{linf_rank_under_l2}[I], \nolinkurl{d_inf_l2}[D],
\nolinkurl{d_inf_grid_best}[D], \nolinkurl{linf_gap_abs}[D], \nolinkurl{linf_gap_rel}[D],
\nolinkurl{linf_trigger}[L], and \nolinkurl{linf_sensitivity_status}[C]. These coefficient columns do
not follow the naming of this manual, so map them explicitly: \nolinkurl{b1}, \nolinkurl{b2}, and \nolinkurl{b3} are the coordinates of the structural SDF-news coefficient vector \(\theta\);
\nolinkurl{theta0} is the conditional-scale intercept \(\eta_1\); \nolinkurl{thetaR1} through
\nolinkurl{thetaR4} are the coordinates of the conditional-scale slope vector \(\eta_R\); and \nolinkurl{scale1} through
\nolinkurl{scale4} are the \emph{reciprocals} of the supplied return scores' standard deviations, that is, the diagonal entries of \(D_R\). In particular the columns named for \(\theta\) do not hold
\(\theta\), and the columns that hold \(\theta\) are named \nolinkurl{b}.

The table \nolinkurl{log_var_eq_joint_gmm.csv} has three rows:
\nolinkurl{graph_replication}, \nolinkurl{z}, and \nolinkurl{log_ppml}. The \nolinkurl{z} and
\nolinkurl{log_ppml} rows record skipped blocks. Its ordered columns are \nolinkurl{schema_version}[C],
\nolinkurl{diagnostic}[C], \nolinkurl{sample_id}[C], \nolinkurl{joint_input_id}[C],
\nolinkurl{decision_spec_id}[C], \nolinkurl{spec_id}[C], \nolinkurl{inference_status}[C],
\nolinkurl{moment_block}[C], \nolinkurl{tau_order}[I], \nolinkurl{tau_display}[D],
\nolinkurl{tau_internal}[D], \nolinkurl{numerical_status}[C], \nolinkurl{search_status}[C],
\nolinkurl{coverage_status}[C], \nolinkurl{rho_scaled_linf}[D], \nolinkurl{b1}[D],
\nolinkurl{b2}[D], \nolinkurl{b3}[D], \nolinkurl{theta0}[D], \nolinkurl{thetaR1}[D],
\nolinkurl{thetaR2}[D], \nolinkurl{thetaR3}[D], \nolinkurl{thetaR4}[D], \nolinkurl{a_l}[D],
\nolinkurl{a_p}[D], \nolinkurl{gap_a_p_a_l}[D], \nolinkurl{beta1}[D], \nolinkurl{beta2}[D],
\nolinkurl{beta3}[D], \nolinkurl{beta4}[D], \nolinkurl{fresh_moment_residual}[D],
\nolinkurl{fresh_constraint_residual}[D], \nolinkurl{crossing_count}[I],
\nolinkurl{min_abs_eps}[D], \nolinkurl{eps_floor}[D], \nolinkurl{guard_slack}[D],
\nolinkurl{floor_active}[L], \nolinkurl{floor_sensitivity_status}[C],
\nolinkurl{patterns_observed}[I], \nolinkurl{patterns_polished}[I],
\nolinkurl{argmin_source}[C], \nolinkurl{winning_start_type}[C],
\nolinkurl{optimizer_status}[C], \nolinkurl{box_half_width}[D], \nolinkurl{box_active}[L],
\nolinkurl{box_audit_status}[C], \nolinkurl{n_obs}[I], \nolinkurl{n_added_moments}[I],
\nolinkurl{n_search_parameters}[I], \nolinkurl{n_moments_unprofiled}[I],
\nolinkurl{n_parameters_unprofiled}[I], \nolinkurl{n_moments_profiled}[I],
\nolinkurl{n_parameters_profiled}[I], \nolinkurl{jac_rank_unprofiled}[I],
\nolinkurl{jac_rank_profiled}[I], \nolinkurl{jac_min_sv_unprofiled}[D],
\nolinkurl{jac_min_sv_profiled}[D], \nolinkurl{jac_status_unprofiled}[C],
\nolinkurl{jac_status_profiled}[C], \nolinkurl{active_constraint_count}[I],
\nolinkurl{active_constraint_rank}[I], \nolinkurl{moment_active_stacked_rank}[I],
\nolinkurl{algebraic_moment_excess}[I], \nolinkurl{param_xtol_rel}[D],
\nolinkurl{constraint_tol}[D], \nolinkurl{root_tol}[D], \nolinkurl{moment_delta}[D], and
\nolinkurl{skip_reason}[C]. Its coefficient columns need their own map, which differs from the table
above. Here \nolinkurl{b1} through \nolinkurl{b3} hold \(\theta\); \nolinkurl{a_l} and
\nolinkurl{a_p} hold the mean-log intercept \(a_L\) and the exponential-mean intercept \(a_P\)
defined with those extensions below, \(a_L\) being the conditional-scale intercept \(\eta_1\) of the
log-OLS conditional-scale estimator mapping when the cross-block restriction is not imposed; \nolinkurl{gap_a_p_a_l}
holds their difference; and \nolinkurl{beta1} through \nolinkurl{beta4} hold the conditional-scale slope vector
\(\eta_R\)---not \nolinkurl{thetaR1} through \nolinkurl{thetaR4}, which this table declares but never
populates, as it does not populate \nolinkurl{theta0}. All five stay at their missing value in every
row.

The table \nolinkurl{approximation_error_quoted_numbers.csv} has the columns \nolinkurl{name}[C] and
\nolinkurl{value}[D]. Its five rows, in order, are \nolinkurl{max_bound_sdf_news},
\nolinkurl{max_bound_expected_sdf}, \nolinkurl{max_error_sd_bp},
\nolinkurl{typical_quarterly_move_one_quarter_yield}, and
\nolinkurl{max_error_sd_over_typical_move_pct}.

The table \nolinkurl{log_var_eq_lad_closure.csv} may have no data rows, but must retain this ordered header: \nolinkurl{coef}[C], \nolinkurl{tau}[D], \nolinkurl{estimator}[C],
\nolinkurl{sample_id}[C], \nolinkurl{spec_id}[C], \nolinkurl{crossing_row}[I],
\nolinkurl{crossing_qtr}[Q or missing], \nolinkurl{residual_side}[C], \nolinkurl{coef_side}[C],
\nolinkurl{limit_value}[D], \nolinkurl{path_id}[I], \nolinkurl{witness_status}[C],
\nolinkurl{M_span}[D], \nolinkurl{tail_slope}[D], \nolinkurl{tail_slope_stability}[D],
\nolinkurl{path_summary}[C], and \nolinkurl{provenance}[C]. Sort its rows by slack, crossing-row
number, coefficient, and path number, in that order.

The reusable result \nolinkurl{bootstrap_stage_draws.rds} has these ordered outer fields:
\nolinkurl{anchor}[K], \nolinkurl{mean}[R], \nolinkurl{volatility_primary}[K],
\nolinkurl{volatility_primary_n_failed}[I], \nolinkurl{volatility_sensitivity}[K],
\nolinkurl{volatility_sensitivity_n_failed}[I], and \nolinkurl{provenance}[R]. The
\nolinkurl{anchor} collection is ordered by PPML and Harvey QML. Each estimator record contains collection records
for slacks \(0,0.05,0.10,0.20\), in that order. Every collection record contains
\nolinkurl{lower}[\(D^5\)], \nolinkurl{upper}[\(D^5\)],
\nolinkurl{lower_status}[\(C^5\)], and \nolinkurl{upper_status}[\(C^5\)]. The zero-slack collection record also
contains \nolinkurl{point}[\(D^5\)] and \nolinkurl{point_status}[\(C^5\)].

The ordered fields of \nolinkurl{mean}[R] are \nolinkurl{point_draws}[\(D^{B\times7}\)],
\nolinkurl{point_status}[\(C^{B\times7}\)], \nolinkurl{n_point_deficient}[I],
\nolinkurl{endpoint_draws}[K], \nolinkurl{n_failed}[I], and \nolinkurl{failure_causes}[T or N]. The
\nolinkurl{endpoint_draws} collection contains three positive-slack collection records. Each collection record
contains \nolinkurl{lower}[\(D^{B\times7}\)], \nolinkurl{upper}[\(D^{B\times7}\)],
\nolinkurl{lower_status}[\(C^{B\times7}\)], and
\nolinkurl{upper_status}[\(C^{B\times7}\)]. When it is not N,
\nolinkurl{failure_causes} is a one-dimensional table of integer counts with character names. The primary and sensitivity
collections are each ordered by PPML and Harvey QML and then by the same four slacks. In the primary
collection, \nolinkurl{lower} and \nolinkurl{upper} each have type
\(D^{B\times5}\); \nolinkurl{lower_status} and \nolinkurl{upper_status} each have type
\(C^{B\times5}\); \nolinkurl{point} has type \(D^{B\times5}\); and
\nolinkurl{point_status} has type \(C^{B\times5}\). The sensitivity matrices have the corresponding types with \(B_s\) in
place of \(B\), where \(B_s\) is the positive integer sensitivity bootstrap-draw count. The current
configuration draws both families the same number of times, so \(B_s=B=10{,}000\).

The ordered fields of \nolinkurl{provenance}[R] are \nolinkurl{resampler}[C],
\nolinkurl{sample_size}[I], \nolinkurl{b_reps}[I], \nolinkurl{block}[I], \nolinkurl{seed}[I],
\nolinkurl{rng_kind}[\(C^3\)], \nolinkurl{block_rule}[C], an identity entry for the primary
bootstrap-index family, an identity entry for the random-number state left after that family is
drawn, \nolinkurl{sens_block}[I], \nolinkurl{sens_reps}[I], the two corresponding identity entries
for the sensitivity bootstrap-index family, \nolinkurl{axes}[R], an identity entry for the resampled
data table, an identity entry for the draw specification with that table removed, an identity entry
for the calculation rules that produce the draws, an identity entry for the presentation rules, an
identity entry for the numerical environment, and \nolinkurl{cache_schema_version}[I]: twenty fields
in that order, recording the primary and sensitivity bootstrap-index families separately. The nine
identity entries are character strings, and none is reproducible from Part~I, because each fixes one
particular stored representation, calculation text, or computing environment rather than a
mathematical quantity. Write each from the reproduction's own corresponding object. The ordered
\nolinkurl{axes}[R] fields are \nolinkurl{mean_coefs}[\(C^7\)],
\nolinkurl{mean_taus}[\(D^3\)], \nolinkurl{volatility_coefs}[\(C^5\)],
\nolinkurl{volatility_taus}[\(D^4\)], and \nolinkurl{volatility_estimators}[\(C^2\)].

The reusable result \nolinkurl{log_var_eq_dynamics_gate.rds} has these ordered outer fields:
\nolinkurl{schema_version}[C], \nolinkurl{protocol}[R], \nolinkurl{sample_id}[C], a
character-string version identity, \nolinkurl{b_point}[\(D^3\)], \nolinkurl{n}[I],
\nolinkurl{tested_lags}[\(I^3\)], \nolinkurl{q_stats}[\(D^3\)], \nolinkurl{p_values}[\(D^3\)],
\nolinkurl{gate_lag}[I], \nolinkurl{gate_alpha}[D], \nolinkurl{verdict}[C],
\nolinkurl{min_abs_eps}[D], \nolinkurl{crossing_status}[C],
\nolinkurl{crossing_qtr}[\(Q^*\) or N], \nolinkurl{max_abs_xi}[D],
\nolinkurl{max_abs_xi_qtr}[Q], \nolinkurl{acf}[\(D^8\)], \nolinkurl{xi_hat}[T],
\nolinkurl{sensitivity}[R], and \nolinkurl{notes}[\(C^3\)]. The ordered
\nolinkurl{protocol}[R] fields are \nolinkurl{version}[C], \nolinkurl{tested_lags}[\(I^3\)],
\nolinkurl{gate_lag}[I], \nolinkurl{alpha}[D], \nolinkurl{acf_max}[I],
\nolinkurl{tie_tol}[D], and \nolinkurl{verdicts}[\(C^3\)]. The tested lags are 1, 4, and 8. The
\nolinkurl{q_stats} and \nolinkurl{p_values} vectors use the names \nolinkurl{lag1},
\nolinkurl{lag4}, and \nolinkurl{lag8}; \nolinkurl{acf} uses \nolinkurl{lag1} through
\nolinkurl{lag8}. The \nolinkurl{xi_hat} table has \nolinkurl{qtr}[Q] and \nolinkurl{xi_hat}[D];
its entries hold the zero-slack log-OLS residual \(u_t^L\), not the linear predictor \(\xi_t\) of
Section~\ref{sec:expcondvar}.
The ordered \nolinkurl{sensitivity}[R] fields are \nolinkurl{rows}[T] and
\nolinkurl{exclusions}[T]. The rows table has \nolinkurl{source}[C], \nolinkurl{p_lag1}[D],
\nolinkurl{p_lag4}[D], \nolinkurl{p_lag8}[D], and \nolinkurl{b_n}[\(D^3\)]. The exclusions table has
\nolinkurl{source}[C] and \nolinkurl{reason}[C].

The reusable result \nolinkurl{log_var_eq_egarch_status.rds} has these ordered fields:
\nolinkurl{schema_version}[C], \nolinkurl{protocol}[R], \nolinkurl{stage}[C],
\nolinkurl{plan}[C], \nolinkurl{sample_id}[C], a character-string version identity,
\nolinkurl{gate_verdict}[C], \nolinkurl{gate_lag}[I], \nolinkurl{gate_alpha}[D],
\nolinkurl{gate_q}[D], \nolinkurl{gate_p}[D], \nolinkurl{min_abs_eps}[D],
\nolinkurl{crossing_status}[C], \nolinkurl{estimand_decision}[C],
\nolinkurl{dependency_decision}[C], \nolinkurl{decision_provenance}[C],
\nolinkurl{decided_at_utc}[C], \nolinkurl{run_dynamic}[L], \nolinkurl{terminal_status}[C],
\nolinkurl{skip_reason}[C], \nolinkurl{decision_pending}[L], \nolinkurl{workstream_status}[C],
\nolinkurl{dynamic_artifacts_absent}[L], \nolinkurl{cleanup_audit}[R], and
\nolinkurl{gate_record_path}[C]. Its \nolinkurl{protocol}[R] record is the residual-dynamics check
defined immediately above. The \nolinkurl{cleanup_audit}[R] record has
\nolinkurl{status}[C], \nolinkurl{artifacts}[T], \nolinkurl{n_existed}[I],
\nolinkurl{n_deleted}[I], and \nolinkurl{all_absent}[L], in that order. Its artifacts table has
\nolinkurl{id}[C], \nolinkurl{path}[C], \nolinkurl{existed}[L], and \nolinkurl{deleted}[L].

The \emph{route-status result} records which conditional calculations the current configuration
requests and completes. Its reusable result \nolinkurl{conditional_route_status.rds} has the ordered outer fields
\nolinkurl{schema_version}[C], \nolinkurl{routes}[K], \nolinkurl{cleanup}[K], and
\nolinkurl{artifacts}[T]. Order both \nolinkurl{routes} and \nolinkurl{cleanup} as
\nolinkurl{conditional_lad}, then \nolinkurl{conditional_egarch}. Each route-status record has
\nolinkurl{decision}[C], \nolinkurl{requested}[L], and \nolinkurl{ran}[L]. Each cleanup record has
\nolinkurl{status}[C], \nolinkurl{artifacts}[T], \nolinkurl{n_existed}[I],
\nolinkurl{n_deleted}[I], and \nolinkurl{all_absent}[L]. Its artifacts table has
\nolinkurl{id}[C], \nolinkurl{path}[C], \nolinkurl{existed}[L], and \nolinkurl{deleted}[L]. The
outer artifacts table has \nolinkurl{id}[C], \nolinkurl{status}[C], \nolinkurl{path}[C], and
\nolinkurl{present}[L].
Under the current configuration, the LAD route-status record states that LAD is approved,
requested, and completed. The dynamic-volatility route-status record records the non-reject
residual-dynamics outcome and states that a dynamic-volatility calculation is neither
requested nor completed. The three LAD output locations must be present after successful
reproduction, and the four dynamic-volatility output locations must remain empty.

The reusable result \nolinkurl{log_var_eq_joint_null.rds} has the ordered outer fields
\nolinkurl{schema_version}[C], \nolinkurl{estimator}[C], \nolinkurl{diagnostic}[C],
\nolinkurl{inference_status}[C], \nolinkurl{sample_id}[C], a character-string version identity,
\nolinkurl{root_tol}[D], \nolinkurl{protocol}[R], \nolinkurl{scales}[R], and
\nolinkurl{rows}[K]. The ordered \nolinkurl{protocol}[R] fields are
\nolinkurl{version}[C], \nolinkurl{objective_chunk_size}[I],
\nolinkurl{objective_machine_epsilon_multiplier}[D],
\nolinkurl{objective_root_tolerance_multiplier}[D], \nolinkurl{grid_n}[I],
\nolinkurl{grid_floor}[I], \nolinkurl{root_tol}[D], \nolinkurl{feasibility_tol}[D],
\nolinkurl{crossing_rel_tol}[D], \nolinkurl{machine_adjacent_rel_tol}[D],
\nolinkurl{coordinate_agreement_rel_tol}[D], \nolinkurl{fresh_map_rel_tol}[D],
\nolinkurl{box_active_rel_tol}[D], \nolinkurl{required_agreeing_starts}[I],
\nolinkurl{separated_starts}[I], \nolinkurl{separated_start_fraction}[D],
\nolinkurl{perturbation_fractions}[\(D^3\)], \nolinkurl{perturbation_rank_tol}[D],
\nolinkurl{perturbation_dedupe_tol}[D], \nolinkurl{perturbation_sign_tol}[D],
\nolinkurl{puncture_ratio}[D], \nolinkurl{off_manifold_root_multiplier}[D],
\nolinkurl{linf_gap_root_multiplier}[D], \nolinkurl{linf_gap_rel_tol}[D],
\nolinkurl{epigraph_xtol_rel}[D], \nolinkurl{epigraph_maxeval}[I],
\nolinkurl{epigraph_bound}[D], and \nolinkurl{epigraph_constraint_tol}[D]. The
\nolinkurl{scales}[R] record has \nolinkurl{d}[\(D^4\)],
\nolinkurl{d_inv2}[\(D^4\)], and \nolinkurl{names}[\(C^4\)], in that order. The rows collection contains
four ordered row records, each with the fields and types of \nolinkurl{log_var_eq_joint_null.csv}.

The reusable result \nolinkurl{log_var_eq_joint_gmm.rds} has the ordered outer fields
\nolinkurl{schema_version}[C], \nolinkurl{diagnostic}[C], \nolinkurl{inference_status}[C],
\nolinkurl{sample_id}[C], \nolinkurl{joint_input_id}[C], \nolinkurl{decision_spec_id}[C],
\nolinkurl{spec_id}[C], \nolinkurl{decision}[R], \nolinkurl{spec}[R],
\nolinkurl{graph_replication}[R], \nolinkurl{projection}[R], and \nolinkurl{rows}[K]. The ordered
\nolinkurl{decision}[R] fields are \nolinkurl{schema_version}[C],
\nolinkurl{decision_spec_id}[C], \nolinkurl{enable_z}[L], \nolinkurl{enable_log_ppml}[L],
\nolinkurl{intercept_target}[C], \nolinkurl{stage_c_requested}[L],
\nolinkurl{moment_delta}[\(D^0\)], \nolinkurl{gaussian_gap_restriction}[L],
\nolinkurl{decided_on}[C], \nolinkurl{decision_provenance}[C], and
\nolinkurl{rationale}[\(C^6\)].
For the current record, \nolinkurl{enable_z}, \nolinkurl{enable_log_ppml},
\nolinkurl{stage_c_requested}, and \nolinkurl{gaussian_gap_restriction} are false;
\nolinkurl{intercept_target} is \nolinkurl{a_L}; and \nolinkurl{moment_delta} is empty.
The ordered \nolinkurl{spec}[R] fields are \nolinkurl{sample_id}[C],
\nolinkurl{joint_input_id}[C], \nolinkurl{decision_spec_id}[C], \nolinkurl{basis_version}[C],
\nolinkurl{profile_version}[C], \nolinkurl{candidate_version}[C],
\nolinkurl{enabled_blocks}[\(C^0\)], \nolinkurl{moment_scales}[\(D^0\)],
\nolinkurl{coord_scales}[\(D^0\)], \nolinkurl{execution_constants}[R],
\nolinkurl{controls}[\(D^3\)], \nolinkurl{budgets}[\(D^5\)],
\nolinkurl{moment_delta}[\(D^0\)], \nolinkurl{combine_blocks}[L],
\nolinkurl{stage_c_enabled}[L], and \nolinkurl{spec_id}[C]. The
ordered \nolinkurl{execution_constants}[R] fields are \nolinkurl{param_xtol_rel}[D],
\nolinkurl{constraint_tol}[D], \nolinkurl{root_tol}[D],
\nolinkurl{box_half_widths}[\(D^3\)], \nolinkurl{pattern_start_cap}[I],
\nolinkurl{per_solve_maxeval}[I], \nolinkurl{candidate_extra_start_cap}[I],
\nolinkurl{candidate_min_separation}[D], \nolinkurl{replication_grid_n}[I],
\nolinkurl{replication_point_cap}[I], \nolinkurl{basis_axis_tol}[D],
\nolinkurl{basis_rank_tol}[D], \nolinkurl{basis_retained_gap_factor}[D],
\nolinkurl{basis_discarded_gap_factor}[D], \nolinkurl{jacobian_rank_tol}[D],
\nolinkurl{jacobian_weak_gap_factor}[D], \nolinkurl{design_rank_tol}[D],
\nolinkurl{epigraph_r_upper}[D], \nolinkurl{box_interior_tol}[D],
\nolinkurl{epigraph_start_abs_slack}[D], \nolinkurl{epigraph_start_rel_slack}[D],
\nolinkurl{nested_floor_tol}[D], \nolinkurl{floor_audit_rtol}[D],
\nolinkurl{projection_equality_rtol}[D], \nolinkurl{n_moments_unprofiled}[I],
\nolinkurl{n_parameters_unprofiled}[I], \nolinkurl{n_moments_profiled}[I], and
\nolinkurl{n_parameters_profiled}[I]. The \nolinkurl{controls} vector has, in order,
\nolinkurl{param_xtol_rel}, \nolinkurl{constraint_tol}, and \nolinkurl{root_tol}. The
\nolinkurl{budgets} vector has, in order, \nolinkurl{pattern_start_cap},
\nolinkurl{per_solve_maxeval}, \nolinkurl{box_h1}, \nolinkurl{box_h2}, and \nolinkurl{box_h3}. The
\nolinkurl{graph_replication}[R] record has \nolinkurl{replication_status}[C],
\nolinkurl{max_moment_residual}[D], \nolinkurl{n_points}[I], \nolinkurl{n_finite}[I], and
\nolinkurl{min_abs_eps}[D]. The disabled \nolinkurl{projection}[R] record has
\nolinkurl{status}[C], \nolinkurl{reason}[C], and \nolinkurl{rows}[K], with no rows. The outer rows
collection contains, in order, the \nolinkurl{graph_replication}, \nolinkurl{z}, and
\nolinkurl{log_ppml} row records with the fields and types of
\nolinkurl{log_var_eq_joint_gmm.csv}.

\end{sloppypar}

Print mean-equation, log-OLS, PPML, and Harvey QML coefficients to three decimal places. Print LAD
coefficients, test statistics, and variance shares to two decimal places. Print plug-in
approximation-error variance bounds in scientific notation with two digits after the mantissa's
decimal point, that is, three significant digits. Scientific
notation writes a number as a leading number times a power of ten. Use one star below a p-value of
0.10, two below 0.05, and three below 0.01. The \emph{unavailable-value symbol} is an en
dash. A mean-equation coefficient cell with an unbounded endpoint state is labeled \emph{unbounded}
and does not display a half-infinite identified range. A conditional-scale cell may show a
half-infinite identified range and, when defined, its associated half-infinite confidence interval.
Leave a finite identified range blank when the presentation calls for unequal endpoints and its
endpoints are equal. In the mean-equation, conditional-scale, and LAD cells this requires exact
equality. In the variance-share table, treat the endpoints as equal when
\(|\,\text{upper}-\text{lower}\,|\le10^{-9}\,(1+|\,\text{upper}\,|)\).

Print the unavailable-value symbol in the statistic row when zero-slack point inference has too few
retained bootstrap draws. Do the same when fewer than 85 percent of the relevant bootstrap draws
have the bounded endpoint state or when a finite coefficient has a nonfinite or nonpositive endpoint
MAD scale. Reserve a blank statistic row for a column for which no statistic is defined.

In the variance-share table, apply the following rule only in the two columns that display a single
coefficient vector, namely the mean least-squares benchmark column and the zero-slack point column,
and within those columns only to a block of individual PC-score contributions whose displayed
entries and block total are all finite. Leave the positive-slack identified-range columns
unrounded by this rule. Within such a block apply the largest-remainder display adjustment exactly as
Section~\ref{sec:shares} defines it. That adjustment makes the displayed individual
contributions sum to the displayed variance-share total without changing the unrounded values. Do not force the joint
variance share to equal the sum of the lagged expected-SDF and SDF-news variance shares.
Equation~\eqref{eq:share-combined} includes a cross term.

\section{Completion and preservation}

Each machine-readable table must survive a \emph{typed table round trip}, meaning that storage,
retrieval, and comparison retain each column's declared mathematical type. Quarter labels remain
quarter labels, integers remain integers, and finite real numbers agree with their stored values to
within relative tolerance \(10^{-9}\). Every column that is not real-valued must return identically,
and missing, infinite, and undefined values remain explicit.

Published tables require successful standalone compilation. Each reusable result must
survive an \emph{exact reusable-result round trip}, which preserves its field order, dimensions,
categorical values, and recorded origin exactly.
A categorical value is one of a fixed set of labels.

At the start of the current reproduction, clear the seven reserved conditional-output locations and record whether any prior
file was removed. At completion, require all three LAD outputs, require all four dynamic-volatility
locations to remain empty, and write the route-status result. Remove auxiliary typesetting files.

\section{Reproduction sequence}
\label{sec:sequence}

This order respects every scientific dependency.

\begin{enumerate}
\item Fix the input identities, mathematical settings, numerical controls, and random-number rules
stated in this part.
\item Verify every preserved external input before creating or clearing anything: confirm that the
preserved consumption copy is present, and confirm that the preserved daily Treasury release is the
one this part names, obtaining that release if no copy is preserved yet. Stop here if the
consumption copy is absent, or if a preserved daily Treasury copy is not the named release, so that a
failed check leaves the existing outputs untouched rather than stopping the reproduction after they have already been
cleared.
\item Create the output folders. Clear the three LAD and four dynamic-volatility output locations,
and retain the cleanup record for each location.
\item Load the quarterly Treasury yields and term premia. Construct \(n_{i,t}\), log bond-price
news, the SDF-news series, and the expected-SDF series with its unpaired mean correction.
\item Construct consumption growth from the quarterly consumption input.
\item Load the preserved daily Treasury input and construct realized yield volatility and the
centered five-year instrument.
\item Align the supplied return scores to quarters.
\item Form the three SDF panel matrices, perform their three PC constructions, and apply the lagged
expected-SDF PC score sign alignment.
\item Fit the preliminary mean regression and the mean least-squares benchmark. Form the
identification sample, estimate both reduced-form specifications, and construct the initial mean
identified set.
\item Calculate variance shares and their attained numerical ranges, then render the variance-share
table.
\item Form the conditional-scale sample and calculate the log-OLS conditional-scale estimator
mapping.
\item Calculate the slack profiles of the mean identified set and critical slack, and render the
figure of mean-equation identified ranges as a function of slack.
\item Calculate the PPML and Harvey QML conditional-scale estimator mappings and their analytic
covariances. Run the PPML coverage audit, Harvey sensitivity audit, and cold-start endpoint
replications.
\item Evaluate the joint-null compatibility check and stacked-moment replication check.
\item Evaluate the residual-dynamics check. Record the non-reject residual-dynamics outcome, close
the residual-dynamics check, and record that no dynamic-volatility calculation is requested.
\item Confirm the approved LAD decision and required numerical reference. Calculate the LAD
conditional-scale estimator mapping, Frisch--Newton nonuniqueness check, and independent endpoint
check.
\item Generate the primary and sensitivity bootstrap-index families. Confirm that no reusable
bootstrap result exists, calculate all 10,000 bootstrap draws in each family, enforce the failure and acceptance
rules, and preserve the validated reusable bootstrap result.
\item Calculate the comparison-specification bootstrap endpoints from the primary bootstrap-index
family.
\item Render the combined inference table, the estimator table pages, and all conditional-scale
identified-range and fitted conditional-scale figures.
\item Construct the mean-identified-set projection geometry and render the two-coordinate and
three-dimensional mean-identified-set figures.
\item Calculate the mean-equation identification tests and checks and render their table.
\item Calculate both plug-in approximation-error variance bounds, apply the quoted-number
acceptance rules, and render their figure and summary table.
\item Construct the descriptive report, including its tables and figures.
\item Verify the required conditional outputs, then write the route-status result, and remove auxiliary
typesetting files. Each table and reusable result undergoes its stated round trip when written, and
each published table is compiled as a standalone document when rendered.
\end{enumerate}

\part{Permitted modifications and declared limits}

\section{Scope of the modification catalogue}

This part exhausts the modifications represented by input choices,
recorded decisions, and named mathematical, numerical, reporting, preservation, or cleanup controls. The
\emph{complete reproduction route} is the full ordered calculation from inputs through accepted
outputs. A \emph{reusable-result identity} is the recorded collection of inputs, dimensions,
procedures, numerical controls, and random-number settings that determines whether a preserved
result may be reused. This part does not cover redesigning the complete reproduction route. The distinction matters because a
recorded field can describe a calculation that the complete reproduction route cannot perform.

A \emph{full-route modification} changes the current configuration while leaving the complete
reproduction route executable. A \emph{component modification} changes a self-contained calculation, but requires coordinated changes elsewhere before the complete reproduction route can
use it. A \emph{routing-only choice} changes the route-status result, but cannot produce the requested calculation. An \emph{accepted-but-refused choice} satisfies the stated field and value
checks, but the complete reproduction route stops before carrying it out. An \emph{invalid choice}
fails those checks. A \emph{declared unavailable extension} has a recorded decision or mathematical
description but no complete reproduction route. A modification is \emph{executable} when the
reproduction can carry out the changed calculation as described, whether or not the complete
reproduction route reaches it; the qualifier separates such a modification from a refused or
invalid choice and from a declared unavailable extension.
Only a full-route modification is an alternative reproduction. Component modifications,
routing-only choices, accepted-but-refused choices, invalid choices, and declared unavailable
extensions are limits, not alternatives to Part I.

Numerical safeguards remain part of a selected procedure. Endpoint-state rules, rank checks,
zero-residual checks, fallback starting values, and convergence checks therefore remain in Part I.
They respond to calculated values and are not selections made before the reproduction.

\section{Input and calendar modifications}

\subsection{Treasury inputs}

The Treasury input procedure permits the following modifications, subject to the stated restrictions.
The first two items state their own class; the remaining items are classified together after the
list.

\begin{enumerate}
\item The monthly Treasury release is chosen by rule: the preserved user copy when that copy exists,
otherwise the bundled copy. Loading validates the column layout and the date column and verifies no
recorded release identity, because the bundled copy is always available and the route therefore never
reaches the acquisition procedure that would perform that verification. Automatic acquisition is on or off. Both values are inert on this route because the resolved copy
always exists, so neither value reaches the acquisition procedure; the setting is an executable
full-route modification that changes nothing while the bundled copy resolves. That procedure also accepts a forced replacement that
obtains the latest monthly release, verifies it against its recorded identity, and may either show
or suppress progress messages; the complete reproduction route never calls it, so forced replacement
is an executable component modification. Reaching it through a coordinated change would replace the
monthly input with a later vintage, changing every yield, term premium, and quantity built on them,
and invalidating every preserved resampling result and every quoted sample span.
\item The daily Treasury input has exactly two origins, and any third value is refused before
calculation. The preserved origin uses the copy already held: a copy whose recorded identity matches
the release Part I names is accepted, a copy whose identity does not match is refused and stops the
reproduction rather than being replaced, and an absent copy is obtained at that recorded release,
verified against the recorded identity, and discarded if verification fails. The retrieving origin
accepts the latest published release instead. Daily observations have no bundled copy, and the New
York Fed provider cannot supply them. Selecting the retrieving origin changes the daily vintage and
therefore realized yield volatility, the instrument, both samples, every estimate, and every interval,
and it invalidates every preserved resampling result. It is an executable full-route modification. The
recorded release the preserved origin verifies against is itself a setting: naming a different release,
together with its identity, adopts that vintage while staying on the preserved origin, and is an
executable full-route modification with the same downstream consequences.
\item A quarterly Treasury panel may retain or discard an incomplete terminal quarter. Retention
uses the latest available monthly observation and labels it with the terminal month of its calendar
quarter. Discarding removes the incomplete quarter. Because the approximation sample is the full
quarterly Treasury panel including that quarter, discarding it changes every approximation-error
variance bound and therefore invalidates the quoted-number acceptance rules.
\item The maturity set may contain distinct integer maturities from 1 through 120 months. The full
SDF construction requires every yield and term-premium maturity used by its news horizon and every
maturity used by the approximation-error calculations.
\item The analysis may change the news horizon and the reported maturity spacing. The
largest requested maturity plus the news horizon cannot exceed 120 months.
\end{enumerate}

Retaining or discarding the incomplete terminal quarter, changing the maturity set, changing the
news horizon, and changing the reported maturity spacing are each executable full-route modifications.
Changing the maturity set or the news horizon changes the SDF panel matrices and therefore every
principal component, every estimate, and every interval built on them. Changing the reported maturity
spacing leaves the panel matrices untouched, because they use every integer maturity in the main
maturity range; it changes the maturities at which the two bound series are evaluated, and therefore the
maturities the approximation-error figure plots, the summary statistics of those series, and the
quoted-number acceptance comparisons built on them, while leaving every estimate, identified range, and interval unchanged. Changing the maturity set or the news horizon invalidates every preserved resampling result.
Changing the reported maturity spacing does not, and neither does retaining or discarding the
incomplete terminal quarter, which falls outside the principal-component window.

The Treasury extraction calculation accepts a first date, a last date, or both. These date filters are component
modifications because the complete reproduction route does not pass date filters to that
calculation. Any coordinated use must preserve enough observations for every lag and sample.

The Treasury input procedure also admits three input families: yields, term premia, and
\emph{risk-neutral yields}. A risk-neutral yield reflects expected future short rates without the
term-premium component. The complete reproduction route requires quarterly yields and term premia
for the SDF construction and daily yields for realized yield volatility. Adding risk-neutral yields
while retaining those required families is a harmless full-route input modification: the added
columns pass through the quarterly input, but do not enter the current calculations. Omitting or
substituting for either required quarterly family is invalid for the complete reproduction route.
Extracting only one family is a component modification.

The Treasury input procedure admits monthly, quarterly, and daily frequencies. Selecting another
frequency for a self-contained extraction is a component modification. Monthly or daily output
cannot replace the required quarterly panel without coordinated changes to the calendar joins. The
Treasury source setting has three values: automatic resolution, the replication archive, and the New
York Fed provider. The first two resolve to the same monthly path, so choosing between them is an
inert executable full-route modification. The New York Fed value is an accepted-but-refused choice:
it passes the source check, but no preserved copy of that provider's data exists and its annual-only
maturity nodes fail the one-through-120-month requirement, so the route stops before using it.
Requesting that provider at daily frequency is an invalid choice, refused outright.

\subsection{SDF construction and expected-SDF lag}

The expected-SDF calculation accepts either the unpaired mean correction in Part I or a paired mean
correction. The paired correction calculates the correction only on observations shared by the
realized one-quarter price and the approximated expected SDF. The calculation also accepts a zero
maturity, for which it returns the exact realized one-quarter price and applies no approximation
correction. Both changes are component modifications because the complete reproduction route fixes
the unpaired correction and positive maturities.

The lag between the current and lagged expected-SDF panel matrices may change from one quarter.
Changing that lag is a component modification. A coordinated change must use a whole number of quarters, relabel
the lagged expected-SDF panel matrix, reconstruct its PC scores and sign alignment, and update every
sample, output label, and reusable-result identity affected by the lag.

\subsection{Consumption and supplied return scores}

The consumption input has exactly two selectable origins, and the choice between them is a full-route
modification. The preserved origin, which Part I uses, reads the stored quarterly consumption copy
and retrieves nothing; a missing stored copy stops the reproduction before any output is created or
cleared. The retrieving origin obtains the series over a stated window and replaces the stored copy
with the newly retrieved vintage. Any other origin is refused before calculation begins. Selecting
the retrieving origin changes the consumption vintage and therefore every consumption growth
observation, both samples, every mean and conditional-scale estimate, every interval, and every
reported number. It also overwrites the preserved copy that the preserved origin would otherwise
reproduce, and it invalidates every preserved resampling result.

The consumption series identifier, retrieval start date, and retrieval end date are consulted only
on the retrieving origin. Under the preserved origin they are inert and changing them changes
nothing. On the retrieving origin they are full-route modifications as long as the retrieved series
is quarterly, positive where growth is calculated, and long enough to supply every required lag.
The retrieving origin also admits the acquisition controls that suit a network environment: an
access credential selects a keyed service, and without one a public retrieval is used. Neither
choice changes any mathematical transformation. The keyed service permits changes to its attempt
count, timeout, initial retry delay, and maximum retry delay. The public retrieval permits changes
to redirect and network-protocol settings, retry count and retry conditions, connection and total
timeouts, retry delay, minimum transfer speed, time below that speed, and message verbosity. Every
one of these controls is inert under the preserved origin. Each is an executable full-route
modification that changes nothing while that origin is selected.

The PC window may change through its first and last quarter. The preliminary mean
sample, identification sample, and conditional-scale sample continue to arise from their stated
complete-row rules. A later start or earlier end reduces the available observations and changes
all estimates, block lengths, and numerical rank checks. Because the block length enters the recorded
reusable-result identity, any change to the window invalidates every preserved resampling result. It
is an executable full-route modification.

The supplied return-score archive contains six score columns. The complete reproduction route uses four because
the conditional-scale design, covariance matrices, bootstrap schemas, and reported tables all have
four return-score coordinates. Selecting different score columns while retaining four coordinates
is a full-route modification after their identities and interpretation are updated. Changing the
number of return-score coordinates is a component modification because the fixed output dimensions
must change with it.

\subsection{Principal-component construction}

The three SDF PC constructions share one preprocessing choice: center the maturity columns, divide
them by their sample standard deviations, do both, or do neither. The supplied return scores have a
separate choice with the same four possibilities. These choices are full-route modifications when
the numbers of scores remain three for each SDF panel matrix and four for the supplied return
scores. Giving the three SDF panel matrices different preprocessing choices is a component
modification because their current PC construction shares one choice.

The PC constructions accept other positive numbers of scores up to the available matrix rank. The
supplied archive limits that count accordingly. These count changes are component modifications. The
mean-equation geometry and its region figures require three PC scores from each SDF panel matrix, and the conditional-scale
schemas require four supplied return scores. A different dimension needs new formulas, figure
geometry, output dimensions, and acceptance rules.

\section{Mean-equation modifications}

\subsection{Specification plan}

The mean calculation permits specification A, specification B, or both. If it contains both, either
order is permitted. The first entry becomes the reported specification, and the later entry becomes
the comparison specification. The conditional-scale calculation must use exactly the reported
specification. Thus the permitted plans are A alone, B alone, A followed by B, and B followed by A.
No plan may repeat a specification or omit both.

Changing the plan changes the meaning of the SDF-news reduced-form residuals. Specification A sets
the SDF-news reduced-form coefficient matrix \(\beta_{2R}\) to zero and sets the SDF-news
reduced-form residual vector \(W_{2t}\) equal to the SDF-news PC score vector \(Y_{2t}\).
Specification B estimates \(\beta_{2R}\) and \(W_{2t}\). The coefficient-recovery equation,
mean identified set, structural mean residual, bootstrap, and every reported label must follow the
selected reported specification.

A plan that drops one of the two specifications also invalidates declared outputs. The comparison
between the reported and comparison specifications is a required step, and the machine-readable
comparison table is fixed at 56 rows, 28 for each specification. Dropping a specification therefore
changes that table's declared row count and removes the comparison bootstrap that fills half of it.
Exchanging the two specifications in either direction is an executable full-route modification.
Dropping either specification is an executable component modification: it removes a required
comparison step and contradicts the declared row count above, so the complete reproduction route
cannot carry it without coordinated changes to that step and that schema.

\subsection{Instrument and slack controls}

Replacing the instrument is an executable full-route modification that the current configuration does
not exercise. It changes the instrument sequence, all seven sample moments, the quadratic terms of
every inequality, the mean identified set, and every estimate and interval built on them, and it
invalidates every preserved resampling result. The instrument may be replaced by another
realized-yield-volatility column produced from the daily Treasury panel. Its maturity must lie between 1 and 120 months, and the replacement must be present
on the identification sample. Its label and reported description must change with the column. A
different observed instrument may also be supplied as a component modification, but it must enter
the quarter join, centering rule, reduced forms, input origin, and every instrument diagnostic before
the complete reproduction route can use it.

The slack system permits these full-route modifications.

\begin{enumerate}
\item The baseline slack may be any positive number below the slack cap.
\item The slack cap may be any number strictly between zero and one and above every requested slack.
\item The display-slack vector may contain any distinct positive values below the cap, provided its
first value equals the baseline slack. It fixes the positive-slack columns of every reported table and
the slack axis of every machine-readable inference table, so its length sizes those rows and columns.
\item Projection slacks must be drawn from the display-slack vector. They fix the slacks drawn in the
two-coordinate projection figures.
\item The slack vector used for fitted conditional-scale envelope figures may contain any distinct
positive slacks below the cap and must contain the baseline slack. Each entry yields one
fitted-envelope figure for each of the two exponential estimators, and four combined panels are drawn
per estimator regardless, so a vector of five slacks gives the current 18 sweep figures.
\item The slack vector drawn by the three-dimensional mean-identified-set figures is separate from
every vector above and is deliberately held outside the recorded analysis settings, so changing it
discards no preserved resampling result. It may contain any distinct positive slacks, each of which
must have a declared output location, because a combination with no declared location cannot be
written. Each entry yields one figure in each of the two unit systems, with and without the projected
mean least-squares benchmark, so the current four slacks give the 16 three-dimensional region
figures. Changing it changes the number and identity of those figures and therefore the stated output
inventory counts, and it leaves every estimate and every interval unchanged. It is an executable
full-route modification.
\item The coarse slack step may be any positive number below the cap. It fixes the coarse slack sequence that
brackets the critical slack, so changing it changes that reported value.
\end{enumerate}

The critical-slack calculation may also change the number of bisection iterations. More
iterations improve numerical resolution, but increase work; the change is an executable full-route
numerical modification that alters the reported critical slack and nothing else. The cap and the
coarse step, classified above, also bound and space the coarse slack sequence that brackets the critical slack, so
they move that value too. Alongside the iteration count sits a second number, the point count of a
fine slack grid. That grid divides into equally spaced parts the interval from zero to the
smallest slack the coarse slack sequence did not certify bounded, or to the largest coarse slack when the
coarse slack sequence certified every one bounded, keeps only the strictly interior division points, and drops any point already on
the coarse slack sequence. Its point count must be at least one. No stage of the complete reproduction route
evaluates that lattice, so changing the count is an executable component modification: the
calculation runs when invoked directly, but no value of it changes any reported quantity.
The same slack definitions must feed the mean identified set, conditional-scale estimator mappings, bootstrap,
figures, schemas, and output labels.

\subsection{Mean-identified-set and variance-share numerical controls}

The quadratic searches permit changes to feasibility tolerances, boundary tolerances, box widths,
box-expansion rounds, multistart counts, candidate-separation thresholds, evaluation limits, and
unboundedness-certification thresholds. Six further tolerances of the same calculation may also
change: the floor below which a constraint's scale is treated as degenerate, the tolerance for
accepting a quadratic term as symmetric, the relative convergence tolerance of the constrained
search, the tolerance below which the mean identified set is treated as a single point, the factor by
which a polished objective may worsen before the polish is rejected, and the tolerance for treating an
attained point as having escaped its search box. All of these are executable full-route numerical
modifications when applied consistently and recorded with the result. Each changes only the numerical
path and the accepted endpoints, never a population definition, and each invalidates the recorded
identity of the quantities it touches, hence every preserved resampling result.

The general constrained quadratic calculation accepts sequential least-squares quadratic
programming or constrained linear approximation. The complete reproduction route does not fix that
choice. Each conditional-scale endpoint polish takes the gradient-based sequential least-squares
programme when the estimator declares a smooth objective and supplies a coefficient gradient, and the
derivative-free constrained-linear-approximation programme otherwise. Every other search fixes the
gradient-based programme: the mean-identified-set box search, the variance-share searches, the shared
profile search, the log-OLS bound calculation, and the joint-null distance search. One further search
fixes the derivative-free programme instead: the joint-null maximum-norm sensitivity refinement.
Because log-OLS,
PPML, and Harvey QML declare smooth objectives with gradients while LAD declares a nonsmooth
objective and supplies no gradient, both programmes execute in the current configuration.

The activating choice is an estimator's declared objective smoothness. Its admissible domain is
exactly two values, smooth and nonsmooth, and withholding a coefficient gradient has the same effect
as declaring the objective nonsmooth. Changing it changes which programme polishes that estimator's
conditional-scale endpoints under identical normalized constraints, and therefore that estimator's
reported identified ranges and every interval built on them. It changes the recorded estimator
identity and invalidates every preserved resampling result. It is an executable full-route
modification, and both of its values are already exercised.

The same calculation can scale the objective by coefficient magnitudes or leave it unscaled, and each
search fixes its own scaling before any reproduction begins. The mean-identified-set box search and the design-functional endpoint searches of
Section~\ref{sec:idset} scale by coefficient magnitudes; the variance-share searches and the
conditional-scale endpoint polish leave the objective unscaled. Both values are therefore already
exercised, and no setting selects either one. The admissible domain is those two values. Changing one
search's scaling changes that search's numerical path and the endpoints it returns, and invalidates
the recorded identity of every quantity that search produces. It is an executable component
modification requiring a coordinated change at each of the four searches, one of which the
conditional-scale estimators and the log-OLS crossing bounds both
route through.

The variance-share calculation uses one common lattice count on every coefficient axis. That count
may be any integer of at least two. Its \emph{covariance-coherence check} requires each group
variance-share range to reach a stated fraction of its largest individual-PC contribution, with a
nonnegative allowance for the lower endpoint. The fraction may lie in \((0,1]\), the allowance may
be nonnegative, and the separate rendering tolerance must be strictly positive. These controls
change the density of the attained inner search and the numerical checks applied to its results;
they do not change the population variance-share definitions. Each is an executable full-route
numerical modification. Each changes the reported variance-share ranges and the rounded contributions
displayed beside them, and each invalidates the recorded identity of the variance-share table while
leaving every mean-equation and conditional-scale estimate unchanged. Tightening the
covariance-coherence check can also stop the reproduction on results the current configuration
accepts.

\section{Conditional-scale modifications}

\subsection{LAD decision}

The \emph{LAD decision} has three values: approved, declined, and unanswered. Approval makes LAD a
required part of a successful complete reproduction route and requires the numerical reference at exactly the recorded
release. An absent numerical reference or a release mismatch stops the reproduction. Declined and
unanswered decisions omit LAD and require all three LAD output locations to remain empty. These decision values are
executable full-route modifications.

The same record carries a second closed field, the consent to install the numerical reference, whose
four values are not asked, not needed, declined, and approved. A value outside that set is refused
before calculation and is an invalid choice, as is a record that names a reference release without
also recording an approved decision, which is refused and stops the reproduction. An absent record
reads instead as unanswered and takes the empty-output route, so no LAD calculation is even loaded.
Selection between the approved route and the empty-output route belongs to the decision field
classified above, not to this one: the decision changes which of the three LAD outputs exist and
whether the route-status result reports LAD as completed, and leaves every mean-equation and non-LAD
conditional-scale quantity unchanged. The consent field itself is inert beyond its own check: no
stage reads it afterwards and it reaches no output, so choosing among its four admissible values is
an executable full-route modification that changes nothing.

The LAD component accepts a quantile and two numerical procedures. The complete reproduction route fixes the
median quantile, uses Barrodale--Roberts for the retained coefficient vector, and uses
Frisch--Newton only to detect nonunique minimizers. Another quantile or a reversal of those roles is
a component modification because it changes the estimated scale and requires new interpretation,
labels, output rules, and validation.

\subsection{Estimator search controls}

The shared conditional-scale search accepts these full-route numerical modifications:
\begin{enumerate}
\item Change the lattice density \(n_{\mathrm{lat}}\), the minimum number of feasible lattice points that triggers a
denser lattice, or the number of calculations performed in one chunk. The lattice density must leave
the densified full lattice within one million points, that is, \((2n_{\mathrm{lat}}-1)^{3}\le10^{6}\)
on the three structural SDF-news coordinates, so \(n_{\mathrm{lat}}\le50\); a larger value is
refused before any conditional-scale calculation.
\item Change the traversal order or the nearest-neighbor traversal limit. Nearest-neighbor traversal
visits the closest unvisited feasible lattice point after each accepted fit so that the preceding fit
can serve as a warm starting value, and the limit is the largest point count for which this traversal
is permitted. The alternative scans the feasible lattice in its own enumeration order and is not
subject to that limit. An independently supplied lattice selector may likewise declare its own scan
order, which also bypasses the limit. Each of the two settings has exactly two values, and both
orders are exercised in the current configuration: LAD declares the lattice enumeration order, and
the PPML coverage audit supplies a selector that declares its own. The order fixes the
warm-start chain and therefore which fits converge, the scan extrema, and the polished endpoints; it
changes the recorded estimator identity and invalidates every preserved resampling result.
\item Change the numbers of primary and independent-check starting values, their minimum
separation, or whether cold-start endpoint replications run.
\item Change the endpoint-agreement tolerance, which compares a primary endpoint with its audit;
the nesting tolerance, which checks that identified ranges do not shrink when slack rises; or the
containment tolerance, which checks that the zero-slack point lies in every positive-slack range.
The return-score-centering tolerance is separate from the three above and does not reject a search
result. It is a positive bound on the absolute sample mean of each centered lagged supplied return
score, checked before the PPML, Harvey QML, and LAD conditional-scale searches and not before the
log-OLS search, as Section~\ref{sec:logvar-scale} states; reaching or exceeding it stops the
reproduction. Raising
or lowering it changes only which samples are admitted, never a reported value.
\item Change the lattice density and retained-point cap used for the fitted conditional-scale
figures of every estimator, the number of starting values used for fitted conditional-scale endpoints, or the
fitted-envelope fresh-fit allowance, which Part~I fixes at 80,000. Each changes which candidates the
envelope search reaches and therefore the reported envelope, and each invalidates the recorded
identity of every quantity that search produces.
\item Change the reduced lattice cap and fresh-fit allowance that the bootstrap stage applies inside
every resampled draw, which Part~I fixes at 1,500 and 8,000. Each changes which candidates every draw
reaches, and therefore every reported interval, and each invalidates every preserved resampling
result.
\item Change the primary lattice cap, the primary fresh-fit allowance, or the sensitivity fresh-fit
allowance of the exponential conditional-variance estimators. Each of these three is one setting shared by
both estimators rather than two independent ones: assigning PPML and Harvey QML different values is refused
before any calculation and is therefore an invalid choice. Only the PPML coverage audit records its
own lattice cap and fresh-fit allowance, which is why it scans more lattice points than the primary
mapping. The Harvey sensitivity audit reuses the shared primary lattice cap together with the shared
sensitivity allowance and exposes no independent cap of its own.
\end{enumerate}
Increasing a count or allowance can find more candidates, but does not certify a global optimum.
Tightening a tolerance can reject a result that the current configuration accepts.

The estimator-specific controls permit these further modifications:
\begin{enumerate}
\item For log-OLS, change the cold-start coefficient-agreement tolerance. This control is inert:
log-OLS performs no cold-start endpoint replication and uses no warm starting values, so no value of
it changes any reported quantity. A shared fallback value
exists for an estimator that declares none, but the three estimators that perform cold-start
endpoint replication all declare their own, so that fallback is inert.
\item Change the crossing-range relative tolerance of the quadratic system, which Part~I fixes at
\(10^{-8}\) and which widens the comparison of an outcome reduced-form residual against the
accepted endpoints of the corresponding \(W_{2t}'\theta\) over the mean identified set. It must be positive. It
changes which quarters are declared crossing quarters and therefore which endpoints become unbounded
or unreliable, and it invalidates the recorded identity of every conditional-scale quantity. It is an
executable full-route numerical modification.
\item For PPML, change the iteration tolerance and limit; estimating-equation-residual, rank, condition-number, Jacobian,
and cold-start tolerances; and the four acceptance checks for a coefficient boundary, a finite
fitted mean, full rank, and cold-start agreement. The ordered fallback starting values may also
change. The deterministic lattice order may change the number of binary digits taken from each
coordinate, provided that count times the number of coefficient axes does not exceed the exactly
representable integer width of the numeric type; a larger product is refused before the ordering is
built.
The PPML normalization check decides whether to divide all squared structural mean
residuals by one fixed positive PPML normalization divisor before fitting; its overflow margin,
condition-number limit, and lattice count may change. The PPML coverage audit may change its lattice cap and
fresh-fit allowance.
\item For Harvey QML, change the first-order-condition-norm tolerance and the tolerances that
determine the rank of the conditional-scale design matrix \(X^V\) restricted to quarters with
positive squared structural mean residuals and the singular-value rank of the corresponding
normalized matrix \(Z_+^H\). The condition-number, Jacobian,
Newton-step, recession-rate, and infeasibility-certificate tolerances may also change. The
objective-noise multiplier sets the increase in the objective treated as numerical noise; the
first-order-condition-progress multiplier sets the improvement in the first-order-condition norm
required after such an increase. Both may
change, as may the line-search depth, iteration limit, relative-change tolerance, auxiliary
linear-programming tolerance, evaluation limit, coefficient bound, and cold-start tolerance. The
starting-value order may change among the warm Harvey QML value, PPML at the current structural
SDF-news coefficient vector, the accepted zero-slack PPML value, the shifted log-OLS benchmark
value, and the intercept-only value. The coefficient scaling rule is
fixed at no scaling; every other value is invalid.
\item For LAD, change the lattice cap, total fresh-fit allowance, and separate allowances for scanning,
hyperplane approaches, nonuniqueness checks, endpoint refinement, and cold-start endpoint
replication.
The objective, coefficient, Frisch--Newton, near-hyperplane, crossing, feasibility, and cold-start
tolerances may change. The crossing limit and the evaluation limit and radius for an LAD side anchor
may also change. The approach ratio multiplies the distance to a structural-mean-residual hyperplane
at each successive probe; the approach count limits the number of probes. The two tail-window lengths
give the numbers of latest probes, closest to that hyperplane, used to assess a finite limit and
persistent divergence. The increment count gives the number of latest coefficient changes that must
be small for a finite-limit classification, and the tail-stability tolerance is the relative size,
measured against the largest coefficient magnitude in the window, below which such a change counts as
small. It must be positive; it changes only which crossings receive a finite-limit classification and
therefore the reported closure limits, and it is an executable full-route numerical modification that
invalidates the recorded LAD identity.

The tail-slope floor prevents division by a nearly zero slope, and the tail-slope agreement
tolerance compares slopes calculated over the two tail windows. The two movement thresholds require
the fitted tail movement to be large both relative to numerical tolerance and relative to the
coefficient level. The minimum tail span is the smallest permitted range of transformed log distance
over a divergence window, and the minimum \(R^2\) is the weakest permitted linear fit over that
window. The edge fractions position crossing-search starting values between the midpoint and edges
of the structural SDF-news coefficient box. The bisection count limits successive halvings used to
restore exact feasibility. The simultaneous-crossing cap limits the number of structural-mean-residual hyperplanes
treated at one crossing. The cone-step fraction is the fraction of the local coefficient-box scale
used to move away from those hyperplanes when checking a feasible direction. Each control may change.
\end{enumerate}
These changes remain full-route numerical modifications when the estimator's defining objective and
reported scale stay unchanged.

A separate LAD refinement calculation accepts its own refinement radii, which determine the sizes
of its local searches. The complete reproduction route does not reach this calculation. Changing
these radii is therefore a component modification.

\subsection{Harvey recession-classification certificate}

Part I states when the Harvey QML recession classification is invoked, its four verdicts, and the
consequence of each. This subsection states the facet certification that produces a verdict in the
one case the two immediate branches leave open: at least one squared structural mean residual
exactly zero and at least one strictly positive. The complete reproduction route enters it only at
an exact zero structural mean residual. The per-slack self-test never enters it, because its two
classification calls are settled by the immediate branches. The enumeration itself is fixed: it runs
over twice as many facets as there are conditional-scale coefficients, with no count, ordering, or
subsetting control. This
subsection specifies apparatus that the complete reproduction route itself performs whenever that
case arises, so the apparatus is not a modification of any class; it is stated here only to keep
Part~I's account of the safeguard compact. What may change are the auxiliary linear-programming
tolerances named above, each an executable full-route numerical modification.

Solve every program below by sequential least-squares quadratic programming with the analytic linear
objective gradient and the analytic inequality and equality Jacobians. Let \(n_+\) be the number of
rows of \(Z_+^H\), a positive integer.

On facet \((j,\sigma)\), the \emph{primal problem} is
\begin{equation*}
 m_{j,\sigma}=\min_{u\in\mathbb R^5}(c^H)'u
 \quad\text{subject to}\quad
 Z_+^Hu\ge0,\quad u_j=\sigma,\quad -1\le u_\ell\le1.
\end{equation*}
Start it from the direction whose coordinates are all zero except coordinate \(j\), which equals
\(\sigma\). Measure the feasibility violation of the returned direction \(u\) by
\begin{equation*}
 \max\left\{\max_{r}\left(-Z_{+,r}^Hu\right),\;|u_j-\sigma|,\;
 \max_{\ell}\left(|u_\ell|-1\right)\right\},
\end{equation*}
where \(r\in\{1,\ldots,n_+\}\) indexes the rows of \(Z_+^H\) and \(\ell\in\{1,\ldots,5\}\) indexes
the coordinates. The facet counts as feasible when this violation is at most \(10^{-8}\). It counts
as apparently infeasible when the violation exceeds \(10^{-8}\) or the solve returns no finite
point.

For a feasible facet whose primal value exceeds \(\rho_H\), verify positivity with a \emph{dual
lower bound}, which combines the primal constraints with nonnegative weights to produce a number no
larger than the primal minimum. Write the facet inequalities as \(G_Hu\le h_H\), where
\begin{equation*}
 G_H=\begin{pmatrix}-Z_+^H\\ I_5\\-I_5\end{pmatrix},
 \qquad
 h_H=(\mathbf0_{n_+}',\mathbf1_5',\mathbf1_5')'.
\end{equation*}
Here \(G_H\in\mathbb R^{(n_++10)\times5}\), \(h_H\in\mathbb R^{n_++10}\), \(I_5\) is the
five-by-five identity matrix, \(\mathbf0_{n_+}\) is a vector of zeros, and \(\mathbf1_5\) is a
vector of ones. Solve for a multiplier pair \((\lambda,\nu)\in\mathbb R^{n_++10}\times\mathbb R\)
starting from the zero vector of length \(n_++11\). Bound every coordinate of \(\lambda\) to
\([0,10^6]\) and the scalar \(\nu\) to \([-10^6,10^6]\). The solve does not itself guarantee
\(\lambda\ge0\) to working precision: replace \(\lambda\)
by its coordinatewise positive part \(\lambda_+=\max(\lambda,0)\) before forming any certificate
quantity. Define the stationarity residual
\begin{equation*}
 q_H=c^H+G_H'\lambda_++\nu e_j,
\end{equation*}
where \(e_j\in\mathbb R^5\) has one in coordinate \(j\) and zeros elsewhere and
\(q_H\in\mathbb R^5\) measures failure of the weighted constraints to reproduce the primal objective
direction. The corrected dual lower bound is
\begin{equation*}
 L_H=-h_H'\lambda_+-\sigma\nu-\|q_H\|_1,
\end{equation*}
where \(\|q_H\|_1\) is the sum of the absolute coordinates of \(q_H\). Certify the facet positive
only when \(L_H>\rho_H\) and the relative difference between \(L_H\) and the primal value is at most
\(10^{-8}\).

For an apparently infeasible facet, solve the \emph{phase-I problem}, which adds a nonnegative
violation variable \(t\) and minimizes it to measure the smallest constraint violation:
\begin{equation*}
 \min_{u,t}t
 \quad\text{subject to}\quad
 -Z_+^Hu\le t\mathbf 1,\quad u_j=\sigma,\quad
 -1\le u_\ell\le1,\quad0\le t\le t_{\max},
\end{equation*}
where \(\mathbf 1\) is a vector of ones and
\begin{equation*}
 t_{\max}=\max\left\{1,\max_r\sum_{\ell=1}^5|Z_{+,r\ell}^H|\right\}.
\end{equation*}
Here \(t_{\max}>0\) bounds the phase-I violation variable. Start the phase-I problem from the pair whose
first five coordinates are zero except coordinate \(j\), which equals \(\sigma\), and whose last
coordinate equals \(t_{\max}\). A phase-I value at most \(10^{-8}\) sends the facet back to the
primal problem; on that return the facet can end only as a negative recession direction, a zero-rate
direction, or unresolved, and never as a certified positive facet. To certify infeasibility, let
\(v=(u',t)'\in\mathbb R^6\), \(a=(\mathbf0_5',1)'\), and
\begin{equation*}
 G_I=\begin{pmatrix}
 -Z_+^H&-\mathbf1_{n_+}\\
 I_5&\mathbf0_5\\
 -I_5&\mathbf0_5\\
 \mathbf0_5'&1\\
 \mathbf0_5'&-1
 \end{pmatrix},
 \qquad
 h_I=(\mathbf0_{n_+}',\mathbf1_5',\mathbf1_5',t_{\max},0)'.
\end{equation*}
Here \(G_I\in\mathbb R^{(n_++12)\times6}\), \(h_I\in\mathbb R^{n_++12}\), and \(G_Iv\le h_I\)
collects all phase-I inequalities. Also \(a'v=t\) is the phase-I objective. Let
\(e_{j,I}=(e_j',0)'\in\mathbb R^6\). Solve for \((\lambda_I,\nu_I)\in\mathbb R^{n_++12}\times\mathbb
R\) from the zero vector of length \(n_++13\), again bounding every coordinate of \(\lambda_I\) to
\([0,10^6]\) and the scalar \(\nu_I\) to \([-10^6,10^6]\), and again replacing \(\lambda_I\)
by its coordinatewise positive part before use. Define
\begin{equation*}
 q_I=a+G_I'\lambda_I+\nu_Ie_{j,I},
 \qquad
 L_I=-h_I'\lambda_I-\sigma\nu_I-\sum_{\ell=1}^5|q_{I,\ell}|-t_{\max}|q_{I,6}|.
\end{equation*}
Here \(q_I\in\mathbb R^6\) is the phase-I stationarity residual and \(L_I\in\mathbb R\) is its dual
lower bound. Certify the facet infeasible when \(L_I>10^{-8}\).

For every primal, dual, and phase-I calculation, require feasibility error and relative
primal--dual disagreement at most \(10^{-8}\), use relative iterate tolerance \(10^{-12}\), and
permit at most 2,000 evaluations.

\subsection{Analytic covariance choices}

For a sequence \(a_t\in\mathbb R^p\), where \(t=1,\ldots,n\), \(n\) is the positive integer
sequence length, and \(p\) is its positive integer dimension, define the lag-\(\Upsilon\) Bartlett
long-run covariance sum
\begin{equation*}
 \widehat M_\Upsilon(a)=
 \sum_{t=1}^{n}a_ta_t'
 +
 \sum_{j=1}^{\min(\Upsilon,n-1)}
 \left(1-\frac{j}{\Upsilon+1}\right)
 \sum_{t=j+1}^{n}(a_ta_{t-j}'+a_{t-j}a_t').
\end{equation*}
Here \(\Upsilon\) is any nonnegative integer lag count, and \(j\) indexes a lag from one through the
smaller of \(\Upsilon\) and \(n-1\). The symmetric matrix
\(\widehat M_\Upsilon(a)\in\mathbb R^{p\times p}\) is the uncentered Bartlett long-run covariance sum.
The second sum is empty when \(\Upsilon=0\), so \(\widehat M_0(a)=\sum_ta_ta_t'\).

PPML admits four analytic covariance choices. At a fixed structural SDF-news coefficient vector
\(\theta\) and an accepted PPML conditional-scale coefficient vector, let
\(v_t^P=\exp(X_t^{V\prime}\widehat\eta_P)\) be the PPML fitted conditional variance and let
\(s_t^P=X_t^V\{q_t(\theta)-v_t^P\}\in\mathbb R^5\) be the PPML estimating-equation contribution. Define
\begin{align*}
 A_P&=\sum_{t=1}^{T_v}v_t^PX_t^VX_t^{V\prime},&
 B_P&=\sum_{t=1}^{T_v}s_t^Ps_t^{P\prime},\\
 \widehat\Xi_P&=\frac{1}{T_v-5}
 \sum_{t=1}^{T_v}\frac{\{q_t(\theta)-v_t^P\}^2}{v_t^P}.
\end{align*}
Here \(A_P,B_P\in\mathbb R^{5\times5}\), and \(\widehat\Xi_P\ge0\) is the scalar Pearson
dispersion estimate. The labels HC0 and HC1 mean heteroskedasticity-consistent covariance without
and with the stated finite-sample multiplier. The label HAC means
heteroskedasticity-and-autocorrelation-consistent covariance. Autocorrelation is correlation between a sequence and its
lagged values. The four choices are
\begin{align*}
 \widehat V_P^{\mathrm{model}}&=\widehat\Xi_PA_P^{-1},&
 \widehat V_P^{\mathrm{HC0}}&=A_P^{-1}B_PA_P^{-1},\\
 \widehat V_P^{\mathrm{HC1}}&=\frac{T_v}{T_v-5}\widehat V_P^{\mathrm{HC0}},&
\widehat V_P^{\mathrm{HAC}}(\Upsilon)&=A_P^{-1}\widehat M_\Upsilon(s^P)A_P^{-1}.
\end{align*}
Each displayed \(\widehat V_P\) is a \(5\times5\) PPML coefficient covariance matrix.
The model choice is the exponential-model covariance. HC0 allows unequal residual variances without
a finite-sample multiplier. HC1 multiplies HC0 by \(T_v/(T_v-5)\). HAC allows both unequal residual
variances and serial correlation. At \(\Upsilon=0\), HAC equals HC0 and has no finite-sample multiplier.

At the same fixed \(\theta\) and an accepted Harvey QML conditional-scale coefficient vector, let
\(v_t^H=\exp(X_t^{V\prime}\widehat\eta_H)\) be the Harvey QML fitted
conditional variance, let \(r_t^H=q_t(\theta)/v_t^H\ge0\) be the Harvey variance ratio, and let
\(s_t^H=\tfrac12(1-r_t^H)X_t^V\in\mathbb R^5\) be its gradient contribution. Define
\begin{align*}
 E_H&=\frac12\sum_{t=1}^{T_v}X_t^VX_t^{V\prime},&
 A_H&=\frac12\sum_{t=1}^{T_v}r_t^HX_t^VX_t^{V\prime},&
 B_H&=\sum_{t=1}^{T_v}s_t^Hs_t^{H\prime}.
\end{align*}
Here \(E_H,A_H,B_H\in\mathbb R^{5\times5}\) are the expected-information,
observed-information, and gradient outer-product matrices. Harvey QML admits
\begin{align*}
 \widehat V_H^{\mathrm{expected}}&=E_H^{-1},&
 \widehat V_H^{\mathrm{observed}}&=A_H^{-1},&
 \widehat V_H^{\mathrm{GOP}}&=B_H^{-1},\\
 \widehat V_H^{\mathrm{robust}}&=A_H^{-1}B_HA_H^{-1},&
\widehat V_H^{\mathrm{HAC}}(\Upsilon)&=A_H^{-1}\widehat M_\Upsilon(s^H)A_H^{-1}.
\end{align*}
Each displayed \(\widehat V_H\) is a \(5\times5\) Harvey QML coefficient covariance matrix. GOP
abbreviates gradient outer product. The expected-information, observed-information, and GOP choices
invert \(E_H\), \(A_H\), and \(B_H\). The robust choice allows unequal residual variances. The HAC
choice also allows serial correlation. At \(\Upsilon=0\), the HAC choice equals the robust choice. None
uses a finite-sample multiplier.

The mean least-squares benchmark and log-OLS share a further covariance rule. Let
\(x_t^{\mathrm{cov}}\in\mathbb R^p\) be the relevant covariance-design vector, let
\(\widehat u_t\in\mathbb R\) be its fitted residual, let
\(s_t=x_t^{\mathrm{cov}}\widehat u_t\in\mathbb R^p\), and set
\(Q_{\mathrm{cov}}=\sum_{t=1}^{n}x_t^{\mathrm{cov}}x_t^{\mathrm{cov}\prime}
\in\mathbb R^{p\times p}\). Without
\emph{prewhitening}, set \(\Omega_\Upsilon^{\mathrm{cov}}=\widehat M_\Upsilon(s)\). Prewhitening first fits the no-intercept
first-order vector regression
\begin{equation*}
 s_t=\widehat B_{PW}s_{t-1}+r_t^{PW},\qquad t=2,\ldots,n,
\end{equation*}
where \(\widehat B_{PW}\in\mathbb R^{p\times p}\) minimizes
\(\sum_{t=2}^{n}\|s_t-B_{PW}s_{t-1}\|_2^2\) over \(B_{PW}\in\mathbb R^{p\times p}\), and
\(r_t^{PW}\in\mathbb R^p\) is the resulting residual vector. Let \(I_p\) be the \(p\)-dimensional
identity matrix, define \(D_{PW}=(I_p-\widehat B_{PW})^{-1}\), and set
\begin{equation*}
 \widetilde r_i^{PW}=r_{i+1}^{PW},\quad i=1,\ldots,n-1,
 \qquad
 \Omega_\Upsilon^{\mathrm{cov}}=D_{PW}\widehat M_\Upsilon(\widetilde r^{PW})D_{PW}'.
\end{equation*}
Here \(\widetilde r_i^{PW}\in\mathbb R^p\) reindexes the prewhitening residuals as a sequence of
length \(n-1\). The Bartlett definition therefore caps its lag at \(n-2\). Prewhitening removes
fitted first-order dependence and restores the original long-run scale with \(D_{PW}\). The matrix
\(\Omega_\Upsilon^{\mathrm{cov}}\in\mathbb R^{p\times p}\) is the resulting Bartlett long-run covariance sum.

Let the positive scalar \(\kappa=1\) without the finite-sample adjustment and
\(\kappa=n/(n-p)\) with it. The covariance is
\begin{equation*}
 \widehat V=\kappa Q_{\mathrm{cov}}^{-1}\Omega_\Upsilon^{\mathrm{cov}}
 Q_{\mathrm{cov}}^{-1}\in\mathbb R^{p\times p}.
\end{equation*}
For the mean least-squares benchmark, use \(n=N_O\), \(p=7\), \(x_t^{\mathrm{cov}}=x_t^O\), and
\(\widehat u_t=\widehat u_t^O\). For log-OLS, use \(n=T_v\), \(p=5\),
\(x_t^{\mathrm{cov}}=X_t^V\), and \(\widehat u_t=\widehat u_t^L\). Every required inverse must exist. The
PPML and Harvey QML inverses must also pass the selected reciprocal-condition-number tolerance.
Every selected covariance rule must be named in the table notes and reusable-result identity. A
covariance selection is activated by naming a different rule, lag count, prewhitening setting, or
finite-sample multiplier in the reporting configuration. The admissible domain is exactly the rules
displayed above: four for PPML, five for Harvey QML, and, for the mean least-squares benchmark and
log-OLS, any nonnegative lag count with prewhitening on or off and the finite-sample multiplier on or
off. Part I fixes the current selection, which is the heteroskedasticity-and-autocorrelation-consistent
rule at four lags for
all four fits, with no prewhitening and no finite-sample multiplier for the two least-squares fits. A
change replaces the displayed covariance matrix and therefore every standard error, test statistic,
p-value, significance star, and analytic confidence statement built on it, while leaving every point
estimate, every identified range, and every resampled endpoint unchanged. It enters the
reusable-result identity and so invalidates every preserved resampling result. Each of these
selections is an executable full-route modification.

The estimator contract reserves a per-side probe entry that an estimator may supply for diagnostics
only. It is explicitly not status-bearing: no endpoint state, identified range, or interval depends on
it, and no estimator on the complete reproduction route supplies one. It is a declared unavailable
extension.

The estimator catalogue can describe another estimator, its required inputs, reported scale, and
outputs. That description alone is a declared unavailable extension. A new estimator becomes a
full-route modification only after it has a defining conditional-scale estimator mapping, a search
over the mean identified set, a bootstrap-draw definition when applicable, output schemas, reporting
rules, and acceptance checks.

\section{Diagnostic modifications}

\subsection{Mean-equation diagnostics}

Exactly two heteroskedasticity collections are configured: the collection Part~I reports, and a
second collection that differs from it only by omitting the Anscombe calculation. The configuration
is refused before any calculation if it holds any other number of them. Either collection may select any subset of the White,
studentized Breusch--Pagan, Goldfeld--Quandt, Harvey, and Anscombe calculations defined in
Section~\ref{sec:mean-diagnostics}. Naming the Cook--Weisberg calculation in either collection is
refused at the same point and is therefore an invalid choice, not merely an omission. Selecting a
different subset is an executable full-route modification: it changes the rows of that collection's
table, and it changes no estimate, no identified range, and no interval.
Separately, the positive threshold on the fitted-standard-deviation ratio that decides whether the
Anscombe calculation is included is its own control, not part of any heteroskedasticity collection. It must be positive; it changes only whether that one row appears, and it is an executable full-route modification. It
does not move the unreported relevance summary, which reads only three
p-values. The Glejser, Breusch--Pagan LM, and ARCH(1)
calculations remain the supplementary calculations defined in that section. The rejection level is
likewise its own control, not part of any heteroskedasticity collection; it names one of the
significance thresholds and takes that threshold's value, so changing those thresholds moves it. So is the list of calculations whose
p-values that summary reads, currently the studentized Breusch--Pagan test, the Goldfeld--Quandt
test, and the ARCH(1) LM test. Changing either is an executable full-route modification that moves only the unreported
relevance summary of Section~\ref{sec:mean-diagnostics}.

For Goldfeld--Quandt, the component accepts any nonconstant predictor in the base regression as the
ordering variable. The configured control is narrower: it is a position among the available instrument
columns, and a position beyond the last column is clamped to the last one. Because the current
configuration supplies a single instrument column, every admissible position selects that column and
the choice is inert. It becomes an executable full-route modification, changing the two group
regressions and hence the reported p-value, only in a configuration with more than one
instrument column. The statistic itself is not reported: only p-values are retained for the table.
Let \(\mathcal F_{q_2,q_1}\) and \(\mathcal T_{GQ}\) retain the meanings defined in
Section~\ref{sec:mean-diagnostics}. The upper-tail, lower-tail, and two-sided p-values are
\begin{equation*}
 1-\mathcal F_{q_2,q_1}(\mathcal T_{GQ}),\qquad
 \mathcal F_{q_2,q_1}(\mathcal T_{GQ}),\qquad
 \min\!\left\{2\mathcal F_{q_2,q_1}(\mathcal T_{GQ}),
 2[1-\mathcal F_{q_2,q_1}(\mathcal T_{GQ})]\right\}.
\end{equation*}
The selected tail changes the alternative hypothesis but not the two group regressions. Selecting a
different tail is an executable full-route modification: it changes the Goldfeld--Quandt p-value
printed in the reported table, and it changes no estimate, no identified range, and no interval.

The complete reproduction route excludes the Cook--Weisberg calculation from the configurable
collection. For a regression with \(n\) observations, let
\(e_t^C\in\mathbb R\) be residual \(t\), let \(Z^C\in\mathbb R^{n\times(p_C+1)}\) be an auxiliary
predictor matrix whose first column is an intercept, and define
\begin{equation*}
 \widehat\sigma_C^2=\frac1n\sum_{t=1}^n(e_t^C)^2,
 \qquad c_t=\frac{(e_t^C)^2}{\widehat\sigma_C^2}.
\end{equation*}
Here \(p_C\) is the positive integer count of nonconstant auxiliary predictors,
\(\widehat\sigma_C^2>0\) is the residual-variance estimate, and
\(c=(c_1,\ldots,c_n)'\in\mathbb R^n\) is the standardized squared-residual vector. The
multiplicative form regresses \(c\) on \(Z^C\). If \(RSS_C\ge0\) is that auxiliary regression's residual
sum of squares, its statistic is
\begin{equation*}
 \mathcal T_{CW}=\frac12\left\{
  \sum_{t=1}^nc_t^2-n\bar c^2-RSS_C
 \right\},
 \qquad \bar c=\frac1n\sum_{t=1}^nc_t.
\end{equation*}
Here \(\mathcal T_{CW}\ge0\) is the scalar Cook--Weisberg statistic. Its p-value is the upper-tail
probability of a chi-squared random variable with \(p_C\) degrees of freedom. The component also
permits an additive form, which replaces each nonconstant column of \(Z^C\) by twice that column.
Doubling a column changes neither the span of the auxiliary predictors, nor the residual sum of
squares, nor the degrees of freedom, so this form leaves the statistic and its p-value unchanged and
is inert. The component also permits
a logarithmic multiplicative form, which replaces each such column by its natural logarithm. These
forms require the transformed auxiliary predictors to be finite.

A separate time-series diagnostic component accepts one finite scalar sequence
\(x_1,\ldots,x_n\), where \(n\) is its observation count. It may report sample autocorrelations
through any positive integer maximum lag smaller than \(n\), or omit all sample autocorrelations.
This component is absent from the complete reproduction route and the descriptive report. Every
choice catalogued for this time-series component and for the Cook--Weisberg component above is
therefore an executable component modification: the
calculations run when invoked directly, but no coordinated change short of adding a caller brings
them onto the complete reproduction route, and none of them changes any reported estimate, identified
range, interval, or output while they remain uninvoked.

The component's \emph{augmented Dickey--Fuller calculation} tests the null that \(x_t\) has a unit
root, meaning that a shock to the sequence does not decay. It estimates
\begin{equation*}
 \Delta x_t=\gamma x_{t-1}+d_t+\sum_{j=1}^{q}\phi_j\Delta x_{t-j}+u_t,
 \qquad \Delta x_t=x_t-x_{t-1}.
\end{equation*}
Here \(\Delta x_t\) is the scalar first difference, \(\gamma\in\mathbb R\) is the unit-root
coefficient, \(q\in\{0,1\}\) is the selected augmented Dickey--Fuller lag count, \(\phi_j\in\mathbb R\) is the coefficient
on lagged first difference \(j\), and \(u_t\in\mathbb R\) is the regression error. The deterministic
term \(d_t\) may be zero, an intercept, or an intercept plus a linear time trend. The fixed-lag,
Akaike-information, and Bayesian-information labels are all accepted. The component fixes the
maximum lag at one and, under each label, uses \(q=1\). The test statistic is the least-squares
t-statistic for \(\gamma=0\).

A response-surface approximation converts that statistic to a p-value. Its reference trend may be
no deterministic term, an intercept, an intercept plus a linear trend, or an intercept plus linear
and quadratic trends. Its reference statistic may be the t-statistic or the sample-size-normalized
coefficient statistic. These reference choices affect only the p-value approximation; they do not
change the estimated augmented Dickey--Fuller regression.

The component's \emph{Kwiatkowski--Phillips--Schmidt--Shin calculation}, abbreviated
\emph{KPSS}, tests a stationarity null, meaning that the sequence returns
to a fixed mean or a fixed trend after a shock. Remove either a constant or a constant and linear
trend from \(x_t\), and call the resulting residual \(r_t\in\mathbb R\). Define
\begin{align*}
 S_t&=\sum_{i=1}^{t}r_i,&
 \Upsilon&=\left\lfloor4(n/100)^{1/4}\right\rfloor,\\
 \widehat\omega^2&=\frac1n\sum_{t=1}^{n}r_t^2+
 \frac2n\sum_{j=1}^{\Upsilon}\left(1-\frac{j}{\Upsilon+1}\right)
 \sum_{t=j+1}^{n}r_tr_{t-j},&
 \mathcal T_K&=\frac{n^{-2}\sum_{t=1}^{n}S_t^2}{\widehat\omega^2}.
\end{align*}
Here \(S_t\in\mathbb R\) is the cumulative residual, \(\Upsilon\) is the nonnegative integer Bartlett lag
count, \(\widehat\omega^2\) is the scalar long-run residual-variance estimate, and
\(\mathcal T_K\ge0\) is the KPSS statistic. Its p-value is linearly interpolated between the 10,
5, 2.5, and 1 percent critical values for the selected constant or trend form and is limited to the
interval \([0.01,0.10]\).

For a selected positive integer lag \(m_S<n\), define the lag-\(j\) sample autocorrelation, for
\(j\in\{1,\ldots,m_S\}\), by
\begin{equation*}
 \widehat\rho_j=
 \frac{\sum_{t=j+1}^{n}(x_t-\bar x)(x_{t-j}-\bar x)}
 {\sum_{t=1}^{n}(x_t-\bar x)^2},
 \qquad \bar x=\frac1n\sum_{t=1}^{n}x_t.
\end{equation*}
Here \(\widehat\rho_j\in[-1,1]\) is the scalar sample autocorrelation at lag \(j\), and
\(\bar x\in\mathbb R\) is the sample mean. The serial-correlation calculation may use
\begin{equation*}
 Q_P=n\sum_{j=1}^{m_S}\widehat\rho_j^2
 \qquad\text{or}\qquad
 \mathcal T_{LB}=n(n+2)\sum_{j=1}^{m_S}\frac{\widehat\rho_j^2}{n-j}.
\end{equation*}
Here \(m_S\) is the positive integer serial-correlation lag, \(Q_P\ge0\) is the lag-\(m_S\)
Box--Pierce statistic, and \(\mathcal T_{LB}\ge0\) is the lag-\(m_S\)
Ljung--Box statistic. Each p-value is the upper-tail probability of a chi-squared random variable
with \(m_S\) degrees of freedom; no fitted-parameter adjustment is subtracted. If a caller invoked
them, all time-series diagnostic settings would change diagnostics rather than the mean or
conditional-scale estimates; while the component stays uninvoked they change nothing.

\subsection{Conditional-scale diagnostic controls}

The joint-null compatibility check permits coordinated changes to these numerical controls:
\begin{enumerate}
\item Change the number of objective values calculated in one chunk and the two multipliers that
translate machine precision and root tolerance into objective-agreement tolerances.
\item Change the lattice density, minimum retained lattice count, root tolerance, feasibility
tolerance, hyperplane-crossing tolerance, coordinate-agreement tolerance, fresh-calculation
tolerance, or active-box tolerance. The active-box tolerance determines when a candidate is close
enough to a limit of the coefficient box to count as box-limited. The machine-precision-adjacent tolerance
is the relative size below which the smallest structural mean residual counts as indistinguishable
from zero, which forces a not-demonstrated outcome; it may change with the rest. Every tolerance in
this item must be positive, and each changes only the joint-null outcome and its recorded witness,
not any estimate.
\item Change the required number of agreeing starting values, the number and separation of those
values, the perturbation sizes, and the rank, duplicate-direction, and direction-sign tolerances.
\item Change the neighborhood removed around a zero residual; the tolerances for departures from the
constraint surface and maximum-norm gaps; and the auxiliary constrained minimization's relative
tolerance, evaluation limit, coefficient bound, and constraint tolerance.
\end{enumerate}
Every joint-null control in the four items above is an executable full-route numerical modification:
all are read on every run and none is pinned by a recorded decision.

The active stacked-moment replication check permits changes to its three-by-three-by-three lattice
density, retained-point cap, and normal-equation tolerance. Each is an executable full-route
numerical modification, since the check runs under the current \emph{no-answer default}, the recorded
decision in which no moment block is enabled. These settings change only the
replication check; they add no identifying moment.

\subsection{Residual-dynamics decisions}

Changing any residual-dynamics gate control is an executable component modification rather than a
full-route one: the committed scope decision binds the gate record's protocol, so a changed control
makes the freshly generated record disagree with that binding and stops the reproduction until the recorded
decision is regenerated against it. With that coordinated change, the check permits different
positive tested lags, one tested lag as the decision
lag, a significance level strictly between zero and one, a nonnegative autocorrelation display limit,
and a nonnegative numerical tie tolerance. Changing any of these values invalidates the recorded
residual-dynamics outcome and requires a new outcome tied to the new residual-dynamics check and inputs.
An unreliable residual-dynamics outcome retains the same reusable-result field order and uses an
empty table for the residual sequence.

The \emph{residual-dynamics decision sequence} applies these rules in order:
\begin{enumerate}
\item An unreliable or non-reject residual-dynamics outcome ends the sequence and records both
later decisions as not asked.
\item A reject residual-dynamics outcome opens two ordered decisions. If either decision is
unanswered, the sequence ends without requesting a dynamic-volatility calculation.
\item The first decision asks whether to replace the coefficient on lagged supplied return scores
in the static conditional-scale projection with its counterpart in a dynamic-volatility model. A
decline ends the sequence. An approval continues it.
\item The second decision asks whether to use the required dynamic numerical reference at its
recorded release. A decline ends the sequence. An approval continues it.
\item If both decisions are approved, the numerical reference must be present at the approved
release. Its absence or a release mismatch stops the reproduction. A matching reference permits the
route-status result to request a dynamic-volatility calculation.
\end{enumerate}

The residual-dynamics decision sequence is a routing-only choice, not a dynamic-estimation
capability. The complete reproduction route has no dynamic-volatility estimator. Even a consistent reject outcome followed by both
approvals can record a request, but cannot construct a dynamic-volatility estimate or its four reserved outputs.
The dynamic-volatility calculation is therefore a declared unavailable extension.

The unavailable extension would change the estimand, meaning the population quantity being
estimated. The static conditional-scale slope vector \(\eta_R\in\mathbb R^4\) is defined by the projection
\begin{equation*}
 \log e_t(\theta)^2=a_L+R_t'\eta_R+u_t^L(\theta),
\end{equation*}
where \(a_L\in\mathbb R\) is the mean-log intercept, \(R_t\in\mathbb R^4\) is the vector of
centered lagged supplied return scores for quarter \(t\), and \(u_t^L(\theta)\in\mathbb R\) is the
log-OLS residual. Both sides carry the same quarter index, as in Part~I. The
dynamic-volatility slope \(\delta_D\in\mathbb R^4\) would instead enter
\begin{equation*}
 \log h_{t+1}=\omega_D+R_t'\delta_D+\alpha_D z_t^D+
 \gamma_D\{ |z_t^D|-\mathbb E|z_t^D|\}+\beta_D\log h_t,
 \qquad z_t^D=\frac{e_t(\theta)}{\sqrt{h_t}}.
\end{equation*}
Here \(h_t>0\) is the conditional variance of the structural mean residual, \(z_t^D\in\mathbb R\)
is its standardized value, \(\omega_D\in\mathbb R\) is the dynamic intercept,
\(\mathbb E|z_t^D|\) is the expected absolute standardized residual,
\(\alpha_D\in\mathbb R\) measures the signed standardized-residual effect,
\(\gamma_D\in\mathbb R\) measures the absolute standardized-residual effect, and
\(\beta_D\in\mathbb R\) measures persistence in log conditional variance. The vector
\(\delta_D\) measures the partial effect of \(R_t\) while holding \(z_t^D\) and \(h_t\) fixed. It is
therefore not the static conditional-scale slope vector \(\eta_R\). The route can record a request for this estimand, but contains no procedure that estimates it.

\section{Resampling and inference modifications}

\subsection{Draw count, random numbers, and workers}

The bootstrap draw count may be any integer from 2 through \(2^{31}-1\), and changing it is an
executable full-route modification; a value that is not a whole number is an invalid choice, refused
before the run begins. A value other than 10,000 requires an
explicit acknowledgement that the resulting published files are a draft rather than the publication
configuration. The resampling seed may be replaced by any integer in
\(\{1,\ldots,2^{31}-1\}\). Changing the seed is an executable full-route modification. Such a change
changes every bootstrap-index family and invalidates every reusable bootstrap result.

The worker count may be any positive integer supported by the machine. Left unset it is derived from
the machine's own processor count, so a reproduction on different hardware runs a different number of
workers unless it is set explicitly. It changes scheduling, but must not change results, and changing it is an executable full-route modification, because every
sampled index sequence is generated before the deterministic
draw calculations begin. Another language must evaluate each bootstrap draw from its pre-generated
sampled index sequence independently of evaluation order.

The pseudo-random-number generator may be changed only as a component modification. The complete reproduction route and
cross-language equations fix MT19937, inversion for normal draws, and rejection for discrete
sampling. Another generator changes every bootstrap-index family and requires new cross-language equations,
recorded identities, and fixed acceptance comparisons.

\subsection{Reusable-bootstrap-result policy and block construction}

The bootstrap mode has exactly two values, and any other value is refused before calculation begins,
which makes a third value an invalid choice. Switching between the two admissible values is an
executable full-route modification: the route stays executable either way, no population definition
or reported estimate changes, and the only consequence is whether the preserved draws are loaded or
discarded and drawn again. \emph{Reuse mode} validates a
reusable bootstrap result against
all input, axis, estimator, numerical, random-number, and schema identities. A valid result is loaded as one
unit. A missing, unreadable, stale, incomplete, or invalid result causes a fresh run and an atomic
replacement after validation. \emph{Rerun mode} performs a fresh run without attempting reuse.
Reuse never combines old and new bootstrap draws.

The recorded identity also covers the calculation code that produces the draws and the computing
environment: the language version, the machine platform, the versions of the required external
calculation libraries, the source and namespace identities of the installed identification
calculations, the linear-algebra libraries in use, and the recorded analysis, inference, reporting,
serialization, and cache-schema control collections. A reproduction that differs in any of these
cannot reuse a preserved result even when every input and every setting matches. The identity of the
code that only formats and summarizes already-calculated draws is recorded for audit and is
deliberately excluded from the reuse decision, so a change confined to reporting keeps the stored
draws while any change to the draw calculation discards them.

The moving-block rule's positive scale and finite exponent are sealed constants of the resampling
protocol rather than named controls. No setting selects them, so activating a change means editing
the sealed protocol itself, and any edit must still produce a positive integer requested block
length. Such an edit changes every block length, every resampled index sequence, and every interval
built on them, and it changes the recorded block-rule description that the reusable-result identity
carries, so it invalidates every preserved resampling result. The complete reproduction route remains
executable afterwards, so it is an executable full-route modification that no configuration exposes.

The resampling protocol names a third index family beyond the primary and sensitivity families. It is
built only by a self-contained resample-and-evaluate calculation that generates its own family and
evaluates its draws immediately, and the resampling stage's own contract admits only the primary and
sensitivity families. The complete reproduction route never invokes that calculation, so the third
family is an executable component modification the route cannot reach.

The sampler's constraints below are derived rather than selected: no setting exposes them, and a
specification that violates one is refused before any draw is taken, which makes such a deviation an
invalid choice. The sampler replaces a requested block longer than the sample with
the identification-sample observation count \(N_O\). The requested sensitivity block length must remain exactly twice the requested
primary block length, and both bootstrap-index families must retain the same draw count. Each
permitted change must enter the reusable-result identity and every displayed sample description.

\subsection{Inference controls}

Every control in this paragraph is an executable full-route modification. Each changes the reported
intervals and gate outcomes and nothing else, except the progress-reporting interval, which changes
neither. The nominal error probability may be any number strictly between zero
and one-half; a value at or above one-half is refused by the Stoye comparison calibration and stops
the reproduction. The minimum valid-draw
share may lie in \((0,1]\), and the bounded-endpoint stability share may lie strictly between zero
and one. The pointwise target tolerance must be positive. The fatal bootstrap-draw calculation
failure share must lie strictly between zero and one, and the progress-reporting interval must be a
positive integer.

The calibrated confidence procedures permit changes to their root-search brackets and positive
tolerances, the correlation threshold used for nearly degenerate bivariate-normal calculations, the
positive integration tolerance, the minimum number of paired endpoint bootstrap draws, and the positive
multiple used to trim median-absolute-deviation outliers. Each bracket must order its lower endpoint
below its upper endpoint and contain a zero of the corresponding calibration equation. The correlation threshold must lie strictly between zero and one, and the
minimum paired-endpoint threshold must be at least two. Because the observed paired-endpoint count
is an integer, a noninteger threshold has the same effect as rounding it upward. These controls are full-route numerical modifications when the nominal confidence
level and all table notes change with them.

Log-OLS and LAD have no mean-identified-set bootstrap in the complete reproduction route. Adding one is a declared
unavailable extension because neither estimator has the required bootstrap-draw definition, storage fields,
confidence procedure, or reporting rule.

Which stages an estimator reaches at all is itself recorded, as a list of capabilities per estimator
covering the identified-range-by-slack figure, the published table, the mean-identified-set bootstrap,
and the fitted conditional-scale envelope. Withdrawing a capability an estimator currently holds is an
executable full-route modification: the corresponding stage no longer runs for that estimator, and the
output inventory loses the files that stage produces. Granting a capability an estimator does not
currently hold is the declared unavailable extension described above whenever the underlying procedure
is missing.

Each estimator also carries a version label that enters its recorded identity. Changing one is an
executable full-route modification: it changes no calculation, but it discards every preserved
resampling result. The same holds for the recorded version of any control collection that enters that
identity; the recorded versions of control collections outside it, such as the figure-rendering controls,
change nothing.

\section{Approximation-error modifications}

The approximation-error calculations permit other distinct positive maturities and other positive
news horizons. Every positive maturity must be an integer multiple of the selected news horizon.
The choices remain subject to the 120-month Treasury limit and the trimming rules in
Section~\ref{sec:approximation}. The reported maturity spacing can also change. These changes are
full-route modifications when both bound calculations, their common finite set, the quoted-number
checks, the figure, and the summary table use the same requested maturities.

The SDF-news plug-in approximation-error variance bound may replace the sample construction of the
exponential envelope \(\mathcal E_i^N\) with a supplied deterministic exponential envelope. The
supplied envelope is one finite, strictly positive number per maturity, given one maturity at a time;
a collection of values in a single request is refused. That replacement is a component modification because the complete reproduction route does not provide such an
input or record its origin. The paired expected-SDF plug-in approximation-error variance bound also
accepts a zero maturity and returns a zero bound with a warning. The complete reproduction route
excludes zero from its maturity sequence, so supplying a zero maturity is a component
modification.

\section{Reporting, preservation, and cleanup modifications}

\subsection{Table and number presentation}

The significance thresholds may change, provided they remain strictly ordered from the weakest to
the strongest evidence. Changing them is an executable full-route modification. It moves the printed
significance marks and, through the rejection level it supplies, the rejection count of the
unreported relevance summary of Section~\ref{sec:mean-diagnostics}, which Part~I records as reaching no
reported output. It changes no estimate, no identified range, and no interval. Because the printed marks are among the values the
numerical-table acceptance comparison must match, a changed threshold also changes what that
comparison requires. The HAC lag counts may also change; they change analytic covariances, standard
errors, test statistics, p-values, and reusable-result identities, but not point estimates.

Coefficient digits, statistic digits, scientific-notation digits, slack digits, the
unavailable-value symbol, endpoint-state labels, half-infinite-range display, open-interval notation, and
equal-endpoint display may change. Every digit count must be a nonnegative integer. Each is an
executable full-route presentation modification: it affects presentation and the comparison of
printed numerical tables defined below, and it does not change unrounded estimates.

Separate nonnegative precision controls govern descriptive summaries, correlations, diagnostic
tables, figure annotations, caption percentages, caption endpoints, and approximation-error bounds.
Table font sizes, cell spacing, row stride, and the table-note label may also change, as may the
number of digits used for console reporting and the typeface family of the three-dimensional
mean-identified-set figures. All of these are executable full-route presentation modifications: they
change only how a value is displayed, never a reported number, an identified range, or an interval.

Three significance-legend layouts exist: one listing the thresholds in ascending order as
percentages, one listing them in ascending order with a colon separator, and one listing them in
descending order by p-value. Only the ascending-percentage layout is reachable from a published
caption, and selecting it is the current full-route setting. The descending order is reachable only
from an inactive conditional-scale table-note builder that the complete reproduction route never
calls, and the colon-separated order is reachable from nowhere; both, and the inactive note builder
itself, are executable component modifications the route cannot reach. A layout change affects only
the wording of a note and changes no number.

\subsection{Figure presentation}

Figure presentation permits these groups of modifications. Every item in the list is an executable
full-route modification. Except where an item states otherwise, each changes only how a figure is
drawn, never a reported number, an identified range, or an interval.
\begin{enumerate}
\item General styling may change colors, transparency, line widths, point sizes, and positive canvas
dimensions.
\item Three-dimensional mean-identified-set figures may change the \emph{surface-mesh drawing
seed}, a valid integer that fixes the display-only random displacement of mesh points. The current seed
is 15,599. The mesh is triangular: successive rows sit \(h_M\sqrt3/2\) apart, alternate rows are
offset by \(h_M/2\), and points within a row sit \(h_M\) apart, where \(h_M\) is the positive
scalar horizontal lattice spacing. In each of the two normalized mesh coordinates, every point then
receives an independent uniform displacement between \(-\varkappa h_M\) and \(\varkappa h_M\),
where \(\varkappa\) is the dimensionless displacement fraction. The current configuration sets
\(h_M=0.16\) and, independently, \(\varkappa=0.16\); the two are unrelated quantities that happen
to share a value. The seed changes the triangular display mesh but not the mean identified set. The
figures may also change the \emph{wall-contour lattice size}, an integer of at least two that gives the
number of equally spaced values on each of the two coordinates of a background wall. Its current value is 440, so each wall uses
\(440^2\) function evaluations before its zero contour is drawn. Camera position, frame padding,
axis limits, tick locations and digits, and axis labels may also change. The figures
may use coefficient units or standard-deviation-scaled units and may include or omit the projected
mean least-squares benchmark.
\item Two-coordinate projection figures may change their lattice density to an integer of at least
two. They may also change colors, range padding, positive canvas dimensions, and panel layout.
\item The bounds-by-slack grid of Part~I admits three controls: its \emph{backbone count}, its
\emph{upper-tail refinement threshold}, and its \emph{upper-tail subdivision count}. The backbone
count must be a whole number of at least three, the threshold must lie strictly between zero and one,
and the subdivision count must be a whole number of at least two; any combination leaving either the
retained lower part or the refined upper part empty is refused before the grid is built. Each of the
three is an executable full-route modification, and this item is the exception the lead-in names:
each of the three changes which slacks the mean-equation stage solves and therefore both that stage's
stored boxes and the drawn figure, while leaving every reported estimate, identified range, and interval at the display slacks unchanged.
\item Conditional-scale figures may appear by
estimator and slack, combine several slacks, use a base-ten logarithmic axis, or normalize lower and
upper endpoint series separately.
\end{enumerate}
Each chosen variant must have a declared output location and must pass visual inspection.

\subsection{Preservation, validation, and cleanup}

Each serialization control below is an executable full-route modification: it enters the recorded
reuse identity, so changing it forces a fresh resampling run while leaving the route executable. The
reusable-result serialization version, 17-significant-digit numeric format, positive
round-trip tolerance, missing-value label, text-escape character, UTF-8 encoding, quarter-label
format, and the pattern that quarter labels are validated against may change. A coordinated change must update every reusable-result identity and round-trip
check.

An optional equivalence check can compare two independently assembled forms of the mean-equation
quadratic system at every evaluation, requiring exact equality of the resulting quadratic terms
and their component values. Enabling it is an executable full-route modification that adds an
exact-identity assertion at each assembly and changes no reported value. It adds validation work, but does not change the retained construction.

The cleanup procedure may remove reusable bootstrap results, route-status results, diagnostics, tables, figures,
or reports as separate compartments. The conditional cleanup and the typesetting-attempt controls run
inside the route and are executable full-route modifications; the six reset compartments are an
executable component modification, because the complete reproduction route never invokes the reset.
It may preserve publication files intended to remain while
clearing temporary working files. Resetting is not part of the current
reproduction. Conditional-output cleanup always precedes scientific calculations so a stale conditional
file cannot be mistaken for a newly produced result. The list of auxiliary typesetting extensions,
the positive number of cleanup attempts, and the nonnegative wait between attempts may change.

\section{Separate acceptance procedures}

The complete reproduction route does not run the three procedures in this section. They are
separate acceptance capabilities.

Output completeness is a manual acceptance step, not a procedure the reproduction performs. The
complete reproduction route proves only that the conditional-output locations agree with their
recorded route status: it stops when a route that ran left one of its outputs missing, and when a
route that did not run left any of its outputs present. Nothing compares the files present against
the whole inventory, so a verifier must make that comparison. For the current configuration the
comparison target is the 80 files of the current output inventory. The visual procedure inspects every table
and figure at its intended page size. Standalone compilation alone does not establish that an
output is current and readable.

The \emph{numerical-table acceptance comparison} compares one complete reference output with one
complete candidate output. It requires the same relative locations for every numerical table. It
then matches printed numbers by row and column, requires the same count of printed numbers, and
requires the same attached significance stars, counting the absence of stars as a value to be matched.
Its reach is narrower than the whole document: it reads only the cells that follow the first column
separator, only in rows that follow the horizontal rule, only inside a tabular block, and only in
typesetting sources below the tables folder. Numbers printed in a row's leading column and in header
rows are never compared. A pair of outputs from which nothing projects compares nothing and is
accepted.

For a printed scalar number \(x\in\mathbb R\), let \(q_x=10^{e_x-d_x}>0\). Here
\(e_x\) is the integer printed
base-ten exponent, set to zero when none is printed, and \(d_x\) is the nonnegative integer count of
digits printed after the decimal point in the leading number. Define the positive scalar \(q_y\) in
the same way for a candidate scalar number \(y\in\mathbb R\). Set
\begin{equation*}
 w_{xy}^R=\frac{q_x+q_y}{2},\qquad
 s_{xy}^R=\max(|x|,|y|,w_{xy}^R).
\end{equation*}
Here \(w_{xy}^R>0\) is the combined rounding width, and \(s_{xy}^R>0\) is the comparison scale. The
pair passes when \(x=y\) or
\begin{equation*}
 |x-y|<w_{xy}^R-8\epsilon_{\mathrm{mach}}s_{xy}^R.
\end{equation*}
The strict inequality supplies a binary64 numerical guard. The comparison can therefore accept
different printed character sequences that represent the same rounded value. It does not compare
labels, headings, captions, notes, figures, descriptive reports, input origins, or rendered
appearance.

\section{Declared unavailable joint-moment extensions}

The available joint-moment settings describe an additional-instrument moment block, a combined
log-and-exponential moment block, a mean-log intercept, an exponential-mean intercept, or both intercepts. Let
\(a_L\in\mathbb R\) denote the scalar mean-log intercept and let \(a_P\in\mathbb R\) denote the
scalar exponential-mean intercept. The settings can also describe a projected moment-compatibility
region, a strictly increasing vector of finite nonnegative moment tolerances, and a Gaussian
intercept-gap restriction.

None of these noncurrent settings is a full-route modification. The complete reproduction route
accepts only the current configuration: neither moment block, target \(a_L\), no projected
moment-compatibility region, an empty tolerance vector, and no Gaussian intercept-gap restriction. The
additional-instrument and combined log-and-exponential moment blocks, alternative intercept targets,
a projected moment-compatibility request, and a valid nonempty tolerance
vector are mathematically valid settings, but are refused before calculation. They are
accepted-but-refused choices. A Gaussian intercept-gap restriction is not accepted and is an
invalid choice. The associated calculations are declared unavailable extensions. Their mathematical
definitions follow without implying that the complete reproduction route can calculate them.

For the additional-instrument moment block, let \(Z^J\in\mathbb R^{T_v\times r_J}\) contain
\(r_J\ge1\) candidate additional instruments aligned to the conditional-scale sample, where
\(r_J\) is an integer. Define
\begin{equation*}
 M_V=I_{T_v}-X^V(X^{V\prime}X^V)^{-1}X^{V\prime},
 \qquad A^J=M_VZ^J.
\end{equation*}
Here \(I_{T_v}\in\mathbb R^{T_v\times T_v}\) is the identity matrix,
\(M_V\in\mathbb R^{T_v\times T_v}\) removes the linear projection on the conditional-scale
design, and \(A^J\in\mathbb R^{T_v\times r_J}\) is the residualized additional-instrument matrix.
Let \(q_J\ge0\) be the integer numerical rank of \(A^J\). When \(q_J\ge1\), let
\(B^J\in\mathbb R^{T_v\times q_J}\) be the deterministic orthonormal basis for its column space.
Orthonormal means \(B^{J\prime}B^J=I_{q_J}\). A zero rank makes the block unreliable. For structural SDF-news coefficient vector
\(\theta\in\mathbb R^3\), define the log-OLS residual vector
\begin{equation*}
 u^L(\theta)=\ell(\theta)-X^V\widehat\eta_L(\theta)\in\mathbb R^{T_v}.
\end{equation*}
The additional-instrument sample restriction is
\begin{equation*}
 g_J(\theta)=\frac1{T_v}B^{J\prime}u^L(\theta)=0_{q_J},
\end{equation*}
where \(g_J(\theta)\in\mathbb R^{q_J}\) is the additional-instrument sample-moment vector and
\(0_{q_J}\in\mathbb R^{q_J}\) is a vector of zeros. The population interpretation is that the
residualized additional instruments are uncorrelated with the log-OLS residual.

For the combined log-and-exponential moment block, write
\(X^V=[\mathbf1_{T_v},R]\), where \(R\in\mathbb R^{T_v\times4}\) contains the four centered lagged
supplied return-score columns of Part~I. Let \(\eta_R\in\mathbb R^4\) be one conditional-scale slope vector shared by both blocks. The
two five-dimensional sample-moment vectors are
\begin{align*}
 g_L(\theta,a_L,\eta_R)
 &=\frac1{T_v}X^{V\prime}
 \left[\ell(\theta)-X^V(a_L,\eta_R')'\right],\\
 g_P(\theta,a_P,\eta_R)
 &=\frac1{T_v}X^{V\prime}
 \left[q(\theta)-\exp\{X^V(a_P,\eta_R')'\}\right].
\end{align*}
Here \(g_L,g_P\in\mathbb R^5\), \(q(\theta)\in\mathbb R^{T_v}\) contains the squared
structural mean residuals, and the exponential is taken element by element. Stacking the two
vectors gives ten sample moments in the nine parameters consisting of \(\theta\), \(\eta_R\),
\(a_L\), and \(a_P\). The shared \(\eta_R\) is the cross-block restriction.

The two intercepts may instead be removed analytically. A \emph{profiled value} is substituted by
an exact formula rather than searched numerically. The profiled intercept values are
\begin{align*}
 a_L(\theta,\eta_R)
 &=\frac1{T_v}\sum_{t=1}^{T_v}
 \left[\log q_t(\theta)-R_t'\eta_R\right],\\
 a_P(\theta,\eta_R)
 &=\log\!\left\{\sum_{t=1}^{T_v}q_t(\theta)\right\}
 -\log\!\left\{\sum_{t=1}^{T_v}\exp(R_t'\eta_R)\right\}.
\end{align*}
Here \(a_L(\theta,\eta_R)\) and \(a_P(\theta,\eta_R)\) are scalars that set the intercept
coordinate of the corresponding sample-moment vector to zero. Removing those two coordinates
leaves eight slope moments in the seven parameters \(\theta\) and \(\eta_R\). An all-zero
\(q(\theta)\) cannot define \(a_P\).

To define a projected moment-compatibility region, stack the selected sample-moment vectors into
\(g(x)\in\mathbb R^m\), where \(x\in\mathbb R^k\) collects the parameters of the selected block,
\(m\) is the positive integer moment count, and \(k\) is the positive integer parameter count. Let
\(S=\operatorname{diag}(s_1,\ldots,s_m)\in\mathbb R^{m\times m}\), where each \(s_j>0\) is the
root-mean-square scale of observation-level moment contribution \(j\), held fixed during the search. A root-mean-square
scale is the square root of the mean squared contributions in that coordinate. Let
\(\epsilon_J>0\) be the structural-mean-residual crossing floor. For a finite nonnegative moment
tolerance \(d\), define
\begin{equation*}
 \mathcal C_d=
 \left\{x:\max_{1\le j\le m}\left|\frac{g_j(x)}{s_j}\right|\le d,
 \ \theta\in\Theta(\tau),\ \min_t|e_t(\theta)|\ge\epsilon_J\right\}.
\end{equation*}
Here \(\mathcal C_d\) is the moment-compatibility region at tolerance \(d\). A search over
\(\mathcal C_d\) can project \(a_L\), each coordinate of \(\eta_R\), and, for the combined block,
\(a_P\) and \(a_P-a_L\). The structural SDF-news coefficient vector is searched, but is not itself a reported projection target. A found point establishes compatibility within the search. Failure to
find one does not prove that \(\mathcal C_d\) is empty.

The Gaussian intercept-gap restriction sets
\begin{equation*}
 a_P-a_L=-\{\psi(1/2)+\log2\}=\gamma_E+\log2\approx1.270362845.
\end{equation*}
Here \(\psi(1/2)\) is the digamma function at one half, the first derivative of the logarithm of
the gamma function, and \(\gamma_E\) is the Euler--Mascheroni constant. This equality is a
distributional restriction implied by a squared standard-normal residual; it is not a numerical
tolerance.

The component calculations for these unavailable extensions expose the following numerical
controls. Changing them cannot make a refused choice executable.
\begin{enumerate}
\item The coefficient, constraint, and root tolerances govern coefficient convergence, feasibility,
and whether a sample-moment residual counts as zero.
\item The expanding-box widths govern successive search regions. The pattern starting-value cap
limits the number of starting values selected from distinct residual-sign patterns. The additional
starting-value cap limits further well-separated starting values. The remaining controls set the
minimum starting-value separation and the maximum evaluations in one search.
\item The replication lattice density and point cap are catalogued with the active stacked-moment
replication check above; they are not part of this unavailable extension.
The Jacobian rank and design rank tolerances determine when the corresponding matrices count as
having full rank. The basis rank tolerance does two other things: it screens out a candidate column
whose residual norm falls below that fraction of its own norm, and it sets the singular-value cutoff
above which a direction counts toward the basis rank, a rank that is reported rather than compared
against a full value. A separate basis-axis tolerance sets the residual norm at which a
projected canonical axis is accepted into the basis. Two basis gap factors decide whether the
retained-versus-discarded singular-value separation is clear enough to report the residualized basis
reliable, and a Jacobian gap factor decides when a Jacobian of full rank is nonetheless reported
weak.
\item The upper auxiliary-objective bound limits the nonnegative variable that represents the
largest scaled sample-moment residual. The absolute and relative starting offsets, which the record's
\nolinkurl{epigraph_start_abs_slack} and \nolinkurl{epigraph_start_rel_slack} fields carry, move this variable
strictly above its initial lower bound.
\item The interior-box tolerance determines whether an attained point is far enough from the edge of its search box to count as interior. The nested-floor tolerance determines when a new minimum is
small enough to replace the minimum carried from a smaller mean identified set. The floor-audit
tolerance compares that carried minimum across alternative numerical floor multipliers.
\item The projection-equality tolerance compares two projected values that should agree under an
enabled projected moment-compatibility calculation.
\end{enumerate}

\section{Output consequences of modifications}

The \emph{complete output inventory} contains 84 locations: 77 locations required regardless of
the LAD and residual-dynamics decisions, three LAD
locations, and four dynamic-volatility locations. The current output inventory contains 80 files
because LAD runs and the dynamic-volatility calculation is unavailable. A declined or unanswered LAD
decision reduces the required output inventory to the 77 unconditional files.

The four dynamic-volatility locations are reserved but unavailable. They are
\nolinkurl{diagnostics/log_var_eq_egarch_x.csv},
\nolinkurl{diagnostics/log_var_eq_egarch_x.rds},
\nolinkurl{state/log_var_eq_egarch_pilot.rds}, and
\nolinkurl{figures/log_var_eq_bounds_tau_egarch_x.pdf}. They must remain empty whenever the complete
reproduction route is executable. A route-status result may report a dynamic request, but a request does not authorize
the presence of these files.

Every full-route modification must update its affected identities, dimensions, notes, output labels,
reusable-result eligibility, and acceptance rules. A declared unavailable extension
remains unavailable until its complete reproduction route, outputs, and acceptance rules exist.

\end{document}
